%==============! USE XELATEX TO COMPILE !==================%
%==========================================================%
% To compile:                                              %
%     1. USE XELATEX                                       %
%     2. USE XELATEX                                       %
%     3. USE XELATEX                                       %
%     4. Use \newcommand{\renderOption}{0} for fast render %
%        (disabling most rendering of mahjong tiles)       %
%==========================================================%

\documentclass[12pt]{article}


\newcommand{\renderOption}{0}
% 0: fast render
% 1: full render


\usepackage[top=2cm, bottom=2cm, left=2cm, right=2cm]{geometry}
\usepackage{graphicx, expl3, minibox, amsmath, stfloats, tablefootnote, tikz, float, threeparttable, xfrac, multirow, etoolbox, verbatim, diagbox, titling, subfiles, subcaption, bibentry, amssymb}
\usetikzlibrary{backgrounds}

\usepackage{siunitx} % for aligning numbers within table
\sisetup{table-format=-4,
         table-space-text-pre={(},
         table-space-text-post={)},
         table-number-alignment = center,
         table-align-comparator = false,
         table-align-text-pre = false,
         table-align-text-post = false,
}

\usepackage{enumitem}
\setlist{noitemsep}

\usepackage{mahjong-tsumogiri/mahjong}
\graphicspath{{mahjong-tsumogiri/tiles/}}

\usepackage[CJKspace]{xeCJK}
\setCJKmainfont{YuMincho}

\usepackage[hidelinks]{hyperref}

\usepackage{babel}
\usepackage{natbib}
%\addto{\captionsjapanese}{\renewcommand{\refname}{参考文献}}

\usepackage[nottoc,numbib]{tocbibind} % to add references into the ToC

\usepackage{caption}
%\captionsetup[table]{name=表}
%\captionsetup[figure]{name=図}

\usepackage[explicit]{titlesec} % formatting \subsubsection* and underline
\usepackage{ulem}
\renewcommand{\ULthickness}{0.7pt} % default: 0.4pt
\renewcommand{\ULdepth}{4pt} % default: 3pt ish
% format only numberless subsubsections
\titleformat{name=\subsubsection,numberless}{\normalfont\normalsize\bfseries}{}{1em}{\hspace{-2ex}\uline{#1}}
%\titlespacing*{command}{left}{before-sep}{after-sep}[right-sep]
\titlespacing*{\subsubsection}{0pt}{3.25ex plus 1ex minus .2ex}{1.5ex plus .2ex}

\edef\restoreparindent{\parindent=\the\parindent\relax}
\usepackage[skip=\glueexpr \baselineskip + 0pt plus 2pt\relax]{parskip}

% store baselineskip
% https://tex.stackexchange.com/questions/392269/default-baselineskip-for-corresponding-font-size
\newlength\stdbls % standard baselineskip
\AtBeginDocument{%
  \setlength\stdbls{\baselineskip}%
  \gdef\stdblsn{\strip@pt\stdbls}%
}

%\usepackage{xr}
%\externaldocument[hairi-]{section/hairi}
%\externaldocument[joban-]{section/joban}
%\externaldocument[furo-]{section/furo}
%\externaldocument[riichi-]{section/riichi}
%\externaldocument[riichicare-]{section/riichicare}
%\externaldocument[furocare-]{section/furocare}
%\externaldocument[read-]{section/read}


%============================== New command and float ==============================

\newfloat{rittai}{t}{ritt}
\floatname{rittai}{Situation}

% mj command in footnote
\newcommand{\mjfn}[1]{\raisebox{-.135\stdbls}{\mahjong[0.75\stdbls]{#1}}}

\newcounter{tehaicount}
\renewcommand{\thetehaicount}{(\arabic{tehaicount})}

\newcommand{\onlyinsubfile}[1]{#1}
\newcommand{\notinsubfile}[1]{}
% subfiles only see commands before begin of document in base doc
% these are later to be redefined

\newcommand{\subfilepreamble}[1]{%
	\onlyinsubfile{\section{#1}}
	\onlyinsubfile{\tableofcontents\vspace{5mm}}
	\onlyinsubfile{\bibliographystyle{../jecon}}
  \onlyinsubfile{\nobibliography{../refmj}}
}
% the preamble that is common to all subfiles
% arg: name of standalone section



%============================== ExplSyntax ==============================
\ExplSyntaxOn

\DeclareExpandableDocumentCommand{\IfNoValueOrEmptyF}{mm}
{
  \IfNoValueF {#1}
  {
    \tl_if_empty:nF {#1} {#2}
  }
}

% adjust mahjong tile height
\NewDocumentCommand{\mj}{m}{
  \ifdefstring{\renderOption}{1}{
    \raisebox{-.18\stdbls}{\mahjong{#1}}
  }{
    \tl_set:Nn \l_tmpa_tl {#1}
    \tl_replace_all:Nnn \l_tmpa_tl { ^ } { / } % text mode does not allow ^
    \tl_use:N \l_tmpa_tl
  }
}

\NewDocumentCommand{\tehai}{smO{}O{}}{
  % DECIDE FAST RENDER OR FULL RENDER
  \ifdefstring{\renderOption}{1}{
    \nankiru_drawTehai:nnnn {#1} {#2} {#3} {#4}
  }{
    \nankiru_drawTehaiFast:nnnn {#1} {#2} {#3} {#4}
  }
}

\cs_new_protected:Npn \nankiru_drawTehai:nnnn #1 #2 #3 #4 {
  \tl_if_empty:nTF {#3} {
    \noindent \IfBooleanF {#1} {\refstepcounter{tehaicount} \thetehaicount \hspace{1em} \phantomsection \IfNoValueOrEmptyF{#4}{\label{mj:#4}}} \raisebox{-.35\height}{\mahjong[1.5\stdbls]{#2}} 
  }{
    \noindent \IfBooleanF {#1} {\refstepcounter{tehaicount} \thetehaicount \hspace{1em} \phantomsection \IfNoValueOrEmptyF{#4}{\label{mj:#4}}} \minibox{#3 \\ \mahjong[1.5\stdbls]{#2}} 
  }
}

\cs_new_protected:Npn \nankiru_drawTehaiFast:nnnn #1 #2 #3 #4 {
  % show plain text
  \tl_set:Nn \l_tmpa_tl {#2}
  \tl_replace_all:Nnn \l_tmpa_tl { ^ } { / } % text mode does not allow ^
  \tl_if_empty:nTF {#3} {
    \noindent \IfBooleanF {#1} {\refstepcounter{tehaicount} \thetehaicount \hspace{1em} \phantomsection \IfNoValueOrEmptyF{#4}{\label{mj:#4}}} \tl_use:N \l_tmpa_tl
  }{
    \noindent \IfBooleanF {#1} {\refstepcounter{tehaicount} \thetehaicount \hspace{1em} \phantomsection \IfNoValueOrEmptyF{#4}{\label{mj:#4}}} \minibox{#3 \\ \tl_use:N \l_tmpa_tl}
  }
}

\cs_generate_variant:Nn \tl_remove_once:Nn { NV }
\cs_generate_variant:Nn \tl_remove_once:Nn { Nx }

% define keys for rittai
\keys_define:nn { rittai }
{
  jicha/  seat  .tl_set:N = \l_rittai_jicha,
  jicha/  seat  .initial:n = E,
  jicha/  hand  .tl_set:N = \l_rittai_jichahand,
  jicha/  hand  .value_required:n = true,
  jicha/  river .tl_set:N = \l_rittai_jichariver,
  jicha/  riverb.tl_set:N = \l_rittai_jichariverb,
  jicha/  riverc.tl_set:N = \l_rittai_jichariverc,
  jicha/  point .tl_set:N = \l_rittai_jichapoint,
  jicha/  point .initial:n = 25000,
  
  toimen/ hand  .tl_set:N = \l_rittai_toimenhand,
  toimen/ hand  .initial:n = XXXXXXXXXXXXX,
  toimen/ river .tl_set:N = \l_rittai_toimenriver,
  toimen/ riverb.tl_set:N = \l_rittai_toimenriverb,
  toimen/ riverc.tl_set:N = \l_rittai_toimenriverc,
  toimen/ point .tl_set:N = \l_rittai_toimenpoint,
  toimen/ point .initial:n = 25000,
  
  kami/   hand  .tl_set:N = \l_rittai_kamihand,
  kami/   hand  .initial:n = XXXXXXXXXXXXX,
  kami/   river .tl_set:N = \l_rittai_kamiriver,
  kami/   riverb.tl_set:N = \l_rittai_kamiriverb,
  kami/   riverc.tl_set:N = \l_rittai_kamiriverc,
  kami/   point .tl_set:N = \l_rittai_kamipoint,
  kami/   point .initial:n = 25000,
  
  shimo/  hand  .tl_set:N = \l_rittai_shimohand,
  shimo/  hand  .initial:n = XXXXXXXXXXXXX,
  shimo/  river .tl_set:N = \l_rittai_shimoriver,
  shimo/  riverb.tl_set:N = \l_rittai_shimoriverb,
  shimo/  riverc.tl_set:N = \l_rittai_shimoriverc,
  shimo/  point .tl_set:N = \l_rittai_shimopoint,
  shimo/  point .initial:n = 25000,
  
    
  kyoku   .tl_set:N = \l_rittai_kyoku,
  kyoku   .initial:n = East-1,
  honba   .tl_set:N = \l_rittai_honba,
  honba   .initial:n = 0,
  kyoutaku.tl_set:N = \l_rittai_kyoutaku,
  kyoutaku.initial:n = 0,
  dora    .tl_set:N = \l_rittai_dora,
  dora    .initial:n = 0m,
}

\NewDocumentCommand{\drawrittai}{m}{
  % DECIDE FAST RENDER OR FULL RENDER
  \ifdefstring{\renderOption}{1}{
    \nankiru_drawRittai:n {#1}
  }{
    \nankiru_drawRittaiFast:n {#1}
  }
}

% draw rittai (ver 2 10/15/2023)
\tl_new:N \l_rittai_sws
\tl_new:N \l_rittai_sw
  
\cs_new_protected:Npn \nankiru_drawRittai:n #1 {
  \keys_set:nn { rittai } {#1}
  
  \begin{tikzpicture}[framed]
    
    %----------------------hand
    \node[anchor=south] (self) at (0,-3.7) {\mahjong[1.2\stdbls]{\tl_use:N \l_rittai_jichahand}};
    % ver 2: others hands are smaller
    \node[anchor=south,rotate=180] (toimen) at (0,3.7) {\mahjong[1.2\stdbls]{\tl_use:N \l_rittai_toimenhand}};
    \node[anchor=south,rotate=270] (kami) at (-5.2,0) {\mahjong[1.2\stdbls]{\tl_use:N \l_rittai_kamihand}};
    \node[anchor=south,rotate=90] (shimo) at (5.2,0) {\mahjong[1.2\stdbls]{\tl_use:N \l_rittai_shimohand}};
    
    %----------------------kyoku
    % ver 2: no rectangular box
    \node[align=center] (center-kyoku) at (0,0) {\tl_use:N \l_rittai_kyoku\\Dora~\raisebox{-.25\height}{\mahjong[1.2\stdbls]{\tl_use:N \l_rittai_dora}}};
    
    %----------------------seat wind and point
    % ver 1: seat wind and point are at center of table
    % ver 2: seat wind and point now appear above hand
    %        honba and kyoutaku appear above own point
    \tl_set:Nn \l_rittai_sws {ESWNESWN}
    
    \int_do_while:nn { \l_tmpa_int < 5 } {
      \tl_set:Nx \l_rittai_sw {\tl_head:N \l_rittai_sws}
      
      \tl_if_eq:NNT \l_rittai_sw \l_rittai_jicha {
        \node[anchor=south~west,align=left] at (-3.6,-2.9) {\,\mahjong[0.6\stdbls]{h}~\tl_use:N \l_rittai_honba \mahjong[0.6\stdbls]{-k}~\tl_use:N \l_rittai_kyoutaku \\ \tl_use:N \l_rittai_jicha \ \tl_use:N \l_rittai_jichapoint};
        
        \tl_remove_once:Nx \l_rittai_sws {\tl_head:N \l_rittai_sws}
        \node[anchor=south~west,align=left,rotate=180] at (3.6,2.9) {\tl_head:N \l_rittai_sws \ \tl_use:N \l_rittai_toimenpoint};
        
        \tl_remove_once:Nx \l_rittai_sws {\tl_head:N \l_rittai_sws}
        \node[anchor=south~west,align=left,rotate=90] at (4.3,-2.5) {\tl_head:N \l_rittai_sws \ \tl_use:N \l_rittai_shimopoint};
        
        \tl_remove_once:Nx \l_rittai_sws {\tl_head:N \l_rittai_sws}
        \node[anchor=south~west,align=left,rotate=270] at (-4.3,2.5) {\tl_head:N \l_rittai_sws \ \tl_use:N \l_rittai_kamipoint};
      }
      
      \int_incr:N \l_tmpa_int
      \tl_remove_once:NV \l_rittai_sws \l_rittai_sw
    }
    
    %----------------------river
    \tl_if_empty:NF \l_rittai_jichariver {
      \node[anchor=north~west, align=left] (riv-self) at (-1.5,-0.6) {
        \mahjong[1.2\stdbls]{\tl_use:N \l_rittai_jichariver} \tl_if_empty:NF \l_rittai_jichariverb {
          \\ \mahjong[1.2\stdbls]{\tl_use:N \l_rittai_jichariverb} \tl_if_empty:NF \l_rittai_jichariverc {
            \\ \mahjong[1.2\stdbls]{\tl_use:N \l_rittai_jichariverc}
          }
        }
      };
    }
    
    \tl_if_empty:NF \l_rittai_toimenriver {
      \node[anchor=north~west, rotate=180, align=left] (riv-toimen) at (1.5,0.6) {
        \mahjong[1.2\stdbls]{\tl_use:N \l_rittai_toimenriver} \tl_if_empty:NF \l_rittai_toimenriverb {
          \\ \mahjong[1.2\stdbls]{\tl_use:N \l_rittai_toimenriverb} \tl_if_empty:NF \l_rittai_toimenriverc {
            \\ \mahjong[1.2\stdbls]{\tl_use:N \l_rittai_toimenriverc}
          }
        }
      };
    }
    
    \tl_if_empty:NF \l_rittai_kamiriver {
      \node[anchor=north~west, rotate=270, align=left] (riv-kami) at (-1.5,1.5) {
        \mahjong[1.2\stdbls]{\tl_use:N \l_rittai_kamiriver} \tl_if_empty:NF \l_rittai_kamiriverb {
          \\ \mahjong[1.2\stdbls]{\tl_use:N \l_rittai_kamiriverb} \tl_if_empty:NF \l_rittai_kamiriverc {
            \\ \mahjong[1.2\stdbls]{\tl_use:N \l_rittai_kamiriverc}
          }
        }
      };
    }
    
    \tl_if_empty:NF \l_rittai_shimoriver {
      \node[anchor=north~west, rotate=90, align=left] (riv-shimo) at (1.5,-1.5) {
        \mahjong[1.2\stdbls]{\tl_use:N \l_rittai_shimoriver} \tl_if_empty:NF \l_rittai_shimoriverb {
          \\ \mahjong[1.2\stdbls]{\tl_use:N \l_rittai_shimoriverb} \tl_if_empty:NF \l_rittai_shimoriverc {
            \\ \mahjong[1.2\stdbls]{\tl_use:N \l_rittai_shimoriverc}
          }
        }
      };
    }
  \end{tikzpicture}
}


\cs_new_protected:Npn \nankiru_drawRittaiFast:n #1 {
  \begin{tikzpicture}
    \filldraw [fill=gray!10, draw=black] (-4, -2.8) rectangle (4, 2.8);
    \node[draw, align=center] (center-kyoku) at (0,0) {\tl_use:N \l_rittai_kyoku\\Dora~\raisebox{-.25\height}{\mahjong[1.2\stdbls]{\tl_use:N \l_rittai_dora}}};
  \end{tikzpicture}
}


% multicitep: cite multiple references in one parenthesis
% ex: \multicitep{adams03, p.~3; collier09, p.~5}
\makeatletter
\NewDocumentCommand{\multicitep}{m}
 {
  \NAT@open
  \mjb_multicitep:n { #1 }
  \NAT@close
 }
\makeatother

\seq_new:N \l_mjb_multicite_in_seq
\seq_new:N \l_mjb_multicite_out_seq
\seq_new:N \l_mjb_cite_seq

\cs_new_protected:Npn \mjb_multicitep:n #1
 {
  \seq_set_split:Nnn \l_mjb_multicite_in_seq { ; } { #1 }
  \seq_clear:N \l_mjb_multicite_out_seq
  \seq_map_inline:Nn \l_mjb_multicite_in_seq
   {
    \mjb_cite_process:n { ##1 }
   }
  \seq_use:Nn \l_mjb_multicite_out_seq { ;~ }
 }

\cs_new_protected:Npn \mjb_cite_process:n #1
 {
  \seq_set_split:Nnn \l_mjb_cite_seq { , } { #1 }
  \int_compare:nTF { \seq_count:N \l_mjb_cite_seq == 1 }
   {
    \seq_put_right:Nn \l_mjb_multicite_out_seq
     { \citeauthor{#1},~\citeyear{#1} }
   }
   {
    \seq_put_right:Nx \l_mjb_multicite_out_seq
     {
      \exp_not:N \citeauthor{\seq_item:Nn \l_mjb_cite_seq { 1 }},~
      \exp_not:N \citeyear{\seq_item:Nn \l_mjb_cite_seq { 1 }},~
      \seq_item:Nn \l_mjb_cite_seq { 2 }
     }
   }
 }


\ExplSyntaxOff




%======================= END OF PREAMBLE ==========================
\begin{document}

\renewcommand{\onlyinsubfile}[1]{}
\renewcommand{\notinsubfile}[1]{#1}

\title{Spaghet Boi's Handbook of Handbuilding}
\renewcommand\maketitlehookc{\vspace{-2ex}}
\maketitle

\tableofcontents

%\section{牌理}
%\subfile{section/hairi}
%
%\section{副露判断}
%\subfile{section/furo}
%
%\section{立直判断}
%\subfile{section/riichi}
%
%\section{立直対応}
%\subfile{section/riichicare}
%
%\section{副露対応}
%\subfile{section/furocare}
%
%\section{河読み}
%\subfile{section/read}

\section{Tile Efficiency} \label{sec:tile-eff}

Tile efficiency is a set of policies and techniques that allows us to reach an optimal point when taking speed, value, win rate, and many other factors into consideration. There is also the concept of \textit{pure} tile efficiency which concerns of only the speed of reaching \textit{tenpai}, one tile away from completing our hand. 

In most rulesets such as M. League ruleset where there are yakus such as dora, akadora, ippatsu, and uradora to inflate the expected value of riichi, the best course of action is often to become the first player to call riichi, and the speed and opportunity of reaching tenpai depend on how wide the acceptance of our hand is. Therefore, pure efficiency is very often prioritized and preached. Indeed, mahjong is not a math contest where the player who can calculate the most acceptance automatically wins the game; there are many other factors that would require us to sacrifice efficiency for a better overall outcome, and we will see how those factors play in a bigger picture of hand building.

The following text assumes that the reader has a basic understanding of mahjong rules such as suits, tiles, win conditions, call mechanisms, and riichi. There are, however, many terms that I cannot find a proper translation in English, so I will use Japanese terms as is. Some of them are listed in Table~\ref{tab:term}. Many other terms such as \textit{ryanmenkanchan} and \textit{agaritoppu} will be introduced later in the text.

\begin{table}[ht]
  \centering
  \begin{tabular}{ll}
    Term & Definition \\
    \hline
    Ryanmen & The shape consisting of two neighboring number tiles such as \mj{23p} (23p) \\
    Kanchan & The shape consisting of tiles whose numbers have difference of 2 such as \mj{24p} (24p) \\
    Penchan & The shape consisting of number tiles of 1 or 9 and a neighboring tile such as \mj{12p} (12p) \\
    Toitsu & The shape consisting of two identical copies such as \mj{11p} (11p) \\
    Taatsu & The set of ryanmen, kanchan, and penchan \\
    Shuntsu & The shape consisting of three consecutive number tiles such as \mj{123p} (123p) \\
    Koutsu & The shape consisting of three identical copies such as \mj{111p} (111p) \\
    Minko & A koutsu resulting from pon call and thus open to all players \\
    Anko & A koutsu in our hand invisible to other players \\
    Mentsu & The set of koutsu and shuntsu, completed shapes \\
    Tenpai & One tile away from completing our hand \\
    1-away & One tile away from tenpai, also called \textit{iishanten} \\
    2-away & Two tiles away from tenpai, also called \textit{ryanshanten} \\
    Shanpon & A wait pattern where we have two pairs in our hand
  \end{tabular}
  \caption{Terminology}
  \label{tab:term}
\end{table}


\subsection{Complex Taatsu} \label{subsec:complex-taatsu}

Complex taatsus are shapes involving simple taatsus (penchans, kanchans, and ryanmens) and an additional tile. Table~\ref{tab:fkgtt-1} lists complex taatsus that uses three tiles.

\begin{table}[ht]
  \centering
  \begin{tabular}{lll}
    Name & Shape & Acceptance \\
    \hline
    Penchan-toitsu & \mj{112m} & \mj{13m} 6 tiles \\
    Kanchan-toitsu & \mj{113m} & \mj{12m} 6 tiles \\
    Ryankanchan & \mj{135m} & \mj{24m} 8 tiles \\
    Ryanmen-toitsu & \mj{223m} & \mj{124m} 10 tiles
  \end{tabular}
  \caption{Types of complex taatsus using three tiles}
  \label{tab:fkgtt-1}
\end{table}

Indeed, there are other kinds of complex shapes such as \mj{12234m}, \mj{2345m}, \mj{334568m}, and so on. Shapes with mentsus involved have more tiles and thus less frequent than shapes with less tiles, but their complexity can still easily confuse us. In Section \ref{subsec:complex-taatsu}, we will deal with complex shapes that has no mentsu inside. Shapes with mentsus will be covered later.


\subsubsection{Complex Taatsus with Pairs} \label{subsubsec:toitsu}

Since complex taatsus from Table \ref{tab:fkgtt-1} (except for ryankanchan) contain a pair, they must serve as head when there is no other head candidate in our hand. On the other hand, when we have more than two complex taatsus containing pairs, the pairs will impair our acceptance, so in most cases we opt to break those pairs and reduce them to single taatsus. In order to increase the probability of reaching good waits, as a rule of thumb, we drop one tile from the strongest complex taatsu \citep{gendai-complex1, gendai-complex2}. Table~\ref{tab:fkgtt-2} lists our actions based on the number of pairs, complex taatsus (excluding ryankanchans), and ryanmen-toitsus in our hand。

\begin{table*}[ht]
  \centering
  \begin{tabular}{llll}
    Pair & Complex taatsu & Ryanmen-toitsu & Action \\
    \hline
    2 & 2 &   & Reduce the weakest complex taatsu to a pair \\
    3 & 2 & 0 & Break the pair in the inner kanchan-toitsu \\
      &   &   & Reduce outer penchan-toitsu to a pair \\
    3 & 2 & 1 & Reduce the ryanmen-toitsu to a ryanmen \\
    3 & 3 & 0 & Reduce the weakest complex taatsu to a pair \\
    3 & 3 & 1 & Reduce the ryanmen-toitsu to a ryanmen \\
    3 & 3 & 2 & Reduce the weakest complex taatsu to pair
  \end{tabular}
  \caption{Action based on the number of pairs, complex taatsus, and ryanmen-toitsus.}
  \label{tab:fkgtt-2}
\end{table*}

Note that Table~\ref{tab:fkgtt-2} is just a rule of thumb. Depending on many other factors such as the existence of yakuhai pairs or table situations, we may take alternatives.


\subsubsection*{Two Pairs}

We must drop a tile from either \mj{233p} or \mj{466s} from \ref{mj:fkgtt-1} since we will not be touching completed groups and simple ryanmens. To keep two pairs within our hand, our choices narrow down to \mj{2p} and \mj{4s}. Their acceptances are listed in Table~\ref{tab:fkgtt-3}, which suggests that dropping \mj{4s} has 4 more tiles than dropping \mj{2p} in terms of acceptance. This is a case of reducing the weaker complex taatsu to a pair and keeping the stronger complex taatsu, and this particular shape of 1-away is called \textit{perfect 1-away} or \textit{kanzen iishanten} in Japanese.

\tehai{12334578m 233p 466s}[][fkgtt-1]

\begin{table}[ht]
  \centering
  \begin{tabular}{lll}
    Drop & \mj{4s} & \mj{2p} \\
    \hline
    Acceptance & \mj{69m 134p 6s} 20 tiles & \mj{69m 3p 56s} 16 tiles \\
  \end{tabular}
  \caption{Acceptance of \ref{mj:fkgtt-1}.}
  \label{tab:fkgtt-3}
\end{table}

\ref{mj:fkgtt-1a} has \mj{112345p} which makes dropping \mj{1p} and reducing it to \mj{45p} ryanmen attractive. But by doing so, we lost flexibility in manzu shape and we must make either \mj{66m} or \mj{77m} our head, thus reducing our acceptance. If we see \ref{mj:fkgtt-1a} as choosing between \mj{667m}\footnote{Strictly speaking, \mjfn{666778m} can be treated as either \mjfn{666m} + \mjfn{778m} or \mjfn{667m} + \mjfn{678m}, depending on which mentsu we extract from it. Either way, manzu shape is made of a ryanmen-toitsu and a mentsu; it does not change the fact that we are comparing between a ryanmen-toitsu and a penchan-toitsu.} and \mj{112p}, it becomes clear that dropping \mj{2p} here gives us the widest acceptance. Now, dropping \mj{5p} also gives the same acceptance as \mj{2p}, but considering the case when we draw \mj{47s} and wait on \mj{69m1p}, dropping \mj{2p} here makes the outside tile \mj{1p} more likely to come out. Hence drop \mj{2p} \citep{shibukawa-twopairs}. Table~\ref{tab:fkgtt-1a} lists our acceptance.

\tehai{666778m 112345p 06s}[East-1, East seat, turn 6, dora \mj{4m}][fkgtt-1a]

\begin{table}[ht]
  \centering
  \begin{tabular}{lll}
    Drop & \mj{1p} & \mj{2p} \\
    \hline
    Acceptance & \mj{36p47s} 15 tiles & \mj{56789m1p47s} 24 tiles \\
  \end{tabular}
  \caption{Acceptance of \ref{mj:fkgtt-1a}.}
  \label{tab:fkgtt-1a}
\end{table}

If our choices are between two ryanmen-toitsus, reduce the one on the inside into a pair to make sure that when we reach tenpai our waits are outside, which has higher win rate. Therefore, between \mj{4m} and \mj{5s} from~\ref{mj:fkgtt-1-1}, we drop \mj{5s} and keep the possibility of \mj{25m} wait \citep[][pp.~208--209]{suphx}.

\tehai{334789m 34p 111566s}[][fkgtt-1-1]

\ref{mj:two-ryanmentoitsu} also has two ryanmen-toitsus \mj{778m} and \mj{556s}, and one is connected to a completed mentsu \mj{345s}. Upon seeing this complex shape, one may be reluctant to touch complex souzu shape and tempted to drop \mj{7m} to fix \mj{78m} ryanmen or drop \mj{8m}. The \mj{7m} drop forces us to treat \mj{55s} as our head and renders \mj{36s} surplus, thus having less acceptance than \mj{8m} drop, while \mj{8m} drop leads us to final wait on either \mj{47p} or \mj{47s}. Our previous conclusion applies here as well: drop \mj{6s} since the ryanmen-toitsu \mj{556s} is inside \citep{shibukawa-twopairs}. Table~\ref{tab:two-ryanmentoitsu} lists our acceptance.

\tehai{778m 12356p 345556s}[East-1, East seat, turn 6, dora \mj{6p}][two-ryanmentoitsu]

\begin{table}[ht]
  \centering
  \begin{tabular}{llll}
    Drop & \mj{7m} & \mj{8m} & \mj{6s} \\
    \hline
    Acceptance & \mj{69m47p} 16 tiles & \mj{7m47p2457s} 22 tiles & \mj{679m47p25s} 23 tiles
  \end{tabular}
  \caption{Acceptance of \ref{mj:two-ryanmentoitsu}.}
  \label{tab:two-ryanmentoitsu}
\end{table}

It should be noted that in all cases with only two pairs, reducing complex taatsus to simple taatsus will not lead to maximal acceptance. When we have two pairs in our hand, we can accept tiles of the same kind as these pairs; any tiles that follows these pairs also increases acceptance. If we only have one pair, say \mj{23p466s}, then the tile that follows the pair becomes useless, since drawing \mj{5s} destroys the pair and produces a mentsu, which renders our hand headless and unready. This means that sometimes we have to counterintuitively destroy ryanmens.


\subsubsection*{Three Pairs, Two Complex Taatsus}

Since we have more than two pairs in \ref{mj:fkgtt-2}, we opt to break one, either from \mj{233p} or \mj{466s}\footnote{Note that we do not drop \mjfn{8m} here since that is equal to dropping entire block and increasing our shanten number.}. Breaking the pair in either complex taatsus will give us the same number of acceptance, but the numbers of draws that lead to ryanmen tenpai are different and listed in Table~\ref{tab:fkgtt-4a}. This is a case of reducing stronger complex taatsu to a simple taatsu \citep{gendai-complex1}.

\tehai{12334588m 233p 466s}[][fkgtt-2]

\begin{table}[ht]
  \centering
  \begin{tabular}{lll}
    Drop & \mj{3p} & \mj{6s} \\
    \hline
    Acceptance & \mj{8m14p56s} 16 tiles & \mj{8m134p5s} 16 tiles \\
    Ryanmen-tenpai draws & \mj{8m 56s} 8 tiles & \mj{5s} 4 tiles \\
  \end{tabular}
  \caption{Acceptance of \ref{mj:fkgtt-2}.}
  \label{tab:fkgtt-4a}
\end{table}

If our choices are between complex taatsus without ryanmens, there is no difference in direct acceptance and ryanmen-tenpai draws. If our complex taatsus contains kanchans, there is a chance for those to upgrade into ryanmens. It is therefore optimal to break the pair from inner kanchan-toitsus, since they have 8 tiles of ryanmen improvement as opposed to 4 tiles for outside kanchans. This means dropping \mj{6s} from \mj{466s} in \ref{mj:fkgtt-3}. Drawing \mj{37s} improves our \mj{46s} kanchan into ryanmen, while only \mj{4p} can produce a ryanmen taatsu out of \mj{13p} \citep{gendai-complex1}.

\tehai{12334588m 133p 466s}[][fkgtt-3]

For complex taatsus with few ryanmen-improvement draws, if the remaining pair is in the middle of a suit, we can expect its ryanmen improvement. We need to reduce a complex taatsu into a pair for that outcome, and it then becomes optimal to fix the weakest complex taatsu into a pair. \ref{mj:fkgtt-4} has \mj{899s} that, when turned into a penchan (simple taatsu), has no ryanmen improvement, compared to \mj{113p} which produces a ryanmen when we draw \mj{4p}, so we drop \mj{8s} and make it a pair \citep{gendai-complex1}.

\tehai{123345m 11366p 899s}[][fkgtt-4]

\ref{mj:fkgtt-4-1} has a more complex shape \mj{11333445p} that can be treated as either \mj{11p} + \mj{345p} + \mj{334p} or \mj{11p} + \mj{333p} + \mj{445p}, which leads to its wide acceptance of \mj{123456p}. We will lose some of its acceptance if we drop any tile from this shape, so instead we reduce \mj{778s} to a ryanmen, since technically we still have three pairs in our hand. Table~\ref{tab:fkgtt-4-1} lists the acceptance after different drops. This leads to the rule of thumb that we keep the more complex shape with more acceptance intact \citep[][p.~24]{uzaku-nankiru}.

\tehai{234m11333445p778s}[East-1, North seat, turn 7, dora \mj{3p}][fkgtt-4-1]

\begin{table}[ht]
  \centering
  \begin{tabular}{llll}
    Drop & \mj{7s} & \mj{4p} & \mj{3p} \\
    \hline
    Acceptance & \mj{123456p69s} 24 tiles & \mj{1236p679s} 21 tiles & \mj{125p679s} 19 tiles
  \end{tabular}
  \caption{Acceptance of \ref{mj:fkgtt-4-1}.}
  \label{tab:fkgtt-4-1}
\end{table}

\ref{mj:fkgtt-4-2} is also choosing between two ryanmen-taatsus \mj{556m} and \mj{566s}, and reducing either to a ryanmen does not make any difference in terms of acceptance. However, if we keep \mj{345566s} and draw another dora \mj{3s}, we can drop \mj{22z} and go for tanyao. In that regard, it is best that we drop \mj{5m} here \citep[][p.~24]{uzaku-nankiru}.

\tehai{556m 678p 345566s 22z}[East-1, East seat, turn 6, dora \mj{3s}][fkgtt-4-2]

There is no difference in terms of immediate acceptance between \mj{8m} drop and \mj{2s} drop from \ref{mj:fkgtt-ssk}. However, in case of drawing \mj{7p} and thus the potential 789 sanshoku, we must fix \mj{78m} ryanmen. With \mj{8m} drop and \mj{7p} draw, we can proceed to drop \mj{3s} from the shape \mj{78m788p223s} \citep{nemata-16}。

\tehai{788m 23488p 223789s}[Dora \mj{4z}][fkgtt-ssk]

\textbf{When there is a yakuhai pair, keep 3 pairs} \citep[][p.~72]{uzaku-haikouritsu}. Yakuhai pairs contribute almost\footnote{Well, it does give 2 fu, but that is negligible compared to the value an addition han brings.} no value unless it becomes a triplet or at least is part of our final waits. Therefore, we do not expect it to be our head, and we must find our head somewhere else. Dropping either of \mj{8p2s} from~\ref{mj:fkgtt-4-3} gives us 2 less tiles of acceptance than maximal number, but it keeps the option of \mj{7z} pon. Since pinzu shape has draws like \mj{2p} that potentially improve our shapes and souzu shape only has \mj{5s} improvement, we drop \mj{8p} here and keep the follow-up tile for \mj{4s} pair.

\tehai{789m 345668p 244s 77z}[][fkgtt-4-3]

Similarly for \ref{mj:fkgtt-4-4}, under the premise that \mj{2z} will be ponned, we must drop either of \mj{8m} and \mj{7s}. If we drop \mj{7s} and draw \mj{7m}, \mj{05s} can only serve as head and cannot accept dora \mj{4s}. Therefore, drop \mj{8m}.

\tehai{135899m 123p 057s 22z}[South-3, East seat, turn 5, dora \mj{4s}][fkgtt-4-4]


\subsubsection*{Three Pairs, Three Complex Taatsus, No Ryanmen-Toitsus}

If there are three complex taatsus, all of which contain no ryanmen, then it is optimal to reduce the one with least ryanmen improvement into a pair, so that when the remaining complex taatsus becomes a mentsu, there will be no surplus tiles. \ref{mj:fkgtt-5-1} is the hand derived from \ref{mj:fkgtt-5} by dropping \mj{8s} and completing \mj{668p}. \ref{mj:fkgtt-5-2} is derived from \ref{mj:fkgtt-5} by dropping \mj{9s} instead and completing \mj{668p}. We can see, in Table~\ref{tab:fkgtt-4}, that dropping \mj{8s} loses \mj{7s} acceptance and we have 2 tiles of acceptance less than dropping \mj{9s}, but once \mj{668p} becomes a mentsu, \ref{mj:fkgtt-5-1} has 4 tiles of acceptance more than \ref{mj:fkgtt-5-2}. We therefore prioritize the width of 1-away acceptance and sacrifice a little bit at 2-away. Drop \mj{8s} \citep{gendai-complex1}.

\tehai{12378m 113668p 899s}[][fkgtt-5]

\tehai{12378m 113678p 99s}[][fkgtt-5-1]

\tehai{12378m 113678p 89s}[][fkgtt-5-2]

\begin{table}[ht]
  \centering
  \begin{tabular}{lll}
    Drop & \mj{8s} & \mj{9s} \\
    \hline
    Acceptance & \mj{69m1267p9s} 22 tiles & \mj{69m1267p7s} 24 tiles \\
    Hand after completing \mj{668p} & \ref{mj:fkgtt-5-1} & \ref{mj:fkgtt-5-2} \\
    Acceptance after completing \mj{668p} & \mj{69m12p9s} 16 tiles & \mj{69m7s} 12 tiles
  \end{tabular}
  \caption{Acceptance of \ref{mj:fkgtt-5}, \ref{mj:fkgtt-5-1}, and \ref{mj:fkgtt-5-2}}
  \label{tab:fkgtt-4}
\end{table}


\subsubsection*{Three Pairs, Three Complex Taatsus, One Ryanmen-Toitsu}

If there is one ryanmen-toitsu among all three complex taatsus, it is optimal to reduce the ryanmen-toitsu to a ryanmen. There will be no difference in immediate acceptance between the drop candidates \mj{26p9s} from \ref{mj:fkgtt-6}, but there will be huge differences when we draw any of the acceptances, as shown in Table~\ref{tab:fkgtt-6}. If we drop \mj{6p}, completing the remaining two complex taatsus \mj{223p} and \mj{899s} will leave us only 3 kinds 12 tiles of acceptance. Likewise for \mj{9s} drop. On the other hand, if we drop \mj{2p}, all effective draws will guarantee us 16 tiles of acceptance. One may be tempted by the \mj{7p} draw after dropping \mj{6p} which gives us perfect 1-away, but also keep in mind that it is more likely to draw ourselves into 3 kinds 12 tiles of acceptance.

\tehai{12378m 223668p 899s}[][fkgtt-6]

\begin{table}[ht]
  \centering
  \begin{tabular}{lll}
    Drop & Next draw & Max acceptance \\
    \hline
    \mj{2p} & \mj{69m14p} & drop \mj{8p} $\to$ \mj{146p79s} 16 tiles \\
            &         & drop \mj{8s} $\to$ \mj{1467p9s} 16 tiles \\
            & \mj{6p} & drop \mj{8p8s} $\to$ \mj{69m14p} 16 tiles \\
            & \mj{7p} & drop \mj{6p8s} $\to$ \mj{69m14p} 16 tiles \\
            & \mj{7s} & drop \mj{8p9s} $\to$ \mj{69m14p} 16 tiles \\
            & \mj{9s} & drop \mj{8p8s} $\to$ \mj{69m14p} 16 tiles \\
    \hline
    \mj{6p} & \mj{69m} & drop \mj{8s} $\to$ \mj{1247p9s} 16 tiles \\
            & \mj{14p} & drop \mj{2p8s} $\to$ \mj{69m7p} 12 tiles \\
            & \mj{2p} & drop \mj{3p8s} $\to$ \mj{69m7p} 12 tiles \\
            & \mj{7p} & drop \mj{8s} $\to$ \mj{69m124p9s} 20 tiles \\
            & \mj{7s} & drop \mj{3p9s} $\to$ \mj{69m7p} 12 tiles \\
            & \mj{9s} & drop \mj{3p8s} $\to$ \mj{69m7p} 12 tiles
  \end{tabular}
  \caption{Differences between \mj{2p} drop and \mj{6p} drop from \ref{mj:fkgtt-6}}
  \label{tab:fkgtt-6}
\end{table}

Since it is difficult to calculate all the acceptance of all possible effective draws after dropping each possible tiles, we can just follow the principle to reduce the strongest complex taatsu to a ryanmen.


\subsubsection*{Three Pairs, Three Complex Taatsus, Two Ryanmen-Toitsus}

If there are two ryanmen-toitsus, then we reduce the remaining complex taatsu to a pair. After that, since we have three pairs in our hand, we will have 2 tiles of acceptance less than reducing the other taatsus to a ryanmen. Considering the acceptance when we reach 1-away, however, it becomes clear that this sacrifice of 2 tiles at 2-away is worth. If we choose to fix any ryanmen in \ref{mj:fkgtt-7}, completing any complex taatsus leaves us 1 surplus tile. On the contrary, by making the weakest complex taatsu our head, we are guaranteed perfect 1-away no matter what we draw, as shown in Table~\ref{tab:fkgtt-7}. Drop \mj{8s}, even though it may be a little bit counterintuitive because we have been hearing that three pairs are the weakest.

\tehai{12378m223677p899s}[][fkgtt-7]

\begin{table}[ht]
  \centering
  \begin{tabular}{lll}
    Drop & Next draw & Max acceptance \\
    \hline
    \mj{2p} & \mj{69m14p} & drop \mj{8s} $\to$ \mj{14578p9s} 20 tiles \\  
            & \mj{58p} & drop \mj{7p8s} $\to$ \mj{69m14p} 16 tiles \\
            & \mj{7p9s} & drop \mj{6p8s} $\to$ \mj{69m14p} 16 tiles \\
            & \mj{7s} & drop \mj{6p9s} $\to$ \mj{69m14p} 16 tiles \\
    \hline
    \mj{8s} & \mj{69m} & drop \mj{2p} $\to$ \mj{14578p9s} 20 tiles \\
            &          & drop \mj{7p} $\to$ \mj{12458p9s} 20 tiles \\
            & \mj{14p} & drop \mj{2p} $\to$ \mj{69m578p9s} 20 tiles \\
            & \mj{2p} & drop \mj{3p} $\to$ \mj{69m578p9s} 20 tiles \\
            & \mj{58p} & drop \mj{7p} $\to$ \mj{69m124p9s} 20 tiles \\
            & \mj{7p} & drop \mj{6p} $\to$ \mj{69m124p9s} 20 tiles \\
            & \mj{9s} & drop \mj{3p} $\to$ \mj{69m2578p} 20 tiles \\
            &         & drop \mj{6p} $\to$ \mj{69m1247p} 20 tiles
  \end{tabular}
  \caption{Differences between dropping \mj{2p} and \mj{8s} from \ref{mj:fkgtt-7}}
  \label{tab:fkgtt-7}
\end{table}


%--------------------------------------------
\subsubsection{Ryankanchan and Its Derived Shapes}

The shape \mj{246p} is called \textit{ryankanchan} in Japanese, sometimes \textit{ryankan} for short. The shape contains two kanchans \mj{24p} and \mj{46p}, and \textit{ryan} in Japanese means ``two'', hence the name. It takes up the space of 3 tiles and accepts at most 8 tiles, and we are left to chose any of its two kanchans when other parts of our hands get filled up, so ryankanchan is weaker than ryanmen in terms of space efficiency and final waits.

\subsubsection*{Comparing Ryankanchans}

Even between different kinds of ryankanchans, there are differences other than the tiles used. Table~\ref{tab:ryankanchan} lists three kinds of ryankanchan \mj{135p}, \mj{246p}, and \mj{357p}. Note that we omit \mj{468p} and \mj{579p} because they are identical to \mj{246p} and \mj{135p} due to symmetry of mahjong tiles. We can see that \mj{357p} has the most improvement draws, and dropping \mj{5p} from \mj{135p} makes our suji \mj{2p} wait better than non-suji \mj{2p} kanchan wait. \mj{246p} and the ryanmen \mj{67p} developed from it are all tanyao tiles, so it is a shape we would like to keep for tanyao.

\begin{table}[ht]
  \centering
  \begin{tabular}{lll}
    Shape & Ryanmen draw & Comment \\
    \hline
    \mj{135p} & \mj{6p} & Dropping \mj{5p} makes \mj{2p} wait suji \\
    \mj{246p} & \mj{7p} & Dropping \mj{2p} still keeps the acceptance of \mj{0p} \\
    \mj{357p} & \mj{28p} & 
  \end{tabular}
  \caption{Comparison between ryankanchans}
  \label{tab:ryankanchan}
\end{table}


\subsubsection*{Comparing Ryankanchan with Other Taatsus}

In the case where there are no floating tiles to drop and we must choose between ryankanchan and other taatsus, when we choose to drop reduce the ryankanchan to a simple kanchan, we lose 4 tiles of acceptance. Therefore, to destroy ryankanchan or not depends on the trade-off of the loss of 4 tiles of acceptance:

\begin{enumerate}
  \item If destroying ryankanchan still gives us the maximal acceptance, it is the right choice
  \item If there is no difference in terms of acceptance,
  \begin{enumerate}
    \item When pinfu is feasible, keep ryankanchan
    \item When pinfu is impossible, drop ryankanchan and set up suji trap
  \end{enumerate}
\end{enumerate}

We must drop a tile from either \mj{135m} or \mj{233345p} \ref{mj:rk-12}. Notice that the complex pinzu shape has an acceptance of \mj{1346p} 12 tiles. If we drop \mj{3p}, we lose the acceptance of \mj{369p} 7 tiles; if we drop \mj{2p}, we lose the acceptance of \mj{14p} 7 tiles. Either way, we lose more than dropping \mj{5m} (which loses \mj{4m} 4 tiles.) Therefore, we destroy ryankanchan in this case. Dropping \mj{1m} or \mj{5m} leaves us the same number of acceptance, but a \mj{5m} in our river\footnote{River, or \textit{kawa} in Japanese, refers to the discards of a player, neatly placed in front of them.} makes \mj{2m} more easily to come out, so \mj{5m} is the best choice here \citep[][pp.~44--45]{fkt}.

\tehai{135m 23334599p 789s}[Dora \mj{9p}][rk-12]

\begin{table}[ht]
  \centering
  \begin{tabular}{llll}
    Drop & \mj{5m} & \mj{3p} & \mj{2p} \\
    \hline
    Acceptance & \mj{2m13469p} 18 tiles & \mj{24m14p} 15 tiles & \mj{24m369p} 15 tiles
  \end{tabular}
  \caption{Acceptance of \ref{mj:rk-12}.}
  \label{tab:rk-12}
\end{table}

\ref{mj:rk-12a} is a hand similar to \ref{mj:rk-12}, but pinzu shape is not complex in a way that would lead to more loss when we drop a tile from it. Notice that, if we keep \mj{135m} ryankanchan, we have a 50-50 chance of getting pinfu by drawing \mj{24m} which makes our hand value 7700 (riichi pinfu dora dora.) If we drop \mj{5m}, only \mj{9p} draw (2 tiles out of the acceptance of 16 tiles total) gives us mangan tenpai; in other cases, we are stuck with 5200 bad wait. Therefore, although there is no difference in acceptance, in terms of expected value, \mj{3p} drop is optimal \citep[][pp.~46--47]{fkt}.

\tehai{135m 23356799p 789s}[Dora \mj{9p}][rk-12a]

\begin{table}[ht]
  \centering
  \begin{tabular}{lll}
    Drop & \mj{5m} & \mj{3p} \\
    \hline
    Acceptance & \mj{2m1349p} 16 tiles & \mj{24m14p} 16 tiles
  \end{tabular}
  \caption{Acceptance of \ref{mj:rk-12a}.}
  \label{tab:rk-12a}
\end{table}

\ref{mj:rk-13} has a \mj{8s} anko and thus pinfu is impossible. With \mj{9m} toitsu tanyao is also far away. Therefore, no matter what we drop, there is no difference in terms of acceptance and value. In this case, we drop \mj{5p} and make our \mj{2p} wait a suji trap \citep[][pp.23--25]{ishibashi-black}.

\tehai{23399m 135p 406888s}[South seat, turn 4, dora \mj{4s}][rk-13]

\ref{mj:rk-1} has a ryankanchan and an anko. If we drop \mj{8s} and make \mj{88s} our head, we will have only 16 tiles of acceptance at 1-away, as opposed to 24 tiles if we drop \mj{26m} while still maintaining 1-away, as shown in Table~\ref{tab:rk-1}. Besides the number of acceptance, when we keep \mj{456888s}, any draw of \mj{367s} improves our souzu shape, and any draw of manzu tile that gives us head leaves us ryanmen tenpai. The problem becomes which of \mj{26m} to drop. Considering the acceptance of \mj{0m} and the improvement draw \mj{7m}, drop \mj{2m} \citep[][p. 84]{uzaku-haikouritsu}.

\tehai{246m 23467p 456888s}[][rk-1]

\begin{table}[ht]
  \centering
  \begin{tabular}{llll}
    Drop & \mj{2m} & \mj{6m} & \mj{8s} \\
    \hline
    Acceptance & \mj{456m5678p} 24 tiles & \mj{234m5678p} 24 tiles & \mj{35m58p} 16 tiles
  \end{tabular}
  \caption{Acceptance of \ref{mj:rk-1}.}
  \label{tab:rk-1}
\end{table}

Since \ref{mj:rk-7} has two ryanmens with overlapping acceptance, dropping ryanmens and dropping any of \mj{268s} from \mj{24668s} will not lead to any difference in terms of the number of acceptance. However, considering the shape of our final waits, we can't touch ryanmens. We must keep \mj{6s} pair in case we draw any of \mj{469m} which gives us perfect 1-away. We also cannot overlook the possibility of 678 sanshoku, so we are incentivized to keep \mj{8s}. In this regard, \mj{2s} becomes the most unwanted tile. Drop \mj{2s} \citep[][p.~14]{uzaku-nankiru}.

\tehai{55678m 3467p 24668s}[East-1, East seat, turn 5, dora \mj{7p}][rk-7]

\begin{table}[ht]
  \centering
  \begin{tabular}{llll}
    Drop & \mj{2s} & \mj{6s} & \mj{8s} \\
    \hline
    Acceptance & \mj{5m258p567s} 24 tiles & \mj{258p357s} 24 tiles & \mj{5m258p356s} 24 tiles
  \end{tabular}
  \caption{Acceptance of \ref{mj:rk-7}.}
  \label{tab:rk-7}
\end{table}


\subsubsection*{When the 5 of 135 Ryankanchan is Akadora}

The draws that are related to the akadora of \mj{130p} ryankanchan are \mj{4567p}, which makes \mj{1p} seemingly useless. That said, when we draw \mj{2p}, \mj{0p} becomes floating and \mj{34567p} draws make \mj{0p} useful. Besides, when we draw another \mj{1p}, it produces a head and a kanchan; drawing \mj{6p} further develops it into a ryanmen. Therefore, we would like to keep it intact, unless we have enough blocks to aim for tenpai \citep{yuusee-13r5}.

For example, in \ref{mj:rk-9} since we only have \mj{345m} + \mj{688m} + \mj{130p} + \mj{1335s} four blocks, we must drop \mj{1z} and aim for draws that give us another block. If we draw \mj{2p}, as in \ref{mj:rk-9a}, then we can drop \mj{1s} and expect blocks around floating \mj{0p} and \mj{5s}.

\tehai{345688m 130p 1335s 1z}[East-1, turn 1, dora \mj{5z}][rk-9]

\tehai{345688m 1230p 1335s}[East-1, turn 2, dora \mj{5z}][rk-9a]

\ref{mj:rk-10} is a similar hand, but since we already have 5 blocks, we can drop \mj{1p} and go for tanyao.

\tehai{23445688m 130p 335s}[East-1, turn 2, dora \mj{5z}][rk-10]

The \mj{10m} shape in \ref{mj:rk-11} is the predecessor of \mj{130m} ryankanchan, so it will be a loss if we drop \mj{1m} over otakaze \mj{2z}. However, this possibility will not make \mj{1m} stronger than seat wind or yakuhai \mj{46z}. So we drop \mj{2z} for now, and drop \mj{1m} later on.

\tehai{10m 1389p 3666s 246z}[East-2, North seat, dora \mj{4m}][rk-11]


\subsubsection*{Vertically Developed Ryankanchans such as 2246 and 2466}

When we have \mj{1135p} and draw another copy of \mj{5p}, the result shape \mj{11355p} is a strong shape that accepts \mj{1245p} 4 kinds 12 tiles. Likewise for \mj{2246p} with \mj{6p} draw, and \mj{1355p} with \mj{1p} draw. If we draw both \mj{46p}, we end up the shape \mj{134556p}, which is called \textit{ryanmenkanchan}\footnote{See Section~\ref{subsec:6-tile} for more details on ryanmenkanchan.} \citep{chirno-kanchan}. Despite all these improvements, we still have to consider the cases where we have more stronger taatsus and must drop one from \mj{1135p} or \mj{2246p}.

We can drop either \mj{4p} or \mj{8p} from \ref{mj:rk-2}, which will lead to the same number of 16 tiles of acceptance, as shown in Table~\ref{tab:rk-2}. Keeping two pairs \mj{99m44p} in our head does not give us more ryanmen improvement draws than keeping \mj{468p} shape (both having only \mj{3p}), and drawing any of these pairs into an anko leaves us rii-nomi\footnote{\textit{Rii-nomi} refers to a state where we have no yaku other than riichi. The term is a portmanteau of riichi and \textit{nomi}, a Japanese word for ``only.'' Rii-nomi itself is not inherently bad; ryanmen rii-nomi is still better than not calling riichi especially when we are the first player to do so. Bad wait rii-nomi would require us to think twice before calling riichi and taking an amount of risk that is not justified by the cheap value. A part of hand building is about how to avoid rii-nomi without sacrificing efficiency too much.}. In that regard, it is slightly more advantageous to keep \mj{468p} shape and hope for pinfu. Drop \mj{4p} \citep[][p. 86]{uzaku-haikouritsu}.

\tehai{12399m 4468p 34789s}[][rk-2]

\begin{table}[ht]
  \centering
  \begin{tabular}{lll}
    Drop & \mj{4p} & \mj{8p} \\
    \hline
    Acceptance & \mj{57p25s} 16 tiles & \mj{9m45p25s} 16 tiles \\
    Pinfu draws & \mj{57p} 8 tiles & \mj{5p} 4 tiles
  \end{tabular}
  \caption{Acceptance of \ref{mj:rk-2}.}
  \label{tab:rk-2}
\end{table}

\ref{mj:rk-3} has \mj{77m} instead of \mj{99m} compared to \ref{mj:rk-2}, which increases the draws that give us perfect 1-away from \mj{3p} to \mj{68m3p}. While it is true that \mj{4p} drop still gives us the maximal 8 tiles of pinfu draw, having 12 tiles of perfect 1-away draw means that it is more likely to get good waits with \mj{8p} drop. Therefore, the best choice here is to drop \mj{8p} \citep[][p. 88]{uzaku-haikouritsu}.

\tehai{12377m 4468p 34789s}[][rk-3]

Dropping \mj{3s} and \mj{7s} from \ref{mj:3577-1} leaves us the same number of acceptance. However, if we drop \mj{3s} and have \mj{99m077s}, drawing \mj{9m7s} means that \mj{0s} is useless. Since we have a \mj{68p} kanchan in our hand, there is no pinfu or even perfect 1-away either. If we are going to riichi anyway, it is best to make full use of \mj{0s} to achieve maximal expected value. Thus drop \mj{7s} \citep[][p.~26]{uzaku-nankiru}.

\tehai{99m 23468p 3077s 444z}[East-1, East seat, turn 7, dora \mj{2p}][3577-1]

Since we have a \mj{77z} pair in \ref{mj:3577-2}, when we call pon on \mj{7z} we must have a pair elsewhere. In that regard, we must keep \mj{77s} pair and drop \mj{3s}. It is likely that \mj{0s} becomes useless, but the difference is only 1000 points. After we open our hand, we go full steam ahead \citep[][p.~26]{uzaku-nankiru}.

\tehai{666789m 34p 3077s 77z}[East-1, West seat, turn 6, dora \mj{2m}][3577-2]

We must drop either \mj{4s} or \mj{8s} from the shape \mj{4468s} from \ref{mj:3577-3}. Since we already have pairs \mj{77m77p}, \mj{8s} drop loses 2 more tiles than \mj{4s} drop. When our hand becomes 1-away with, say, \mj{6m} draw, there is no difference in terms of acceptance (ryanmen + pair + kanchan-toitsu versus ryanmen + head + ryankanchan), but if we have \mj{468s} shape in our hand, we have a higher chance of getting pinfu than \mj{446s} shape. Therefore, based on the acceptance at 2-away and pinfu chance at 1-away, the best choice is \mj{4s} drop \citep{zundamon-3head}.

\tehai{23577m 12377p 4468s}[East-1, West seat, turn 6, dora \mj{1p}][3577-3]


\subsubsection*{Bulging Ryankanchans such as 1335 and 2446}

To handle shapes like \mj{1335p} and \mj{2446p}:

\begin{itemize}
  \item 2-away
  \begin{itemize}
    \item Drop the outside tile, aim for ryanmen improvement
  \end{itemize}
  \item 1-away
  \begin{itemize}
    \item If pinfu is within sight, drop the bulging part
    \item If pinfu is infeasible, drop the inside tile
  \end{itemize}
\end{itemize}

Suppose that a part of our hand is \mj{77m1335p} and the rest of our hand being completed blocks and a ryanmen. Table~\ref{tab:rk-1335} lists the pinfu draws, ryanmen improvement draws, and final waits (when the ryanmen is completed) of different drops from \mj{1335p}. When we reduce \mj{1335p} to a regular ryankanchan, it has the most draws that give us pinfu; if we cannot draw \mj{24p} in time, we can still cut \mj{5p} and riichi and make our \mj{2p} wait a suji. When pinfu is not feasible, we should cut \mj{5p} in advance to make our suji wait \mj{2p} more likely to come out\footnote{The argument that the suji of a tile cut in advance is more likely to come out than the suji of a riichi declaration tile is based on the deal-in rate difference of sujis. Sujis of riichi declaration tiles have higher deal-in rates than other kinds of sujis \citep[][pp.~152--153]{miinin-stat}. If a player's policy is to cut the suji with lower deal-in rate, they will tend not to cut riichi declaration tile suji.}. When our hand is callable, the \mj{3p} pair becomes more valuable.

\begin{table}[ht]
  \centering
  \begin{tabular}{llll}
    Drop & \mj{1p} & \mj{3p} & \mj{5p} \\
    \hline
    Result shape & \mj{335p} & \mj{135p} & \mj{133p} \\
    Pinfu draws & \mj{4p} 4 tiles & \mj{24p} 8 tiles & \mj{2p} 4 tiles \\
    Ryanmen improvement draws & \mj{68m6p} 12 tiles & \mj{6p} 4 tiles & \mj{68m} 8 tiles \\
    Final waits & \mj{7m3p} shanpon & \mj{2p} kanchan & \minibox{\mj{2p} kanchan or\\\mj{7m3p} shanpon}
  \end{tabular}
  \caption{Comparing drops from the shape \mj{77m1335p}.}
  \label{tab:rk-1335}
\end{table}

\ref{mj:rk-1335a} has pinfu chance if we drop \mj{3p} and draw \mj{24p}. Meanwhile, \ref{mj:rk-1335b} cannot have pinfu no matter what we drop; in this case, drop \mj{5p} and make \mj{2p} more likely to come out. \ref{mj:rk-1335c} has no pinfu either, and it has \mj{555z} anko, which means that, after dropping \mj{5p}, calling pon on both \mj{7m3p} will give us tenpai \citep{yuusee-1335}.

\tehai{12377m 1335p 45789s}[East-1, East seat, turn 5, dora \mj{9s}][rk-1335a]

\tehai{77m 1335p 45789s 333z}[East-1, East seat, turn 5, dora \mj{9s}][rk-1335b]

\tehai{77m 1335p 45789s 555z}[East-1, East seat, turn 5, dora \mj{9s}][rk-1335c]


\subsubsection*{Horizontally Developed Ryankanchans such as 2357}

\textbf{Drop the kanchan part; generally, drop 5 first; when we have no head, drop 7} \citep{yuusee-2357}. Even if tanyao is within reach, keeping ryankanchan takes up the space which we could use for safe tiles instead. It is also possible that our ryankanchan ends up a kanchan bad wait, which leaves us no reason to not destroy it when we have a ryanmen.

When we add \mj{4m} to \mj{2357m}, since \mj{23m} also accepts \mj{4m}, we actually end up with a mentsu and a kanchan. Compared to other inside kanchans\footnote{Inside kanchans are kanchans with 8 tiles of ryanmen improvement draws. For pinzu suit, for example, that would be \mjfn{35p}, \mjfn{46p}, and \mjfn{57p}.}, this \mj{57m} kanchan is inferior. In this regard, drop \mj{57m} and keep \mj{46p} in \ref{mj:rk-2357a}.

\tehai{2357m 46p 22256788s}[East-1, turn 3, dora \mj{4z}][rk-2357a]

In addition to ryanmen improvement draws, considering final waits, the kanchan part of \mj{2357m} is inferior to outside kanchans\footnote{Outside kanchans are kanchans with 4 tiles of ryanmen improvement draws. For pinzu suit, for example, that would be \mjfn{13p}, \mjfn{24p}, \mjfn{68p} and \mjfn{79p}.}. In this regard, drop \mj{57m} and keep \mj{79s} in \ref{mj:rk-2357b}.

\tehai{2357m 2225699p 79s}[East-1, turn 3, dora \mj{4z}][rk-2357b]

When we drop \mj{7p} from \mj{2357p}, the \mj{235p} shape produces a head if we draw any of \mj{25p}. Drawing \mj{4p} gives us a yonrenkei shape\footnote{See Section \ref{subsub:yonrenkei}.}, which also produces a head easily. Since \ref{mj:rk-2357c} has no head, drop \mj{7m}.

\tehai{2357m 234678p 3456s}[East-1, turn 3, dora \mj{4z}][rk-2357c]


\subsubsection*{Ryankanchans with Following Ryanmens such as 34468}

The \mj{34468p} shape produces a mentsu if we draw any of \mj{2457p}, as shown in Table~\ref{tab:rk-34468}, and therefore is a strong shape \citep{clearrain-34468, yuusee-34468}.

\begin{table}[ht]
  \centering
  \begin{tabular}{lll}
    Draw & Result shape & Composition \\
    \hline
    \mj{2p} & \mj{234468p} & Shuntsu \mj{234p} + ryankanchan \mj{468p} \\
    \mj{4p} & \mj{344468p} & Anko \mj{444p} + kanchan \mj{68p} \\
            &              & or ryanmen-toitsu \mj{344p} + ryankanchan \mj{468p} \\
    \mj{5p} & \mj{344568p} & Ryanmenkanchan (see \ref{subsec:6-tile}) \\
    \mj{7p} & \mj{344678p} & Ryanmen-toitsu \mj{344p} + shuntsu \mj{678p}
  \end{tabular}
  \caption{Draws that produce a mentsu out of \mj{34468p} shape.}
  \label{tab:rk-34468}
\end{table}

Since dropping any tile from \mj{34468p} annihilates the benefits mentioned earlier, when compared to yonrenkei or nakabukure\footnote{See Section \ref{subsub:nakabukure}.}, we tend to keep \mj{34468p} and reduce other 4-tile shapes to a mentsu. As shown in Table~\ref{tab:rk-34468a}, dropping either of \mj{36p} also gives us far more acceptance than dropping any tile from \mj{34468m}. Therefore, drop \mj{3p} from \ref{mj:rk-34468a}.

\tehai{34468m 3456p 11789s}[][rk-34468a]

\begin{table}[ht]
  \centering
  \begin{tabular}{llll}
    Drop & \mj{4m} & \mj{8m} & \mj{3p} \\
    \hline
    Acceptance & \mj{257m} 12 tiles & \mj{25m} 8 tiles & \mj{2457m1s} 16 tiles
  \end{tabular}
  \caption{Acceptance of \ref{mj:rk-34468a}.}
  \label{tab:rk-34468a}
\end{table}

When we have 6 blocks, we still tend to keep \mj{34468p} shape. When there are only ryanmens in the rest of our hand, drop \mj{8p}. In this regard, since there is a weaker \mj{79s} kanchan in \ref{mj:rk-34468b}, we drop \mj{79s}; since there is no weaker taatsu in \ref{mj:rk-34468c}, drop \mj{8m}.

\tehai{34468m 45789p 2279s}[][rk-34468b]

\tehai{34468m 45789p 2399s}[][rk-34468c]

When we have 5 blocks, we still tend to keep \mj{34468p} shape. When there is no head in our hand, how we handle \mj{34468p} shape will largely depend on table situation and turn number. When we are still early in the game, we can drop \mj{8p} and focus on \mj{34p} ryanmen. When we are late and there is no room for improvement, make \mj{44p} our head and drop \mj{3p}. Therefore, even for the same hand such as \ref{mj:rk-34468d}, our decision changes depending on other information.

\tehai{34468m 123789p 357s}[][rk-34468d]

\textbf{Dropping \mj{5m} from \mj{13556m} in advance at 1-away is not recommended} \citep{yuusee-13556}. When we drop \mj{5m} from \ref{mj:rk-34468e}, our ryanmen-tenpai draw reduces to only \mj{2m}, opposed to \mj{24m} if we keep it, while both options have the same number of bad-wait draws, as shown in Table~\ref{tab:rk-34468e}. Keeping a safe tile here worsens our wait quality too much.

\tehai{13556m 33789p 123s 4z}[East-1, West seat, turn 6, dora \mj{8p}][rk-34468e]

\begin{table}[ht]
  \centering
  \begin{tabular}{lll}
    Drop & \mj{5m} & \mj{4z} \\
    \hline
    Result shape & \mj{1356m33p4z} & \mj{13556m33p} \\
    Ryanmen-tenpai draws & \mj{2m} 4 tiles & \mj{24m} 8 tiles \\
    Bad-wait draws & \mj{47m} 8 tiles & \mj{57m3p} 8 tiles
  \end{tabular}
  \caption{Acceptance of \ref{mj:rk-34468e}.}
  \label{tab:rk-34468e}
\end{table}

That said, we can drop the ryankanchan part when we draw \mj{6m}. Even though the shape \mj{5566m} has overlapping acceptance \mj{47m}, our bad-wait draws reduce to \mj{56m3p} 3 kinds 6 tiles, 2 tiles less than keeping \mj{13556m}, as shown in Table~\ref{tab:rk-34468f}. We also gain space for safe tiles. Therefore, drop \mj{3m} and then \mj{1m} from \ref{mj:rk-34468f}.

\tehai{135566m 33789p 123s}[East-1, West seat, turn 6, dora \mj{8p}][rk-34468f]

\begin{table}[ht]
  \centering
  \begin{tabular}{lll}
    Drop & \mj{3m} & \mj{6m} \\
    \hline
    Result shape & \mj{15566m33p} & \mj{13556m33p} \\
    Ryanmen-tenpai draw & \mj{47m} 8 tiles & \mj{24m} 8 tiles \\
    Bad-wait draw & \mj{56m3p} 6 tiles & \mj{57m3p} 8 tiles
  \end{tabular}
  \caption{Acceptance of \ref{mj:rk-34468f}.}
  \label{tab:rk-34468f}
\end{table}

The same principle applies to inner \mj{34468s} shape \citep{yuusee-34468}. Drop \mj{6s} from \ref{mj:rk-34468fa} and \ref{mj:rk-34468fb}.

\tehai{444m 33p 678p 334468s}[East-1, West seat, turn 4, dora \mj{9p}][rk-34468fa]

\tehai{11 789p 334468s 444z}[East-1, West seat, turn 4, dora \mj{9p}][rk-34468fb]

%Another case where we destroy \mj{13556m} shape at 1-away is when we have a 3--7 pair elsewhere and we draw \mj{1m}. When we draw \mj{1m} and drop \mj{5m}, we end up a kanchan-toitsu, a ryanmen, and a pair. This gives us the same number of acceptance, but the number of ryanmen-improvement draws increases, as shown in Table~\ref{tab:rk-34468g}. Therefore, drop \mj{5m} from \ref{mj:rk-34468g}.
%
%\tehai{113556m 33789p 123s}[East-1, West seat, turn 6, dora \mj{8p}][rk-34468g]
%
%\begin{table}[ht]
%  \centering
%  \begin{tabular}{lll}
%    Drop & \mj{5m} & \mj{1m} \\
%    \hline
%    Result shape & \mj{11356m33p} & \mj{13556m33p} \\
%    Ryanmen-tenpai draws & \sfrac{2}{8}: \mj{12m3p} & \sfrac{2}{8}: \mj{24m} \\
%    Bad-wait draws & \sfrac{2}{8}: \mj{47m} & \sfrac{2}{8}: \mj{47m} \\
%    Ryanmen-improvement draws & \sfrac{2}{8}: \mj{24p} & \sfrac{1}{3}: \mj{6m}
%  \end{tabular}
%  \caption{Acceptance of \ref{mj:rk-34468g}.}
%  \label{tab:rk-34468g}
%\end{table}

When we have yakuhai anko such as \ref{mj:rk-34468h} and draw \mj{1m} with \mj{13556m} shape in our hand, we can drop \mj{5m} and keep \mj{113m} so that when we pon either \mj{1m1p} or draw \mj{2m}, we have a good ryanmen wait. Drop \mj{5m} from \ref{mj:rk-34468h}.

\tehai{113556m 11789p 555z}[East-1, West seat, turn 6, dora \mj{8p}][rk-34468h]


\subsubsection*{Detached Ryankanchan such as 245679}

The \mj{245679p} shape in \ref{mj:rk-4} is dubbed ``detached ryankanchan'' since it accepts \mj{38p} far apart. If we draw \mj{56p}, we can further improve it into a ryanmenkanchan. On the other hand, dropping either \mj{29p} will leave us a mentsu and kanchan; which to drop largely depends on the rest of our hand.

In \ref{mj:rk-4}, dropping either of \mj{29p} from detached ryankanchan or dropping \mj{6m} from \mj{566m} will leave us the same number of acceptance at 1-away. However, considering the draws that will give us pinfu and the sanshoku chance, fixing \mj{56m} ryanmen here is better, as shown in Table~\ref{tab:rk-4} \citep[][p. 89]{uzaku-haikouritsu}.

\tehai{566m 245679p 45699s}[][rk-4]

\begin{table}[ht]
  \centering
  \begin{tabular}{llll}
    Drop & \mj{6m} & \mj{2p} & \mj{9p} \\
    \hline
    Acceptance & \mj{47m38p} 16 tiles & \mj{467m8p9s} 16 tiles & \mj{467m3p9s} 16 tiles \\
    Pinfu draws & \mj{38p} 8 tiles & \mj{8p} 4 tiles & \mj{3p} 4 tiles \\
    Sanshoku-chance draws & \mj{4m8p} & \mj{4m8p} & None
  \end{tabular}
  \caption{Acceptance of \ref{mj:rk-4}.}
  \label{tab:rk-4}
\end{table}


\subsubsection*{Choosing between two ryankanchans}

\ref{mj:rk-5} has two ryankanchans, one bulging and one vertically developed on one end. Since there is no pair in the rest of our hand, we will focus on how we can make tiles following \mj{55p} and \mj{66s} pair function. Dropping \mj{2s} makes \mj{4s} useless since completing \mj{466s} shape leaves us no pair\footnote{Technically, there is, if we treat it as \mj{45s} + \mj{66s}, but we are then required to drop either of \mj{37p} and we are still not tenpai yet despite \ref{mj:rk-5} already at 1-away.}. Dropping either of \mj{357p} will make full use of all tiles while giving us the same number of acceptance. Taking ryanmen improvement into consideration, \mj{5p} drop becomes optimal, as shown in Table~\ref{tab:rk-5}.

\tehai{123789m 3557p 2466s}[][rk-5]

\begin{table}[ht]
  \centering
  \begin{tabular}{llll}
    Drop & \mj{3p} & \mj{5p} & \mj{2s} \\
    \hline
    Acceptance & \mj{56p36s} 12 tiles & \mj{46p3s} 12 tiles & \mj{46p} 8 tiles \\
    Ryanmen-improvement draws & \mj{8p5s} 8 tiles & \mj{28p5s} 12 tiles & \mj{28p} 8 tiles
  \end{tabular}
  \caption{Acceptance and improvement of \ref{mj:rk-5}.}
  \label{tab:rk-5}
\end{table}

When we must choose between two simple ryankanchans, drop the tile that will make our wait better \citep[][p.~26]{feng-intermediate} Thus, we drop \mj{5p} from \ref{mj:rk-6} and make our wait suji. Note that while \mj{6m} drop also makes \mj{3m} wait suji, \mj{3m} is more inside than \mj{8p} and thus still has an overall lower win rate.

\tehai{246999m 33567s 579p}[][rk-6]

% NEED SOURCE
%\subsubsection*{リャンカンのそばに暗刻}
%
%Dropping \mj{1s} from \ref{mj:rk-8} seems to leave us a ryankanchan \mj{307s} and \mj{77s} head that makes full use of \mj{0s}. However, drawing \mj{3s} renders \mj{0s} useless and ends up rii-nomi. Drop \mj{68p}?
%
%\tehai{23468p 130777s 444z}[東1局　東家　7巡目][rk-8]
%
%\begin{table}[ht]
%  \centering
%  \begin{tabular}{ll}
%    切る牌 & 受け入れ \\
%    \hline
%    \mj{1s} & \sfrac{7}{24}：\mj{678p3456s} \\
%    \mj{8p} & \sfrac{6}{21}：\mj{6p12456s}
%  \end{tabular}
%  \caption{手牌\ref{mj:rk-8}の受け入れ}
%  \label{tab:rk-8}
%\end{table}

%--------------------------------------------
\subsubsection{134 Shape}

It is a common pitfall to only see the ryanmen part of a \mj{134p} shape and drop \mj{1p} early. The possibilities of \mj{134p} can help us tremendously especially when it is still early and our hand lacks shapes \citep{yoteru-134}. In short,

\begin{itemize}
\item Makes our hand \textit{kuttsuki}\footnote{\textit{Kuttsuki} is a Japanese term whose literal meaning is ``sticky,'' In mahjong context, it refers to a state where tiles next to or 1 tile away from floating tiles all become accepted.}
\item Easy to produce a pair
\end{itemize}

\begin{table}[ht]
  \centering
  \begin{tabular}{llp{12cm}}
    Draw & Shape & Comment \\
    \hline
    \mj{1m} & \mj{1134m} & A pair and a ryanmen. \\
    \mj{2m} & \mj{1234m} & A yonrenkei, which produces a pair easily or functions as nobetan when we reach tenpai. \\
    \mj{3m} & \mj{1334m} & A kanchan and a ryanmen. \\
    \mj{4m} & \mj{1344m} & A kanchan and a pair. Decent as our final wait. \\
    \mj{5m} & \mj{1345m} & A floating \mj{1m} and a mentsu. However, drawing \mj{2m} gives us \mj{45m} ryanmen, drawing \mj{3m} also gives us a decent kanchan. Drawing \mj{6m} and we have the \mj{13456m} shape which can easily further develop into a mentsu and a ryanmen. \\
    \mj{6m} & \mj{1346m} & Two kanchans.
  \end{tabular}
  \caption{Development of \mj{134m}.}
  \label{tab:134}
\end{table}

Therefore, when our hand does not have 5 blocks ready, we keep \mj{134p} over single yakuhais. When there is no head in our hand, we also prioritize \mj{134p}. Otherwise, dropping \mj{1p} is fine.

\ref{mj:134-1} does not have enough blocks, so we keep \mj{134m} and expect it to develop. Drop \mj{56z} from \ref{mj:134-1}.

\tehai{134m 340899p 248s 56z}[East-1, West seat][134-1]

\ref{mj:134-2} already has 5 blocks to aim for tenpai, so we can drop \mj{1m} and keep \mj{56z} as safe tiles.

\tehai{13488m 3468p 345s 56z}[East-1, West seat][134-2]

\ref{mj:134-3} has no head confirmed, so we keep \mj{134m} to increase our chances of getting a head. Drop \mj{2z}.

\tehai{134m 1236789p 888s 2z}[East-1, West seat, dora \mj{8s}][134-3]


\subsubsection*{Play Defensively when Our Hand is Too Messed Up}

\ref{mj:134-4} is 5-away and only has one ryanmen \mj{34m}. No matter how good \mj{134m} can develop, it is more likely that some one will slam riichi earlier than us. In this case, we drop \mj{1m} and keep safe tiles \mj{567z}.

\tehai{1348m 4899p 136s 567z}[][134-4]


\subsubsection*{Avoid Bad-wait Rii-Nomi}

\ref{mj:134-3a} has the same tiles as \ref{mj:134-3} but dora is \mj{7z} instead of \mj{8s}. The possible kanchan that \mj{134m} brings us becomes unfavorable, so we drop \mj{1m} and fix ryanmen.

\tehai{134m 1236789p 888s 2z}[East-1, West seat, dora \mj{7z}][134-3a]


\subsubsection*{When the 1 is Also a Dora}

As shown in Table~\ref{tab:134-dora}, when we have no akadora in \ref{mj:134}, keeping dora \mj{1p} or taking \mj{34p} ryanmen does not make too much difference in terms of hand value. Hence drop \mj{1p} and take ryanmen to increase win rate. When we have one or more akadoras, keeping dora \mj{1p}, and possibly calling pon \mj{1s} and pivoting to \mj{1p} tanki or drawing another \mj{1p} and waiting on shanpon greatly increase hand value. When it is still early, with akadoras in our hand, it is better to drop \mj{4p} \citep{yuusee-134}.

\tehai{?67m 134p 11?67s -7'77z}[East-1, West seat, dora \mj{1p}][134]

\begin{table}[ht]
  \centering
  \begin{tabular}{lll}
    \mj{?} & Drop \mj{1p} & Drop \mj{4p} \\
    \hline
    \mj{5m} + \mj{5s} & 1000 points: \mj{25p} 7 tiles; 2000 points: \mj{0p} 1 tile & 2000 points \\
    \mj{5m} + \mj{0s} & 2000 points: \mj{25p} 7 tiles; 3900 points: \mj{0p} 1 tile & 3900 points \\
    \mj{0m} + \mj{0s} & 3900 points: \mj{25p} 7 tiles; 7700 points: \mj{0p} 1 tile & 7700 points
  \end{tabular}
  \caption{Value of \ref{mj:134} depending on drop and draw.}
  \label{tab:134-dora}
\end{table}


%--------------------------------------------
\subsubsection{When the 5 of 578 Shape is an Akadora}

\begin{table}[ht]
  \centering
  \begin{tabular}{ll}
    Hand state & Action \\
    \hline
    Far away from tenpai & Keep \mj{078m} \\
    1-away & Drop \mj{0m} if tanyao is infeasible \\
    Tenpai & Generally drop \mj{0m} \\
  \end{tabular}
  \caption{Handing \mj{078m}.}
\end{table}

When we are still far away from tenpai, we can expect to draw tiles around \mj{0m}. Since \ref{mj:5-1} has the risk of ending up no dora bad-wait rii-nomi, we drop \mj{7s} from \mj{778s}. \ref{mj:5-2} also has the same risk, and \mj{557p} has more ryanmen-improvement draws than \mj{224s}, so we drop \mj{4s} \citep{yuusee-578}.

\tehai{078m 23557p 224 778s}[][5-1]

\tehai{078m 123557p 22478s}[][5-2]

When we are 1-away, acceptance and final wait become more important than the value the dora brings. \ref{mj:5-3} already has \mj{5z} anko and we do not have to cling to \mj{0m} for value. Drop \mj{0m}

\tehai{078m 22456p 688s 555z}[][5-3]

However, as an exception, when keeping \mj{07m} gives us sanshoku or tanyao which enables us to call and speed up our hand, we drop \mj{8m}. When reaching tenpai, if dropping \mj{8m} guarantees sanshoku and tanyao, we can also dama our hand instead of dropping \mj{0m} and riichi. In this regard, \ref{mj:5-4} is at tenpai but does not have sanshoku confirmed, so we drop \mj{0m} and riichi; on the other hand, \ref{mj:5-5} has confirmed sanshoku, and its value is already mangan, we drop \mj{8m} without riichi.

\tehai{078m 22456p 22256s -7s}[][5-4]

\tehai{078m 22567p 22256s -7s}[][5-5]


%--------------------------------------------
\subsubsection{3568 Shape}

It is easy to get a head out of \mj{3568p} shape \citep{yuusee-3568}. Not just stacking \mj{38p}, but drawing \mj{47p} also gives us yonrenkei. \mj{3p} can also function as a floater. Thus, it is generally treated as 2 blocks. Table~\ref{tab:3568} lists our policy when taking the rest of our hand into consideration.

\begin{table}
  \centering
  \begin{tabular}{lll}
    Number of block & Compared with & Action \\
    \hline\hline
    4 & & Keep \mj{3568p} \\
    \hline
    5 & Floaters (no head in hand) & Keep \mj{3568p} \\
      & Floater tiles other than 34567 & Keep \mj{3568p} \\
      & Penchan (bottleneck of tanyao) & Keep \mj{3568p} \\
      & Penchan (not bottleneck of tanyao) & Drop \mj{3568p} \\
      & Floaters tiles of 34567 & Drop \mj{3568p} \\
    \hline
    6 & & Drop \mj{3568p} \\
    \hline\hline
  \end{tabular}
  \caption{Handling \mj{3568p} or \mj{2457p} \citep{yuusee-3568, clearrain-3568}}
  \label{tab:3568}
\end{table}

\ref{mj:3568-1} already has \mj{5z} anko, and directly accepting \mj{3m} using \mj{12m} penchan is valuable. Thus, we drop \mj{8p} first, and then drop \mj{3p} if it does not give us ryanmen.

\tehai{12m 3568p 22788s 555z}[East-1, West seat, turn 3, dora \mj{1s}][3568-1]

\ref{mj:3568-2}, however, gains tanyao if \mj{12m} is removed. Therefore, we drop \mj{1m} and expect to draw \mj{8p} or tiles around \mj{3p}.

\tehai{12m 3568p 22455788s}[East-1, West seat, turn 3, \mj{1s}][3568-2]

Compared to 3--7 floaters, generally, we drop \mj{8p} from \mj{3568p} since it has 4 less tiles of acceptance as a floater. However, for hands like \ref{mj:3568-3} where there is no confirmed head, we drop \mj{7m} and expect to find a head in shapes developed from \mj{3568p}.

\tehai{237m 3568p 1236778s}[East-1, West seat, turn 3, dora \mj{1s}][3568-3]

There are 6 blocks in \ref{mj:3568-4}, which makes us choose between \mj{2m89p} to drop. If we simply treat \mj{2457m} as two kanchans, we are tempted to drop the weakest taatsu, which is \mj{89p} penchan. However, dropping penchan has the following demerits:

\begin{itemize}
\item \mj{45m} cannot function as ryanmen (as long as we do not draw \mj{278m})
\item We lose \mj{7p} acceptance
\end{itemize}

On the other hand, if we keep \mj{457m89p} shape, we will not lose \mj{3m} acceptance because we still have \mj{45m} ryanmen, and drawing \mj{7p} completes a block. The merits of keeping \mj{2m} here, which is to increase the probability of confirming a head, becomes irrelevant due to \mj{11p} pair, so we can drop \mj{2m} without much loss.

\tehai{2457m 1189p 667s 444z}[][3568-4]


%--------------------------------------------
\subsubsection{35689 Shape}

In situations where we are choosing between a 3--7 floater and \mj{35689p} shape, drop the latter, but not its penchan part. \textbf{Generally, drop \mj{3p} and keep penchan at 1-away} \citep{yuusee-35689}. With \mj{35689p}, riichi chance on \mj{7p} draw is more valuable than forming a taatsu around \mj{3p}. Besides, this floater \mj{3p} is inferior to an isolated floater 3 because:

\begin{itemize}
\item Drawing \mj{2p} produces a ryanmen with overlapping acceptance
\item Drawing \mj{4p} does not give us a ryanmen; it produces a mentsu \mj{345p} and a kanchan \mj{68p}
\end{itemize}

While it is true that \mj{34568p} has \mj{245p} draws that give us ryanmen, from dropping \mj{9p} to drawing any of \mj{245p} upgrade will inevitably take a couple of turns.

That said, there are two cases where we can drop penchan \mj{89p} and go back to 2-away after dropping \mj{3p}:

\begin{enumerate}
\item Penchan is the bottleneck of tanyao
\item We have other strong shapes such as yonrenkei and nakabukure
\end{enumerate}

When we are comparing between floaters in \ref{mj:35689-1} at 2-away\footnote{While generally we take penchan over floaters at 1-away, when we are 2-away instead, we have some leeway to seek taatsu around floaters and hope for better shapes. Indeed \citet{yuusee-389} states that we can drop penchan if we have two 3--7 floaters at 2-away. However, as stated previously, the \mjfn{3p} floater is actually weaker than an isolated 3--7 floater, so, in the case of \mjfn{35689p}, we do not really meet the standard of dropping penchan.}, \mj{3s} from \mj{35689s} is inferior to isolated \mj{6p}. Therefore, drop \mj{3s}.

\tehai{22378m 1136p 35689s}[East-1, West seat, turn 5, dora \mj{8m}][35689-1]

Likewise, for \ref{mj:35689-2} which is 1-away, we keep \mj{89s} penchan and cut \mj{3s} floater.

\tehai{234789m 115p 35689s}[East-1, West seat, turn 5, dora \mj{8m}][35689-2]

If we later draw \mj{6m} and get yonrenkei \mj{6789m} in \ref{mj:35689-3}, we can drop \mj{89s} penchan. Since this penchan still has overlapping acceptance, when we have very strong shapes in the rest of our hand, we can opt to drop penchan at 1-away.

\tehai{2346789m 115p 5689s}[East-1, West seat, turn 6, dora \mj{8m}][35689-3]

If this \mj{89s} penchan is the only bottleneck of tanyao, we can drop it in exchange of value and flexibility even if it increases our shanten number. Drop \mj{9s} from \ref{mj:35689-4}.

\tehai{22367m 2247p 35689s}[East-1, West seat, turn 5, dora \mj{8m}][35689-4]


%--------------------------------------------
\subsubsection{Sankanchan and Yonkanchan}

Shapes like \mj{1357p} are dubbed \textit{sankanchan} and \mj{13579p} \textit{yonkanchan} in Japanese, because they have three kanchans and four kanchans respectively. In general, the order of importance is ``even-numberred sankanchan \textless{} floaters \textless{} odd-numbered sankanchan \textless{} outside kanchan \textless{} yonkanchan \textless{} inside kanchan'' \citep{chirno-kanchan}.

For odd-numbered sankanchan like \mj{1357p}, drawing \mj{17p} produces a ryankanchan and a head, drawing \mj{9p} produces a yonkanchan, drawing \mj{35p} produces a ryankanchan and a kanchan, drawing \mj{8p} produces a ryanmen, and drawing \mj{246p} completes a mentsu. Therefore, the acceptance of \mj{1357p} is \mj{123456789p}. However, if we must choose between an outside kanchan and a sankanchan, dropping \mj{1p} from \mj{1357p} while keeping the kanchan gives us the same number of acceptance but more ryanmen-improvement draws, as shown in Table~\ref{tab:1357}. Hence the conclusion that sankanchans are weaker than outside kanchans.

\begin{table}
  \centering
  \begin{tabular}{lll}
    Drop & \mj{1p} & \mj{13s} \\
    \hline
    Acceptance & \mj{46p2s} 12 tiles & \mj{246p} 12 tiles \\
    Ryanmen-improvement draws & (\mj{2p})\tablefootnote{Since we just cut \mjfn{1p}, this could make us end up furiten.}\mj{8p4s} & \mj{8p}
  \end{tabular}
  \caption{Acceptance of \mj{1357p13s}.}
  \label{tab:1357}
\end{table}

Yonkanchans like \mj{13579p} accept \mj{2468p} but lack any ryanmen improvement. When we choose between yonkanchan and an inside kanchan, we can reduce yonkanchan to a sankanchan without losing acceptance while increasing our ryanmen improvement draws.

Even-numbered sankanchans such as \mj{2468m} also lack any ryanmen improvement. In most cases, we reduce it to a ryankanchan. The only case where we would keep it as is is when we are aiming for tanyao and treat it as \mj{24m} and \mj{68m} two kanchans. Therefore, we drop \mj{8s} from \mj{2468s} in \ref{mj:3-1} for the faint possibility of 234 sanshoku, and we drop \mj{11m} pair from \ref{mj:3-2} to aim for tanyao \citep{yuusee-2468}.

\tehai{24m 44699p 2468s 777z}[][3-1]

\tehai{11357m 46688p 2468s}[][3-2]


%--------------------------------------------
\subsubsection{Shapes with Overlapping Acceptance} \label{subsub:nidouke}

When two taatsus share the same tiles as their acceptance, our effective number of acceptance reduces. This is known as \textit{nidouke} in Japanese, which literally means ``to accept twice.'' Shapes like \mj{3467p} shares \mj{5p} acceptance, so only \mj{258p} 3 kinds 12 tiles complete a mentsu out of this shape, whereas \mj{2569s} 4 kinds 16 tiles complete a mentsu out of \mj{3478s} shape. Shapes with overlapping acceptance are not always bad or weak; there are cases where we still keep them over other weak taatsus.


\subsubsection*{Neighboring Pairs}

Shapes like \mj{3344p} can be treated as two neighboring pairs \mj{33p} and \mj{44p} or two ryanmens \mj{34p} and \mj{34p}, therefore shapes with overlapping acceptance. We can easily produce a ryanmen out of this shape, so when pinfu is within sight, its priority is higher than weak taatsus such as kanchans, but lower than other pairs and ryanmens. Therefore, we drop one tile from \mj{3344p} in \ref{mj:narabtt-1}, and we drop \mj{68s} kanchan from \ref{mj:narabtt-2} \citep[][pp. 112--113]{uzaku-haikouritsu}.

\tehai{33678m 3344p 23478s}[][narabtt-1]

\tehai{33678m 3344p 23468s}[][narabtt-2]

Since \ref{mj:narabtt-3} has no possibility of pinfu, and dropping \mj{99p} gives us tanyao, it becomes optimal to drop \mj{9p}.

\tehai{33678m 8899p 23468s}[][narabtt-3]


\subsubsection*{Skipping Pairs}

Shapes like \mj{7799p} can be treated as two pairs \mj{77p} and \mj{99p} with a gap in between or two kanchans \mj{79p} and \mj{79p}, therefore shapes with overlapping acceptance. Since \mj{8p} draw gives us iipeikou chance, and tiles around the pair on the inside gives us perfect 1-away, we still keep it over other kanchans and penchans, and destroy it when the rest of our hand are completed mentsus, pairs, and ryanmens \citep[][pp. 114--117]{uzaku-haikouritsu}.

We can drop either of \mj{79m} from \ref{mj:tobitt-1}. Their respective merits are:

\begin{itemize}
  \item \mj{7m} drop:
  \begin{itemize}
    \item When \mj{9m} becomes our final wait, it has higher win rate
  \end{itemize}
  \item \mj{9m} drop:
  \begin{itemize}
    \item \mj{6m} draw gives us perfect 1-away and higher chance of pinfu
  \end{itemize}
\end{itemize}

Especially for \ref{mj:tobitt-1} where \mj{47p} draw gives us tanyao potential, in terms of expected value, \mj{9m} drop is the best.

\tehai{7799m 123456p 5578s}[][tobitt-1]

Drop \mj{68s} from \ref{mj:tobitt-2} for the possibility of perfect 1-away with \mj{6m} draw. Drop \mj{8s} first so that we still keep the possibility of \mj{5s} which gives us a ryanmen.

\tehai{237799m 123456p 68s}[][tobitt-2]

Shapes like \mj{224466p} consists of three skipping pairs can be treated in a similar fashion. However, since \mj{35p} draw give iipeiko chance, we will discuss it in the context of building iipeiko yaku in Section \ref{subsec:iipeiko}.

%\subsubsection*{2335の形}
%
%例え手牌の中に \mj{2335m} の形があったら、要すると、1ブロックにしたい時は打 \mj{5m}、1面子と一つの頭が欲しいなら打 \mj{2m}。\mj{5m} がなくなると受け入れが \mj{134m} になったが、\mj{134m} のどちらをツモったら \mj{33m} の部分が頭になれない。\mj{2m} を切ったら \mj{33m} が頭になって、受け入れが \mj{5m} のくっつきで広がるが、ツモ \mj{37m} でピンフが消え、ツモ \mj{7m} で嬉しくない愚形が残す。
%
%\tehai{123m 2335p 123789s -9m}[][nssg-1]
%
%\begin{table*}[!tb]
%  \centering
%	\begin{threeparttable}
%  \begin{tabular}{llll}
%    切る牌 & \mj{2p} & \mj{5p} & \mj{9m} \\
%    \hline
%    受け入れ & \mj{789m34567p} & \mj{789m1234p} & \mj{1234567p} \\
%    枚数 & 8種28枚 & 7種24枚 & 7種24枚 \\
%    ピンフがつくツモ & \mj{46p} (28.6\%) & \mj{9m3p}\tnote{1} (20.8\%) & \mj{3p}\tnote{1} \mj{46p} (41.7\%)
%  \end{tabular}
%  
%  \smallskip
%  \footnotesize
%  \begin{tablenotes}
%  \item[1] 待ちが\mjfn{124p}になるが\mjfn{2p}で和了ればピンフがない
%  \end{tablenotes}
%  
%  \caption{手牌\ref{mj:nssg-1}の受け入れ}
%  \label{tab:nssg-1}
%  \end{threeparttable}
%\end{table*}


%\subsubsection*{1224 Shape}
%
%1224の辺張の部分を払って辺張の受けの \mj{3m} を引いても損にならない\citep{gendai-penchan}。多くの場合で6889の形の中の9を余り牌として切ると大した差は出ないが、\mj{5m} と \mj{7m} を引いたら \mj{9m} が残る方が少し有利（\mj{8m}とあと一つのペアの双碰または \mj{56m} の両面）。ですので、\mj{9m} を残す基準は孤立の么九牌や役牌との比較\citep{yuusee-6889}。
%
%\begin{table}[ht]
%  \centering
%  \begin{tabular}{lll}
%    形 & \mj{6889m} & \mj{688m} \\
%    \hline
%    枚数 & 6枚 (\mj{7m}$\times4$ + \mj{8m}$\times2$) & 6枚 (\mj{7m}$\times4$ + \mj{8m}$\times2$)
%  \end{tabular}
%  \caption{1面子になる枚数の比べ}
%\end{table}
%
%\begin{table}[ht]
%  \centering
%  \begin{tabular}{lll}
%    形 & \mj{6889m} & \mj{688m} \\
%    \hline
%    変化 & \mj{567m} + \mj{88m} & \mj{567m} + \mj{88m} \\
%     & \mj{56m} + \mj{789m} & 
%  \end{tabular}
%  \caption{\mj{5m} と \mj{7m} を引いた場合（頭はもう一つできている上に）}
%\end{table}
%
%一向聴に絞るとき基本的に \mj{9m} 切りは効率的な方であるが、手牌\ref{mj:2-1}または\ref{mj:2-2}のような切る候補は孤立の么九牌や役牌の場合 \mj{9m} を残す方は少し有利である。
%
%\tehai{6889m1235p167s44z-8s}[東1局　西家　ドラ\mj{7s}][2-1]
%
%\tehai{6889m1235p67s445z-8s}[東1局　西家　ドラ\mj{7s}][2-2]
%
%手牌\ref{mj:2-3}と\ref{mj:2-4}のようなテンパイまで遠い場合は孤立の么九牌やオタ風を切る方がマシである。
%
%\tehai{6889m125p1668s44z-5z}[東1局　西家　ドラ\mj{7s}][2-3]
%
%\tehai{6889m125p668s445z-2z}[東1局　西家　ドラ\mj{7s}][2-4]
%
%手牌\ref{mj:2-5}のように切る候補は \mj{9m} と孤立の役牌の場合、役牌を残し方がいい。しかし、\ref{mj:2-6}のように横引きが期待でき立直を目指す手牌なら孤立の役牌を切る方がいい。
%
%\tehai{6889m125p2668s44z-5z}[東1局　西家　ドラ\mj{7s}][2-5]
%
%\tehai{26889m237p456s44z-5z}[東1局　西家　ドラ\mj{7s}][2-6]
%
%%--------------------------------------------
%\subsubsection{112246の形}
%
%\begin{table}[ht]
%  \centering
%  \begin{tabular}{ll}
%    切る牌 & 条件 \\
%    \hline
%    \mj{1m} & タンヤオに向かう \\
%    \mj{2m} & ピンフに向かう \\
%    \mj{6m} & タンヤオもピンフもならない
%  \end{tabular}
%  %\caption{キャップション}
%\end{table}
%
%\tehai{112246m34p34566-6s}[][1-1]
%
%\tehai{112246m789p3467-8s}[][1-2]
%
%\tehai{112246m999p3467-8s}[][1-3]
%
%\tehai{112246m34678s11-1z}[][1-4]
%
%手牌\ref{mj:1-1}はタンヤオが見えるので、\mj{1m} を切って仕掛けでもしてタンヤオを目指す。手牌\ref{mj:1-2}にはタンヤオはちょっと無理だけど、ピンフがつく可能性が高いなら \mj{2m} を切ってリャンカンを残す。手牌\ref{mj:1-3}のような立直のみを目指す場合は \mj{6m} を切って \mj{12m} の双碰を狙う。手牌\ref{mj:1-4}は役牌に絡んでいるのでポン材を残し \mj{46m} の嵌張を落とす\citep{yuusee-112246}。
%
%%--------------------------------------------
%\subsubsection{12235の形}
%
%\tehai{4567m 12235p 67788s}[東1局　西家　7巡目　ドラ \mj{2s}][12235-1]
%
%\textbf{12235の形はほぼ面子固定でいい}\citep[][p.~32]{uzaku-nankiru}。手牌\ref{mj:12235-1}は頭のない形。頭ができやすい四連形や \mj{67788s} に手をかけるより \mj{12235p} から \mj{2p} を外すの方がいい。


% TODO: below here

%============================================
\subsection{4-Tile Shapes}

%--------------------------------------------
\subsubsection{Yonrenkei} \label{subsub:yonrenkei}

Shapes like \mj{3456p} are called \textit{yonrenkei} in Japanese, which means ``the shape of 4 consecutive numbers.'' If we draw \mj{36p}, we get a pair and a mentsu; if we draw \mj{2457p}, we get a ryanmen and a mentsu. It is therefore a strong shape. Depending on the rest of our hand, we may have to sacrifice this shape and reduce it to a mentsu.

\ref{mj:yonren-1} has two yonrenkei shapes. If we drop either of \mj{13s}, we get the most acceptance, but it is also very likely that we end up kanchan or shanpon bad-wait rii-nomi. The \mj{122233s} consists of \mj{123s} + \mj{223s}, which gives us iipeikou possibility. Therefore we will not touch \mj{122233s} and we opt to drop one tile from the weakest yonrenkei. Drop \mj{1m} \citep[][p.~14]{uzaku-nankiru}.

\tehai{1234m 5678p 122233s}[East-1, West seat, turn 7, dora \mj{9s}][yonren-1]

\ref{mj:yonren-2} has tanyao ready, and with \mj{07m} kanchan we have tanyao dora 2, a very valuable hand. On top of that, we already have 6 blocks in our hand, the merit of producing head or ryanmen out of \mj{5678s} yonrenkei becomes very little. Considering also the possibility of 567 sanshoku, we drop \mj{8s} here \citep[][p.~16]{uzaku-nankiru}.

\tehai{2207m 3466p 225678s}[East-1, East seat, turn 6, dora \mj{7m}][yonren-2]


\subsubsection*{Reaching Tenpai with 3456 Shape}

\textbf{Generally, drop the outermost tile} \citep{yoteru-yonrenkei}. If we have \mj{3456p} and we must dropping one to reach tenpai, dropping \mj{3p} works in general cases. Merits of cutting the outermost tile are:

\begin{itemize}
\item Outer tiles generally have less deal-in rate
\item Not giving other players easy time calling
\item Strength of discards: a \mj{6p} signals that \mj{39p} are relatively safe, while a \mj{3p} gives way less information about our hand.
\end{itemize}

Exceptions are:

\begin{itemize}
\item Doras and sanshoku are involved
\item To keep the possibility of catching kamichi's akadora
\item To keep possibilities of improvements
\item Inner tiles are safer according to table situation
\item Early in a round
\end{itemize}


\subsubsection*{Cases involving doras and sanshoku}

\ref{mj:yonren-3} is an open hand with \mj{3456s}. Our general policy is to drop \mj{3s}. However, since dora is \mj{2s} and, if we keep \mj{345s} shape and drawing one \mj{2s}, we can drop \mj{5s} instead. This is called \text{sliding} and is something one must keep in mind when they have a shuntsu close to dora. Drop \mj{6s} from \ref{mj:yonren-3}.

\tehai{3456s 2256p -2'13m 77'7z}[Dora \mj{2s}][yonren-3]

If we drop \mj{2s} form \ref{mj:yonren-4}, we will not be able to call chi on kamicha's \mj{0s} discard. If we instead keep \mj{234456s}, we can chi \mj{0s} with \mj{46s} and discard \mj{2s}. If \mj{0s} is already discarded, follow the general policy and drop \mj{2s}.

\tehai{2344556s 3367p -55'5z}[][yonren-4]

\ref{mj:yonren-5} has a slight possibility of 123 or 234 sanshoku. In this case, drop \mj{5p}.

\tehai{123m 2345p 2399s -55'5z}[][yonren-5]


\subsubsection*{Iipeikou-Like Shapes Has Less Improvements}

Dropping either of \mj{36p} from \ref{mj:yonren-6} will give us the same waits \mj{47p}. If we drop \mj{6p} and keep \mj{34556p}, drawing \mj{2p} and dropping \mj{5p} gives us \mj{23456p} sanmenchan. Dropping \mj{3p} means that we will lose this upgrade. Therefore, drop \mj{6p} here.

\tehai{345566p 23488s -6'66z}[][yonren-6]

Note that this is for open hands only since iipeikou shapes gives us no additional value when our hand is open. If this is a closed hand, definitely go for iipeikou since we cannot change our tiles after we riichi.


\subsubsection*{Half Suji}

When another player already declared riichi and we must drop either \mj{3p} or \mj{6p} in order to take tenpai, and \mj{9p} is genbutsu, half suji \mj{6p} is safer than musuji \mj{3p} \citep[][pp.~140--141]{miinin-stat}. Likewise, when \mj{8p} is genbutsu, half suji \mj{5p} is safer than musuji \mj{2p}.

The main reason why half suji is safer than musuji is because, in terms of possibility of the incomplete shape we deal in, musuji \mj{3p} deals into \mj{12p} penchan, which \mj{6p} does not deal into. The number of possible combinations of kanchan waits and ryanmen waits are the same for both \mj{36p}, so \mj{3p} has more possible shapes to deal into than \mj{6p}, hence higher deal-in rate.


\subsubsection*{Early in a Round}

To make sure that, when we draw \mj{6p} back later in the game we can drop \mj{3p}, we can drop \mj{3p} when it is very early in a round.

\subsubsection*{Cases of 1234 and 2345}

To be able to slide when we draw middle tiles, in most cases, we still drop outermost tiles from \mj{1234p} or \mj{2345p}. Other general rules such as keeping sanshoku chance or dora slide still apply. \citet[pp.~171--172]{kawamura-naki} also gives similar merits and demerits and suggests to drop innermost tile instead in the following situations:

\begin{itemize}
\item Hand is cheap
\item We prioritize defense over winning, and 1 (or 2) are generally much safer than 4 (or 5)
\item No aka dora acceptance, or aka dora can still be picked up by other blocks
\item No tanyao chance
\item Shimocha still stays menzen
\item We are in huge lead, and we don't mind shimocha opening their hand
\end{itemize}

For \mj{112344m} however, we should drop \mj{4m} and use \mj{234m} mentsu to catch possible \mj{0m} draw.


%--------------------------------------------
\subsubsection{Nakabukure} \label{subsub:nakabukure}

The shape \mj{3445p} is called \textit{nakabukure} in Japanese due to the duplicate of the middle tile of the sequence making it look like it is bulging (and hence nicknamed ``the bulging shape.'') 

\textbf{Nakabukure works the best with the head of our hand confirmed.} \mj{3445p} nakabukure accepts \mj{23456p}, which, except for \mj{4p} draw, make a mentsu and a ryanmen. If we are 1-away and have the head ready, upon drawing any of \mj{2356p}, we are ready for ryanmen tenpai. Contrarily, if \mj{3445p} becomes our final wait, we are forced to wait on \mj{4p} tanki, a very \textit{very} weak wait. Therefore, in cases where we have nakabukure and a ryanmen-toitsu, reducing the latter to a pair often is the best choice.

Dropping \mj{3p} from \ref{mj:nkbkr-1} and \ref{mj:nkbkr-2} is the best choice \citep[][pp.~96--97]{uzaku-haikouritsu}. Table~\ref{tab:nkbkr-2} lists the acceptance of \ref{mj:nkbkr-2} and it shows that confirming our head is better than destroying nakabukure or aryanmen.

\tehai{5667m 344678p 4556s}[][nkbkr-1]

\tehai{5667m 344678p 4456s}[][nkbkr-2]

\begin{table}[ht]
  \centering
  \begin{tabular}{lll}
    Drop & Acceptance & Good-wait draw \\
    \hline
    \mj{3p} & \mj{45678m4p234567s} 38 tiles & \mj{4578m345s} 23 tiles \\
    \mj{6m} & \mj{12345p234567s} 37 tiles & \mj{245p35s} 18 tiles \\
    \mj{4s} & \mj{45678m12345p} 33 tiles & \mj{45678m2p} 20 tiles
  \end{tabular}
  \caption{Acceptance of \ref{mj:nkbkr-2}.}
  \label{tab:nkbkr-2}
\end{table}

4-tile shapes such as \textit{ryantan}\footnote{Shapes such as \mjfn{4555p} which can be treated as a \mjfn{45p} ryanmen and \mjfn{55p} pair, or a \mjfn{4p} tanki and \mjfn{555p} anko, hence the name -- ryan(men and )tan(ki). As a final wait, it waits on \mjfn{346p}.} and \textit{kantan}\footnote{Shapes such as \mjfn{3555p} which can be treated as a \mjfn{35p} kanchan and \mjfn{55p} pair, or a \mjfn{3p} tanki and \mjfn{555p} anko, hence the name -- kan(chan and )tan(ki). As a final wait, it waits on \mjfn{34p}.} work the best without a head, and thus are incompatible with nakabukure. Since the number of waits of ryantan and kantan are close to or greater than ryanmens, we destroy nakabukure for better final waits. Therefore, dropping \mj{3s} from \ref{mj:nkbkr-3} and \mj{6s} from \ref{mj:nkbkr-4} are the best choices here \citep[][pp.~98--99]{uzaku-haikouritsu}. Especially in the case of \ref{mj:nkbkr-4}, even if \mj{0m} is not double dora, knowing that dropping any of \mj{5m4p6s} gives us the same number of 34 tiles of acceptance, reducing nakabukure to a mentsu is still slightly better as suggested by Table~\ref{tab:nkbkr-4}.

\tehai{6777m 234788p 2334s}[][nkbkr-3]

\tehai{0777m 234455p 5667s}[Dora \mj{5m}][nkbkr-4]

\begin{table}[ht]
  \centering
  \begin{tabular}{ll}
    Drop & Good-wait (waiting on $\geq 6$ tiles) draw \\
    \hline
    \mj{5m} & \mj{245p4578s} 21 tiles \\
    \mj{4p} & \mj{46m5p4578s} 24 tiles \\
    \mj{6s} & \mj{456m23456p} 25 tiles
  \end{tabular}
  \caption{Acceptance of~\ref{mj:nkbkr-4}.}
  \label{tab:nkbkr-4}
\end{table}


\subsubsection*{Nakabukure and a Tile One Tile Away}

Shapes such as \mj{23346p} can be treated as a \mj{246p} ryankanchan and \mj{33p} pair, or \mj{2334p} nakabukure and a floating  \mj{6p}. Therefore, when there are no heads confirmed in our hand, we keep this shape and hope that we get a head with \mj{356p} draws.

Drop candidates of \ref{mj:nkbkr-1away-1} are \mj{6m3p6p}, but dropping \mj{6m} gives us the widest acceptance as is suggested by Table~\ref{tab:nkbkr-1away-1} \citep[][p.~100]{uzaku-haikouritsu}.

\tehai{2346m 23346p 45678s}[][nkbkr-1away-1]

\begin{table}[ht]
  \centering
  \begin{tabular}{ll}
    Drop & Acceptance \\
    \hline
    \mj{6m} & \mj{356p369s} 20 tiles \\
    \mj{3p} & \mj{6m6p369s} 17 tiles \\
    \mj{6p} & \mj{6m3p369s} 16 tiles
  \end{tabular}
  \caption{Acceptance of \ref{mj:nkbkr-1away-1}. All the candidates share the same acceptance of \mj{369s} 11 tiles, which unfortunately give us tanki wait.}
  \label{tab:nkbkr-1away-1}
\end{table}

In the case of \mj{233468p} in \ref{mj:nkbkr-1away-2}, as is the case of taking headless 1-away, dropping \mj{3p} gives us the widest 1-away. In terms of the number of good wait draws, as shown in Table~\ref{tab:nkbkr-1away-2}, there is only 1 tile of difference between \mj{8p} drop and \mj{3p} drop. Considering the possibility of improvements with \mj{124p} draw, dropping \mj{8p} seems better. As long as no yakus or doras are involved, when there is no head confirmed, we prioritize nakabukure and a tile one tile away \citep[][p.~101]{uzaku-haikouritsu}.

\tehai{234m 233468p 45678s}[][nkbkr-1away-2]

\begin{table}[ht]
  \centering
  \begin{tabular}{lll}
    Drop & Acceptance & Draws giving us $\geq 6$ tiles of tenpai \\
    \hline
    \mj{2p} & \mj{57p369s} 19 tiles & \mj{57p} 8 tiles \\
    \mj{3p} & \mj{678p3456789s} 33 tiles & \mj{678p} 10 tiles \\
    \mj{8p} & \mj{356p369s} 20 tiles & \mj{356p} 9 tiles
  \end{tabular}
  \caption{Acceptance of \ref{mj:nkbkr-1away-2}.}
  \label{tab:nkbkr-1away-2}
\end{table}


\subsubsection*{Cases where We Drop Yonrenkei or Nakabukure Before Tenpai}

We have learned that both yonrenkei and nakabukure are good at making good shapes. This, however, forces us to treat it as 2 blocks, and to complete 2 mentsus out of these 4-tile shapes, we must draw at least 2 turns. Therefore, when we have enough blocks and they are all sufficient for advancing our hand towards tenpai, we can drop yonrenkei or nakabukure without sacrificing too much acceptance. According to \citet{clearrain-nkbkr-yrk}, when there are no other floating tiles to dispose of, the conditions to drop yonrenkei or nakabukure are,

\begin{itemize}
  \item 2-away
  \begin{itemize}
    \item Other 4 blocks are all ryanmens and pairs.
  \end{itemize}
  \item 1-away
  \begin{itemize}
    \item There is at least 1 ryanmen block.
  \end{itemize}
\end{itemize}

Hand \ref{mj:ny-break-1} already has 5 blocks (treating yonrenkei as 1 here.) All blocks are good shapes, so we drop either \mj{3p} or \mj{6p} from yonrenkei shape.

\tehai{2256m 3456p 334788s}[][ny-break-1]

On the other hand, there is a \mj{68m} kanchan in \ref{mj:ny-break-2}. In such cases, we keep yonrenkei or nakabukure in hope for better shapes, while also keeping the bad block so as to not lose its direct acceptance. Therefore, we turn to ryanmen-toitsus and drop \mj{3s}. We keep \mj{788s} so that we have a better chance of completing this block with tanyao tiles.

\tehai{2268m 3456p 334788s}[][ny-break-2]

While it is tempting to drop bad blocks and rely on yonrenkei or nakabukure to give us good shapes, when we are at 1-away, we have to compromise and accept the chance of bad waits. Therefore, we do not drop \mj{799m} from \ref{mj:ny-break-3}; we drop \mj{5p} instead.

\tehai{11799m 4556p 34s 333z}[][ny-break-3]

There is also a \mj{34m} ryanmen in \ref{mj:ny-break-5}, so we drop \mj{2s} from \mj{2345s} \citep[][pp.~99--100]{ooi-gundou}. Dropping \mj{9p} loses the direct tenpai chance on \mj{8p} draw.

\tehai{34m 66779p 2345888s}[West seat, turn 7, dora \mj{8m}][ny-break-5]

However, when the remaining blocks are so bad that there is just no chance for good wait, we can opt to drop one of the bad blocks and fall back on yonrenkei or nakabukure. Therefore, we drop \mj{2s} from \ref{mj:ny-break-4} and hope for improvement with \mj{5s} draw. We do not touch \mj{799m} since it accepts \mj{89m} 6 tiles, which is more than a simple kanchan like \mj{24s} \citep{clearrain-nkbkr-yrk}.

\tehai{11799m 4556p 24s 333z}[][ny-break-4]


%--------------------------------------------
\subsubsection{Aryanmen}

Shapes like \mj{3345p} are called \textit{aryanmen} in Japanese. The \textit{a-} prefix in \textit{aryanmen} means ``next to,'' and considering that, as a final wait, the shape has a ryanmen wait but 2 tiles less than an ordinary ryanmen, we can see how this shape is treated as something inferior than a ryanmen.

\mj{1123p} is also an aryanmen. Since it has no ryanmen improvements, nor is it a strong wait (since we already have two copies of the outermost wait \mj{1p}), it is the weakest aryanmen of all shapes in this family. In most cases, it is simply treated as \mj{123s} sequence and a floating \mj{1s}.


\subsubsection*{Aryanmen is the Weakest 4-Tile Shape}

In terms of good shape improvement, aryanmens have the least number of draws, as shown in Table~\ref{tab:yonmaikei}. That is why it is often treated as a weak shape and, in literatures such as Uzaku's efficiency book, you will find advices such as ``when in doubt, drop aryanmen, and more often than not you'll be right.''

\begin{table}[ht]
  \centering
  \begin{tabular}{lll}
    Shape & Acceptance & Good shape draw \\
    \hline
    Yonrenkei \mj{3456p} & \mj{12345678p} 28 tiles & \mj{2457p} 14 tiles \\
    Nakabukure \mj{3445p} & \mj{23456p} 16 tiles & \mj{2356p} 14 tiles \\
    Aryanmen \mj{3455p} & \mj{234567p} 20 tiles & \mj{456p} 9 tiles
  \end{tabular}
  \caption{Acceptance of 4-tile shapes.}
  \label{tab:yonmaikei}
\end{table}

Hand \ref{mj:aryanmen-1} has 3 kinds of 4-tile shapes. We must drop one of them, and we opt to drop the weakest of them all. Drop \mj{4s} \citep[][p.~102]{uzaku-haikouritsu}.

\tehai{3456m 6778p 4456s 44z}[East round, South seat][aryanmen-1]

Dropping \mj{6p} from \ref{mj:sanmenchan-1} and having both sanmenchan and the acceptance of dora \mj{2p} confirmed seem good. However, as is shown in Table~\ref{tab:sanmenchan-1}, dropping \mj{5s} actually gives us the widest acceptance and the most number of good wait draws. % どこからの牌姿？

\tehai{4567m 345667p 3455s}[East-1, East seat, turn 4, dora \mj{2p}][sanmenchan-1]

\begin{table}[ht]
  \centering
  \begin{tabular}{lll}
    Drop & \mj{5s} & \mj{6p} \\
    \hline
    Acceptance & \mj{23456789m2356789p} 51 tiles & \mj{47m258p25s} 23 tiles \\
    Good-wait draws & \mj{345678m2358p} 26 tiles & \mj{47m258p25s} 23 tiles
  \end{tabular}
  \caption{Acceptance of~\ref{mj:sanmenchan-1}.}
  \label{tab:sanmenchan-1}
\end{table}


\subsubsection*{Let Go of the ``Ryanmen'' in Aryanmens}

Aryanmens does not work well with headless hands. In a 14-tile hand with 2 uncompleted taatsu and 1 aryanmen, if we choose to keep the pair and the ryanmen in the aryanmen shape and drop one tile from other taatsus, at that moment there is one surplus tile from the destroyed taatsu, and thus our acceptance is narrowed. In such cases, it is better to simply reduce the aryanmen to a mentsu. Indeed, headless hands come with the risk of tanki wait. With ankos or other shapes from which we can easily get a head, there is no need to fret about such risk.

To make all of our tiles useful in \ref{mj:aryanmen-2}, we drop \mj{6p} from aryanmen instead of \mj{8m} or \mj{8p}. Drawing any of \mj{368m} also gives us ryanmen tenpai \citep[][p.~105]{uzaku-haikouritsu}.

\tehai{34568m 6678p 34888s}[][aryanmen-2]

Since there is a \mj{22s} pair and a \mj{344568s} ryanmenkanchan in \ref{mj:aryanmen-3}, dropping \mj{6m} from the aryanmen shape leaves us no surplus tiles. Indeed, dropping \mj{8s} and treating \mj{2234s} also as an aryanmen leave us no surplus tile as well, and the difference in acceptance is just 1 tile. However, keeping \mj{6678m} aryanmen shape means that we must compromise on \mj{9m} draw which makes us lose tanyao \citep[][p.~104]{uzaku-haikouritsu}.

\tehai{6678m 34p 22344568s}[][aryanmen-3]

There is 2 kanchan blocks in \ref{mj:aryanmen-4}. As shown in Table~\ref{tab:aryanmen-4}, dropping any tiles from the kanchan blocks leaves us a surplus tile and narrowed acceptance. Taking headless 1-away indeed has the widest acceptance but also the risk of tanki wait, and, in the case of \ref{mj:aryanmen-4}, there is no shape that gives us a head or serves as a strong wait when our hand is headless. However, notice that for our \mj{23s} ``ryanmen'', we already have two copies of \mj{1s} in our hand, and, if there is another copy of \mj{1s} on the table, this \mj{1123s} shape is actually almost as weak as a kanchan waiting on \mj{4s}. Considering also the fact that \mj{1123s} has no improvement while kanchans at least have ryanmen improvements, dropping \mj{1s} is slightly better here \citep{hirasawa-aryanmen}.

\tehai{24678m 34579p 1123s}[][aryanmen-4]

\begin{table}[ht]
  \centering
  \begin{tabular}{ll}
    Drop & Acceptance \\
    \hline
    \mj{2m} & \mj{4m8p14s} 13 tiles \\
    \mj{1s} & \mj{234m789p} 20 tiles
    \end{tabular}
  \caption{Acceptance of \ref{mj:aryanmen-4}.}
  \label{tab:aryanmen-4}
\end{table}


%--------------------------------------------
\subsubsection{Confirming a Mentsu or Confirming a Head}

In cases where our hand is full of ryanmens and shuntsus and the only pair is buried in nakabukure or aryanmen, we must either drop the sequential part or drop the pair part. While it is true that dropping the pair and taking headless 1-away gives us the widest acceptance and thus higher chance of reaching tenpai, we also face the risk of tanki wait. Tanki waits (except for terminal and honor tile tanki waits) have terrible win rate; considering that our goal is to slam that riichi stick with best outcome possible, we would treat this as undesirable outcome\footnote{Practically speaking, even if we unfortunately arrive at tanki wait, we can still choose to dama our hand, regardless of it having yaku or not, and wait for our draws that give us better waits.} and focus on the rest of our acceptance. The theme of this section is thus to maximize the number of good-wait ($\geq 6$ tiles) tenpai depending on the actual composition of our hand. According to \citet[pp.~46--55]{uzaku-haikouritsu} and \citet{yuusee-1223fix},

\begin{itemize}
  \item Confirming head is better when
  \begin{itemize}
    \item It is still early and other players are still dealing with useless terminal tiles.
    \item We are in a table situation where we are required to force good-wait or yaku.
    \item Taatsus are involved in yakus.
  \end{itemize}
  \item Confirming mentsu is better when
  \begin{itemize}
    \item It is late and we would like to reach tenpai as soon as possible or to reach keiten\footnote{\textit{Keiten}, or \textit{keishiki tenpai}, refers to a state where we are tenpai without a yaku in late turns. At exhaustive draw, players that are not in tenpai must pay those who are in tenpai, so keiten is something we must keep in mind unless someone already called riichi.}.
    \item Drawing duplicate copies of tiles in taatsus is more likely than completing those taatsus.
    \item The shapes of the rest of our hand makes it salvageable even if we complete the taatsus first.
    \item Making it a mentsu enables dora slide.
    \item We have an anko and thus not required to confirm a head.
  \end{itemize}
\end{itemize}

\begin{table}[t]
  \centering
  \begin{tabular}{l|l|ll}
       & Confirming head & Confirming mentsu \\
    Shape & Ryanmen tenpai & Ryanmen tenpai & 6-tile tenpai \\
    \hline
    Ryanmen + ryanmen + shuntsu & 16 tiles & 12 tiles & 0 tile \\
    Ryanmen + sanmenchan        & 19 tiles & 18 tiles & 8 tiles \\
    Ryanmen + \mj{56789p}       & 15 tiles & 18 tiles & 8 tiles \\
    Ryanmen + \mj{66778p}       & 15 tiles & 14 tiles & 8 tiles \\
    Ryanmen + \mj{77889p}       & 15 tiles & 10 tiles & 0 tile \\
    Ryanmen (nidouke) + shuntsu & 12 tiles & 12 tiles & 4 tiles 
  \end{tabular}
  \caption{The difference between confirming head and confirming mentsu in terms of the number of draws that would give us tenpai with good quality. Bad waits such as shanpon and tanki are omitted here.}
  \label{tab:mentsu-or-head}
\end{table}

In case of the uncompleted taatsus being unrelated ryanmens, headless 1-away has the widest acceptance of 28 tiles, 16 tiles of which complete one of the ryanmen taatsus and force us to drop one tile from the other taatsu and take tanki wait. On the other hand, confirming mentsu leaves us 12 tiles of good-wait tenpai draws, as listed in Table~\ref{tab:mentsu-or-head}. Thus, it generally is better to drop \mj{13m} from \ref{mj:moh-0} and confirm our head.

\tehai{1223m 34678p 34578s}[][moh-0]

For shapes like \mj{66778p} in \ref{mj:moh-1}, even if we draw \mj{14s} first, we can still drop \mj{7p} and make the resulting aryanmen shape as our final wait, not too bad. If this shape sits on the edge like \mj{77889p} in \ref{mj:moh-2}, then we lost the leisure of aryanmen tenpai since \mj{7789p} does not accept ``10p.'' Therefore, we confirm its mentsu in \ref{mj:moh-1} while we confirm its head in \ref{mj:moh-2}.

\tehai{1223m 66778p 23678s}[][moh-1]

\tehai{1223m 77889p 23678s}[][moh-2]

In case of sanmenchan, it has more tiles of acceptance than a ryanmen and thus more likely to get completed. However, even then, we can still expect to get good waits when we draw tiles around the completed sanmenchan. For example, after \mj{2m} drop from \ref{mj:moh-3}, upon drawing \mj{8p}, we can first dama\footnote{\textit{Dama} refers to a state where we are in menzen tenpai but no riichi is called.} and expect \mj{235689p} draws to give us sanmenchan or irregular three-sided waits.

\tehai{1223m 34567p 23678s}[][moh-3]

While it is true that, even in the case of \ref{mj:moh-0}, suppose that we take headless 1-away anyway and end up tanki wait, we can still wait several turns before we get nobetan waits\footnote{Waits with yonrenkei shape \mj{3456p}.}. However, in such cases, we have less draws that give us nobetan waits than the case of \ref{mj:moh-3}, and we can barely draw into three-sided waits. Therefore, completed sanmenchan is regarded as salvageable, while unrelated shuntus are not.

In case of two ryanmens with overlapping acceptance, the number of ryanmen tenpai draw after having our head confirmed decreases from 16 tiles to 12 tiles, the same number when we confirm our mentsu. Note that, for shapes like \mj{3467p} in \ref{mj:moh-3a}, if our hand is headless, the nidouke shape actually additionally accepts \mj{5p} since we can get nobetan wait out of \mj{34567p} shape. Therefore, in terms of good-wait draws, having our mentsu confirmed is actually better. Drop \mj{2m} from \ref{mj:moh-3a}

\tehai{1223m 3467p 123678s}[][moh-3a]

As a rule of thumb, when the ryanmen taatsus are unrelated and do not compound with existing mentsus, having our head confirmd gives us the most probability of having good waits. However, for shapes like \mj{34567p}, \mj{34556p}, \mj{33445p} which can still give us aryanmen or nobetan wait even when our hand is headless, taking headless 1-away is optimal.


\subsubsection*{Cases where Taatsus are not All Ryanmens}

In cases where there are taatsus other than ryanmen, the number of draws that give us ryanmen tenpai decreases, but the number of draws that turn incomplete taatsus into head remains the same, so it becomes clear that confirming mentsu is better than confirming head.

For example, if we confirm our head in \ref{mj:moh-4}, the only draw that gives us ryanmen tenpai is \mj{4p} 4 tiles. On the other hand, if we confirm our mentsu, our ryanmen tenpai draws become \mj{35p} 6 tiles.

\tehai{1223m 35789p 23678s}[][moh-4]

In the case of \mj{67889p} in \ref{mj:moh-5}, if we drop \mj{9p}, we will have either aryanmen wait or ryanmen wait on any draws, as shown in Table~\ref{tab:moh-5}. On the other hand, if we confirm our mentsu with \mj{2m} drop, we will have wider acceptance, and more tiles to reach good-wait tenpai. Therefore, \mj{2m} drop is the best choice here. Note that the shape \mj{67889p} becomes a mentsu and a head with \mj{5689p} draws when our hand is headless.

\tehai{1223m 67889p 23678s}[][moh-5]

\begin{table}[ht]
  \centering
  \begin{tabular}{lll}
    Drop & \mj{2m} & \mj{9p} \\
    \hline
    Acceptance & \mj{56789p1234s} 29 tiles & \mj{2m58p14s} 16 tiles \\
    Good-wait draws & \mj{5689p14s} 20 tiles & \mj{2m58p14s} 16 tiles
  \end{tabular}
  \caption{Acceptance of \ref{mj:moh-5}.}
  \label{tab:moh-5}
\end{table}

In the case of \mj{66788p} shape in \ref{mj:moh-6}, confirming mentsu here is better since we can either drop \mj{6p} or \mj{8p} for aryanmen tenpai when \mj{23s} ryanmen gets filled; we can dama with \mj{7p} draw which gives us iipeiko yaku. Its acceptance is listed in Table~\ref{tab:moh-6}.

\tehai{1223m 66788p 23678s}[][moh-6]

\begin{table}[ht]
  \centering
  \begin{tabular}{lll}
    Drop & \mj{2m} & \mj{6p} \\
    \hline
    Acceptance & \mj{56789p1234s} 29 tiles & \mj{2m58p14s} 16 tiles \\
    Good-wait draw & \mj{5689p14s} 20 tiles & \mj{2m58p14s} 16 tiles \\
    Kanchan tenpai & \mj{23s} 6 tiles & 0 tiles \\
    Tanki tenpai & \mj{7p} 3 tiles & 0 tiles
  \end{tabular}
  \caption{Acceptance of \ref{mj:moh-6}.}
  \label{tab:moh-6}
\end{table}

However, if this iipeiko kanchan shape sits on the edge, like \mj{77899p} shape in \ref{mj:moh-7}, confirming \mj{123m} mentsu or \mj{7899p} aryanmen do not differ in terms of good-wait draws, as listed in Table~\ref{tab:moh-7}. In terms of hand value, however, for the same menzen-exclusive 1 han, iipeiko gives more fu than pinfu. The number of draws that give us yaku after \mj{2m} drop is also greater than the acceptance after dropping \mj{7p}. Therefore, dropping \mj{2m} is slightly better here.

\tehai{1223m 77899p 23678s}[][moh-7]

\begin{table}[ht]
  \centering
  \begin{tabular}{lll}
    Drop & \mj{2m} & \mj{7p} \\
    \hline
    Acceptance & \mj{6789p1234s} 25 tiles & \mj{2m69p14s} 16 tiles \\
    Good-wait tenpai & \mj{679p14s} 16 tiles & \mj{2m69p14s} 16 tiles \\
    Kanchan tenpai & \mj{23s} 6 tiles & 0 tile \\
    Tanki tenpai & \mj{8p} 3 tiles & 0 tile \\
    Pinfu draw & \mj{14s6p} 12 tiles & \mj{2m69p14s} 16 tiles \\
    Iipeiko draw & \mj{8p23s} 9 tiles & 0 tile
  \end{tabular}
  \caption{Acceptance of \ref{mj:moh-7}.}
  \label{tab:moh-7}
\end{table}

When the manzu part becomes \mj{1123m} aryanmen, even though dropping \mj{7p} gives us less acceptance, as shown in Table~\ref{tab:moh-8}, we actually end up 1-away that cannot draw into bad wait. Therefore, the best move here is not to confirm mentsu or to confirm head from \mj{1123m} aryanmen, but to drop \mj{7p} and create two aryanmens.

\tehai{1123m 77899p 23678s}[][moh-8]

\begin{table}[ht]
  \centering
  \begin{tabular}{lll}
    Drop & \mj{1m} & \mj{7p} \\
    \hline
    Acceptance & \mj{6789p1234s} 25 tiles & \mj{14m69p14s} 20 tiles \\
    Good-wait tenpai & \mj{679p14s} 16 tiles & \mj{14m69p14s} 20 tile \\
    Kanchan tenpai & \mj{23s} 6 tiles & 0 tile \\
    Tanki tenpai & \mj{8p} 3 tiles & 0 tile
  \end{tabular}
  \caption{Acceptance of \ref{mj:moh-8}.}
  \label{tab:moh-8}
\end{table}


%============================================
\subsection{6-Tile Shapes} \label{subsec:6-tile}

6-tile shapes are a family of shapes that consist of 6 tiles in a suit and, in the context of tile efficiency, compounded with a pair to work with two-pair principle. These particular shapes either are copies of tiles stacked together, like \mj{334456p33z}, or have one seemingly useless tile, such as \mj{334568p33z}, in a way that often puzzles people and makes them cut tiles that could cost them acceptance and tenpai chance. 

\begin{table}
  \centering
  \begin{tabular}{lll}
    Shape & Acceptance & Improvement \\
    \hline
    \mj{444556p33z} & \mj{34567p3z} 16 tiles \\
    \mj{334567p33z} & \mj{2358p3z} 15 tiles \\
    \mj{344456p33z} & \mj{2457p3z} 14 tiles \\
    \mj{444566p33z} & \mj{4567p3z} 12 tiles &  Draw \mj{3p}  drop \mj{6p} \\
    \mj{455566p33z} & \mj{4567p3z} 12 tiles &  Draw \mj{3p}  drop \mj{6p} \\
    \mj{344568p33z} & \mj{257p} 11 tiles & Draw \mj{4p} drop \mj{8p} \\
    \mj{244456p33z} & \mj{347p3z} 11 tiles &  Draw \mj{56p}  drop \mj{2p} \\
    \mj{334556p33z} & \mj{347p3z} 11 tiles &  Draw \mj{25p}  drop \mj{3p} \\
    \mj{334456p33z} & \mj{235p3z} 11 tiles &  Draw \mj{47p}  drop \mj{3p} \\
    \mj{444557p33z} & \mj{56p3z} 8 tiles &  Draw \mj{3p}  drop \mj{7p} \\
    \mj{334568p33z} & \mj{37p3z} 8 tiles &  Draw \mj{245p}  drop \mj{8p} \\
                    &                &  Draw \mj{4p}  drop \mj{3p} \\
    \mj{134568p33z} & \mj{27p} 8 tiles &  Draw \mj{4p}  drop \mj{1p} \\
                    &              &  Draw \mj{5p}  drop \mj{8p} \\
    \mj{344457p33z} & \mj{46p3z} 7 tiles &  Draw \mj{235p}  drop \mj{7p} \\
    \mj{133455p33z} & \mj{24p} 7 tiles &  Draw \mj{356p}  drop \mj{1p} \\
                    &              &  Draw \mj{6p}  drop \mj{3p} \\
    \mj{334566p33z} & \mj{36p3z} 6 tiles &  Draw \mj{2457p} drop \mj{3p} or \mj{6p} \\
  \end{tabular}
  \caption{6-tile shapes, with a head included \citep{6plus2}.}
  \label{tab:6plus2}
\end{table}

Among all the shapes listed in Table~\ref{tab:6plus2}, the shape \mj{344568p} is the most popular one called \textit{ryanmenkanchan} and is the one that is mentioned the most in books and materials on tile efficiency. There are also cases where we have 7-tile shapes such as \mj{3345568p} where we can either make it \mj{334556p} or \mj{344568p}, and the best way of handling this shape depends on many other factors.

 
\subsubsection*{Dropping Surplus Taatsus and Forming Ryanmenkanchan} 

One of the many cases where ryanmenkanchan emerges is when we have surplus taatsus and must drop one. When our taatsus and mentsus are positioned in a particular way, the remaining shape forms a ryanmenkanchan and we get to keep part of the acceptance of the destoryed taatsu, thus minimalizing our loss. Examples include:

\begin{itemize}
\item Dropping \mj{2s} from \mj{2356778s}
\item Dropping \mj{1s} from \mj{1356778s}
\item Dropping \mj{3s} from \mj{3356778s}
\end{itemize}

In \ref{mj:rmkc-1}, if we drop either of \mj{67m3578s}, at that moment we have a surplus tile in our hand, and we end up 4 kinds 16 tiles of acceptance, as shown in \ref{tab:rmkc-1}. The key here is to identify the \mj{356778s} ryanmenkanchan, which makes use of all souze tiles and gives us one more kind of tiles of acceptance. Drop \mj{2s} \citep[][p.~118]{uzaku-haikouritsu}.

\tehai{67m 11456p 2356778s}[][rmkc-1]

\begin{table}[ht]
  \centering
  \begin{tabular}{lll}
    Drop & \mj{2s} & \mj{578s} \\
    \hline
    Acceptance & \mj{58m469s} 19 tiles & \mj{58m14s} 16 tiles
  \end{tabular}
  \caption{Acceptance of \ref{mj:rmkc-1}.}
  \label{tab:rmkc-1}
\end{table}

\ref{mj:rmkc-31} has 6 blocks consisting of mentsus, ryanmens, and pairs. We intend to drop either of the pair. The pinzu shape may look completed, but if we drop \mj{8p}, there is actually \mj{344568p} ryanmenkanchan in our hand and it gives us more acceptance than dropping \mj{8s}.

\tehai{345m 3445688p 3488s}[][rmkc-31]

\ref{mj:rmkc-34} is a similar case where dropping \mj{2m} gives us the ryanmenkanchan \mj{245667m} which makes full use of every tiles \citep{yoteru-joui}.

\tehai{2245667m 22678p 45s}[East-1, turn 4][rmkc-34]


\subsubsection*{Floating Tiles and Ryanmenkanchan Potential}

When comparing floating tiles in our hand, we tend to overlook some floating tiles that position near other shapes in a way that has ryanmenkanchan potential. For example,

\begin{itemize}
\item \mj{36778p} in want of \mj{5p}
\item \mj{35677p} in want of \mj{8p}
\end{itemize}

Both \ref{mj:rmkc-5} and \ref{mj:rmkc-6} have floating \mj{8m} and \mj{8p}, and \mj{8m} is more isolated than \mj{8p}, so we drop \mj{8m} in hope that we can get ryanmenkanchan in pinzu suit \citep[][p.~122]{uzaku-haikouritsu}.

\tehai{348m 34458p 357788s}[][rmkc-5]

\tehai{348m 44568p 357788s}[][rmkc-6]


\subsubsection*{Choosing between Pinfu and Tanyao} 

\textbf{For the same 1 han, tanyao has a probability of scoring higher than pinfu, and we are ready to make calls and speed up, so we prioritize tanyao over pinfu} \citep[][p.~119]{uzaku-haikouritsu}. Dropping either of \mj{9p2s} from \ref{mj:rmkc-3} leaves us no surplus tile and gives us the same 19 tiles of acceptance. Dropping \mj{9p} gives us tanyao chance, and dropping \mj{2s} confirms pinfu. Considering the merits mentioned earlier, we opt to drop \mj{9p} and go for tanyao.

\tehai{678m 455679p 2238 -8s}[][rmkc-3]

Likewise for \ref{mj:rmkc-32}, we drop \mj{9p} and choose tanyao \citep[][p. 10]{uzaku-nankiru}.

\tehai{23488m 455679p 556s}[][rmkc-32]


\subsubsection*{Cases where Pinfu in Impossible}

Due to the characteristic of ryanmenkanchan that it always produces two shuntsu when completed, it works very well with pinfu. For cases where pinfu is not within sight such as having yakuhai pair in our hand, the importance of maintaining ryanmenkanchan decreases.

Hand \ref{mj:rmkc-yaku-3} has \mj{77z} pair and thus cannot have pinfu. Dropping \mj{3m} and \mj{8p} do not have any difference in terms of acceptance, but if we drop \mj{8p} here, we end up with just rii-nomi. Therefore, we drop \mj{3m} here. We trade \mj{4m} acceptance with \mj{8p7z} acceptance \citep{yoteru-rmkc}.

\tehai{356778m 788p 234s 77z}[][rmkc-yaku-3] 

It does not matter too much which to drop for rii-nomi hands such as \ref{mj:rmkc-yaku-4}. Adjust according to actual table situation \citep{yoteru-rmkc}.

\tehai{356778m 788p 33399s}[][rmkc-yaku-4]


\subsubsection*{Choosing between Guaranteed Pinfu and Iipeiko Chance}

For shapes like \mj{3556778p} or \mj{3567788p}, we can either make it ryanmenkanchan, or make it other kinds of 6-tile shapes. Either way, we still have 11 tiles of acceptance with the remaining shape. Therefore, the difference lies in the value the remaining shape could bring us.

Drawing \mj{6m} with \mj{55677m} or \mj{67788m} in \ref{mj:rmkc-yaku-3a} or \ref{mj:rmkc-yaku-3b} gives us iipeiko, so, in both cases, dropping \mj{3m} gives us the most possible value. It is true that, since the remaining shapes after \mj{3m} drop all have shanpon acceptance, we still could end up rii-nomi. However, The increase in value with one more han is greater than the decrease in value with one less han, so, in terms of expected value, \mj{3m} drop is still optimal \multicitep{uzaku-haikouritsu, pp.~120--121; yoteru-rmkc}. The cases where we would choose ryanmenkanchan and guaranteed pinfu instead are,

\begin{itemize}
\item \textit{Agaritoppu}, a kind of table situation where winning our hand confirms our first place.
\item In lead, and we are going to crush others' renchan\footnote{To win a hand as oya, and thus continue their oya status.} chance
\item Already valuable, such as dora 3.
\end{itemize}

\tehai{3556778m 45p 12388s}[][rmkc-yaku-3a]

\tehai{3567788m 78p 23499s}[][rmkc-yaku-3b]

\begin{table}[ht]
  \centering
  \begin{tabular}{lll}
    Drop & \mj{3m} & \mj{5m} \\
    \hline
    Acceptance & \mj{569m36p8s} 19 tiles & \mj{469m36p} 19 tiles \\
    Pinfu drawn & \mj{69m36p} (78.9\%) & \mj{469m36p} (100\%) \\
    Iipeikou draws & \mj{6m} (15.8\%) & None
  \end{tabular}
  \caption{Acceptance of \ref{mj:rmkc-yaku-3a}.}
  \label{tab:rmkc-yaku-3a}
\end{table}

For hands like \ref{mj:ipk-3} that has yakuhai anko, there is no pinfu, so the merit of keeping ryanmenkanchan diminishes. For merits such as \mj{0p} acceptance, we opt to drop \mj{3p} here \citep[][p.~35]{zero-teyaku}.

\tehai{23m 3556778p 33s 777z}[][ipk-3]


\subsubsection*{Dropping Ryanmenkanchan for a Better Chance of a Head}

The shape \mj{34556m} in \mj{134556m} ryanmenkanchan produces a head easily. To produce a head out of \mj{134556m} ryanmenkanchan, however, the only draw is \mj{1m}. In hands such as \ref{mj:rmkc-33} where we seek to confirm our head, we can drop \mj{1m} without losing the opportunity to incorporate \mj{2m}, while leveraging the flexibility of \mj{34556m} shape.

\tehai{134556m 67p 307889s}[][rmkc-33]

\begin{table}[ht]
  \centering
  \begin{tabular}{lll}
    Drop & \mj{1m} & \mj{9s} \\
    \hline
    Acceptance & \mj{1234567m5678p34568s} 49 tiles & \mj{247m58p46s} 27 tiles
  \end{tabular}
  \caption{Acceptance of \ref{mj:rmkc-33}.}
  \label{tab:rmkc-33}
\end{table}


\subsubsection*{Ryanmen-Toitsu and 6-Tile Shape}

When we are choosing between ryanmen-toitsu and 6-tile shape (with also a pair in our hand,) reducing ryanmen-toitsu to a ryanmen is, in most cases, the best choice in terms of acceptance and, when the number of accepted tiles are the same, the number of good-wait draws. Hand has \mj{334568p} which may seem week, but it indeed is a 6-tile shape according to Table~\ref{tab:6plus2}. Dropping \mj{3p} and \mj{5s} have the same acceptance of 16 tiles, but, as shown in Table~\ref{tab:6plus2-1}, keeping \mj{334568p} has the maximal 8 tiles of good-wait draws. By dropping \mj{5s}, we are actively denying the possibility of \mj{9m5s} draws that destroy our ryanmen.

\tehai{99m 334568p 123556s}[][6plus2-1]

\begin{table}[ht]
  \centering
  \begin{tabular}{llll}
    Drop & Acceptance & Good-wait draws \\
    \hline
    \mj{3p} & \mj{9m7p457s} 16 tiles & \mj{7p} 4 tiles \\
    \mj{8p} & \mj{9m3p457s} 14 tiles & \mj{9m3p} 4 tiles \\
    \mj{5s} & \mj{9m37p47s} 16 tiles & \mj{9m37p} 8 tiles
  \end{tabular}
  \caption{Acceptance of \ref{mj:6plus2-1}.}
  \label{tab:6plus2-1}
\end{table}

\ref{mj:6plus2-2} too follows the same principle and \mj{7m} drop gives us the maximal acceptance. Even though \mj{111223p} sits on the edge and loses 4 tiles than its middle counterpart, in terms of total acceptance, \mj{7m} drop still has more acceptance than, say, \mj{2p} drop \citep[][p.~134]{uzaku-haikouritsu}.

\tehai{677m 111223p 11567s}[][6plus2-2]

\begin{table}[ht]
  \centering
  \begin{tabular}{ll}
    Drop & Acceptance \\
    \hline
    \mj{7m} & \mj{58m1234p1s} 20 tiles \\
    \mj{2p} & \mj{578m14p1s} 17 tiles \\
    \mj{6m} & \mj{7m1234p2s} 14 tiles 
  \end{tabular}
  \caption{Acceptance of \ref{mj:6plus2-2}.}
  \label{tab:6plus2-2}
\end{table}


\subsubsection*{From 7 Tiles to 6 Tiles}

In the following cases, we will deal with 7-tile shapes that give us possibility of different 6-tile shapes. The point is:

\begin{itemize}
  \item Always pick the 6-tile shape that gives us the maximal acceptance.
  \item If the result 6-tile shapes have the same number of acceptance, pick the one with most yaku potential, be it tanyao, pinfu, or iipeiko.
\end{itemize}

We must drop one tile from pinzu part of~\ref{mj:6plus2-from7-1}. We can either drop any of \mj{57p} and have what is called ``sanmen-toitsu,''\footnote{A name I made up since, for \mjfn{455678p}, it is a sanmenchan \mjfn{45678p} and a toitsu \mjfn{55p}.} or drop any of \mj{48p} and have what is called ``\textit{niinii} shape.''\footnote{From Japanese where, for \mjfn{556778p}, \textit{ni} ``two'' refers to \mjfn{57p} which we have two copies and \textit{i} ``one'' refers to \mjfn{68p} which we have one copy.} Either way, we have \mj{69m1s} acceptance, so the difference is in the acceptance of sanmen-toitsu and niinii shape. As is shown in both Table~\ref{tab:6plus2} and Table~\ref{tab:6plus2-from7-1}, niinii shapes have 4 less tiles than sanmen-toitsu, so the correct drop is \mj{57p}. If \mj{0p} is still live, drop \mj{7p} to accept it \citep[][p.~132]{uzaku-haikouritsu}.

\tehai{78m 4556778p 11345s}[][6plus2-from7-1]

\begin{table}[ht]
  \centering
  \begin{tabular}{llll}
    Drop & Acceptance of pinzu shape \\
    \hline
    \mj{5p} & \mj{3679p} 13 tiles \\
    \mj{7p} & \mj{3569p} 13 tiles \\
    \mj{4p} & \mj{569p} 9 tiles \\
    \mj{8p} & \mj{367p} 9 tiles
  \end{tabular}
  \caption{Acceptance of \ref{mj:6plus2-from7-1}.}
  \label{tab:6plus2-from7-1}
\end{table}

When the shape sits on the edge like in \ref{mj:6plus2-from7-2}, the sanmen-toitsu shape loses 4 tiles of acceptance and there is no difference, in terms of acceptance, between \mj{68p} drop and \mj{9p} drop. However, by keeping \mj{566788p} niinii shape, we have iipeiko on \mj{7p} draw and thus the highest expected hand value. Therefore, drop \mj{9p} \citep[][pp.~132--133]{uzaku-haikouritsu}.

\tehai{78m 5667889p 11345s}[][6plus2-from7-2]

Shapes like \mj{234445s} accepts 1 less tile than sanmen-toitsu but still 3 more tiles than niinii shape. Therefore, we drop \mj{2s} from \ref{mj:6plus2-from7-3} for maximal acceptance. If \mj{234445s} becomes \mj{123334s} and loses 4 tiles of acceptance as in \ref{mj:6plus2-from7-4}, taking \mj{112334s} niinii shapes instead gives us the maximal acceptance and also iipeiko chance \citep[][p.~133]{uzaku-haikouritsu}.

\tehai{67m 34599p 2234445s}[][6plus2-from7-3]

\tehai{67m 34599p 1123334s}[][6plus2-from7-4]


%===========================================================
\subsection{Other Shapes}

%--------------------------------------------
\subsubsection{1345 Shape}

\textbf{Keep \mj{1p} in \mj{1345p} over isolated \mj{28s}} \citep{ittaku}. Drawing \mj{2m} gives us \mj{45m} ryanmen, and drawing \mj{36m} gives us a \mj{13m} kanchan which can be a decent wait. Compared to an isolated \mj{2s} which produces a penchan if we draw \mj{1s}, kanchan if we draw \mj{4s}, and a ryanmen if we draw \mj{3s}, \mj{1345m} has slight advantage. In this regard, we drop \mj{8p} from \ref{mj:1345-1}.

\tehai{1345m 2338p 5889s 11z}[East-1, West seat, turn 3, \mj{1p}][1345-1]

%--------------------------------------------
\subsubsection{13457 Shape}

\textbf{Keep isolated \mj{28s} over \mj{1p} in \mj{13457p}} \citep{ittaku}. Drawing \mj{7p} with \mj{13457p} makes \mj{1p} much weaker since, if we drop \mj{1p}, we will not lose sanmenchan chance with \mj{6p} draw, and we can still incorporate \mj{2p} since \mj{23457p} is a shape that has \mj{1348p} ryanmen-improvement draws. Therefore, drop \mj{1m} from \ref{mj:13457-1}.

\tehai{13457m 2668p 23699s}[East-1, East seat, turn 4, dora \mj{6p}][13457-1]

%--------------------------------------------
\subsubsection{Entotsu}

The shape \mj{22234p22s} is called \textit{entotsu} in Japanese whose literal meaning is ``chimney.'' Why this particular shape is associated with chimneys remains disputed. Since \mj{22234p} can be treated as either \mj{22p} + \mj{234p} and \mj{222p} + \mj{34p}, with \mj{2s} pair in our hand, we end up \mj{25p2s} wait.

\ref{mj:entotsu} is an open hand and we draw \mj{7m}. Since dora is \mj{8m}, we would like to work our shapes closer to dora. On top of that, since we have a kanchan \mj{68p} in our hand, it is more likely to draw another copy of \mj{68p} (6 tiles) than tsumo our winning tile \mj{7p}; in the former case, we can upgrade our hand to entotsu shape and catch dora \mj{8m}. Therefore, drop \mj{4m}\footnote{This is from one of my games and in the end I did pick entotsu shape and draw \mjfn{8m}.}.

\tehai{40556m 68p -7m -1'23p222z}[South-2, South seat, turn 4, dora \mj{8m}][entotsu]

\ref{mj:entotsu-2} has 6 blocks (seeing pinzu part as \mj{22p} + \mj{234p} ) and dropping any taatsu will not cause too much difference -- any of \mj{8m234p} drop gives 16 tiles of acceptance and any of \mj{34m78s} gives 15 tiles. According to Section \ref{subsub:6block-order}, the best course of action is to drop a pair. Although there is no immediate difference between \mj{8m} drop and \mj{2p} drop, but keep in mind that \mj{22234p} can become shapes listed in Table~{tab:6plus2} with one additional pinzu tile, so \mj{8m} drop is the best here \citep[][p.~130]{uzaku-haikouritsu}.

\tehai{3488m 22234p 12378s}[][entotsu-2]

\ref{mj:entotsu-3} is similar to \ref{mj:entotsu-2} but with its shapes more to the middle and one additional anko, tanyao and sananko weigh than pure efficiency and the improvements with future draws. Since \mj{78s} is now our bottleneck when it comes to scoring tanyao, we can drop it here \citep[][p.~131]{uzaku-haikouritsu}.

\tehai{3488m 44456p 22278s}[][entotsu-3]


%===========================================================
\subsection{6 Block Theory} \label{subsec:6block}

We have seen the term ``block'' here and there previously without defining its exact meaning. In this section, we will cover the term ``block'' thoroughly and explore the possibility of keeping multiple blocks in our hand, even cases where we have more blocks than needed.

A block is a group of tiles that we expect to become mentsu when we complete our hand. To complete a hand in mahjong using 14 tiles, we need 4 mentsus and 1 head\footnote{Except for kokushi musou and chiitoitsu.}. In that regard, we may think that we only need to keep 5 blocks in our hand. However, when we are still far away from completion and each block does not take up too much space, we can fit 6 blocks in our hand. The method of developing our hand while keeping all 6 blocks is called ``6 block theory'' \citep[][p. 16]{uzaku-haikouritsu}.

Merits of keeping 6 blocks include \multicitep{uzaku-haikouritsu, p.~17; yuusee-6block-1; yuusee-6block-2}:

\begin{itemize}
\item Keeping maximal choices between yakus and waits
\item Able to adapt to new table situation such as pon or kan from other players
\item Dropping tiles and fixing ryanmens make their outside wait come out more easily
\end{itemize}

Demerits of keeping 6 blocks include \multicitep{uzaku-haikouritsu, p.~18; yuusee-6block-1}:

\begin{itemize}
\item Keeping 6 blocks leaves us no room for follow-up tiles for each block, thus harder to reach perfect 1-away
\item The moment we drop 1 block to take 1-away, the destroyed block leaves us 1 surplus tile, thus making our acceptance narrower and lowers our defense
\end{itemize}

%---------------------------------------------
\subsubsection{Criteria of Keeping 6 Blocks}

Our policy regarding block number in our hand is \multicitep{uzaku-haikouritsu, p.~22--23; yoteru-6block;  yuusee-6block-1; yuusee-6block-2}:

\begin{itemize}
  \item Keep 6 blocks when
  \begin{itemize}
    \item Going straightforward leaves us high probability of cheap bad wait
    \item Blocks involve hand value and yaku
    \item There is no huge difference between blocks
  \end{itemize}
  \item Keep 5 blocks when
  \begin{itemize}
    \item There is huge difference between blocks
    \item We prioritize the speed of reaching tenpai
    \item Our hand is open and we want pairs for pon calls
    \item We prioritize defense
  \end{itemize}
\end{itemize}


\subsubsection*{Cases where We Keep 6 Blocks}

\textbf{Keep 6 blocks if going straightforward with blocks we currently have leaves us bad waits or cheap hand} \citep[][p. 23]{uzaku-haikouritsu}. For \ref{mj:block-3}, if there is no difference in player points, drop \mj{3p} and take perfect 1-away. However, if we are far behind other players, keeping \mj{3p} gives us a slight chance of sanshoku and dropping \mj{7m} confirms pinfu. This is the case where even a single tile can be treated as a whole block.

\tehai{234 677m 3 88p 234 78s}[][block-3]

When there is no discernible difference between taatsus, we can delay our decision by keeping 6 blocks \citep{yoteru-6block}. \ref{mj:block-make6a} has 6 taatsus, and dropping any of \mj{78m23p45s} leaves us 5 blocks. We can also reduce \mj{334m} to a ryanmen and progress with 6 blocks.

\tehai{334 78m 23 678p 45 88s}[East-1, South seat, turn 3, dora \mj{3z}][block-make6a]

We can drop \mj{4s} from \ref{mj:block-make6b} and keep 6 blocks, or drop \mj{6z} and reduce it to 5 blocks\footnote{If we drop \mjfn{7s} we lose pinfu and sanshoku becomes uncertain. Do not drop \mjfn{7s}.}. If we drop \mj{4s}, drawing \mj{7s} will not destroy sanshoku, and we keep the possibility of drawing another copy of \mj{6z}. If any other ryanmen gets filled up before we complete \mj{66z}, we can drop \mj{66z} and aim for pinfu. This is the case where keeping 6 blocks gives us more choices between yakus \citep{yuusee-6block-1}.

\tehai{345m 34 78p 445 77s 66z}[East-1, South seat, turn 3, dora \mj{8p}][block-make6b]

There are 3 possible options in \ref{mj:block-make6d},

\begin{itemize}
\item Drop \mj{3m} and wait for improvement of \mj{3440p}
\item Drop \mj{4p} and follow pure tile efficiency
\item Drop \mj{9p} and force tanyao
\end{itemize}

Considering the current turn and the value of our hand,

\begin{itemize}
\item It is already in the middle of a round $\to$ we cannot afford going back in shanten with \mj{9p} drop.
\item There are 3 pairs and we could possibly end up rii-nomi without tanyao $\to$ dropping \mj{4p} is not favorable.
\item It is more easily to make another group out of \mj{3440p} than out of \mj{233m}.
\end{itemize}

Therefore, dropping \mj{3m} is the best choice here \citep[][pp.~57--58]{fkt}. If we draw around \mj{3440p}, we can drop \mj{9p} and go for tanyao. If we draw \mj{9p4s} first, we can still drop \mj{4p} and riichi.

\tehai{233678m 344099p 44s}[Turn 8, dora \mj{7p}][block-make6d]


\subsubsection*{Cases where We Keep 5 Blocks}

\ref{mj:block-2} has a penchan that is clearly weaker than kanchan, so we drop \mj{89s} \citep[][p. 21]{ishibashi-black}.

\tehai{11 667m 68p 34 89s 333z}[][block-2]

\ref{mj:penchan-1} has \mj{6667p} shape which can be treated as 2 blocks \mj{66p} + \mj{67p}. If we reduce it to \mj{666p} 1 block, we have the widest acceptance of 7 kinds 24 tiles. However, if we don't draw any of \mj{789m} 9 tiles (1 copy of \mj{8m} is the dora indicator) first, we end up a weak penchan wait. We therefore drop \mj{8m}. If we later draw another copy of dora \mj{9m}, we can drop \mj{7p} \citep[][p.~16]{uzaku-nankiru}.

\tehai{23489m 6667p 34789s}[East-1, East seat, turn 7, dora \mj{9m}][penchan-1]

\textbf{Being able to pon is crucial for open hands} \citep{yoteru-6block}. We just called pon \mj{7z} in \ref{mj:block-make5a}. 6 block theory would require us to drop \mj{3m} here and abandon the ability to call \mj{3m7p} pon, which defies the purpose of going open in the first place. Open hands also come with the ability to change tiles and waits, so having tiles following taatsus becomes more important and contributes to our flexibility. If \ref{mj:block-make5a} develops into \ref{mj:block-make5b}, calling \mj{1245m} chi gives us two ryanmens, drawing \mj{5p} also improves our hand into ryanmens. Therefore, drop either of \mj{24p68s} depending on table situation.

\tehai{33478m 2477p 68s -77'7z}[][block-make5a]

\tehai{233478m 2477p -77'7z}[][block-make5b]

At the moment we draw \mj{5p} in \ref{mj:block-7}, we have 6 blocks. The only purpose of \mj{12m} is to accept \mj{3m}; if the other taatsus get completed first, we end up \mj{3m} penchan wait rii-nomi. It is therefore close to useless, and we would expect to draw another copy of \mj{16z} to give our hand more value; even for defense purposes, \mj{16z} is generally more safe than \mj{12m}. Thus we drop \mj{12m} here \citep[][pp.~154--155]{suphx}.

\tehai{12m 44578p 67s 1336z -5p}[East 1, East seat, turn 1, dora \mj{7p}][block-7]

\ref{mj:block-8} already has 5 blocks necessary for manzu honitsu. If we keep \mj{23p}, at worst our hand becomes haku nomi 1500 points. Therefore we drop \mj{23p}. Keep in mind that dora is \mj{7s} and keeping \mj{6s} makes us able to incorporate dora \citep[][pp.~156--157]{suphx}.

\tehai{3357m 23p 6s 114456 -5z}[East-1, East seat, turn 4, dora \mj{7s}][block-8]


%---------------------------------------------
\subsubsection{Order of Dropping a Block} \label{subsub:6block-order}

\subsubsection*{Comparing Different Shapes}

\textbf{The order of dropping a block is penchan \textgreater{} kanchan \textgreater{} pair \textgreater{} ryanmen} \citep[][pp. 21--22]{ishibashi-black}. Dropping either \mj{6m} (keeping 6 blocks) or \mj{1p} (reducing to 5 blocks) from \ref{mj:block-5} gives us the same acceptance of 24 tiles. However, dropping \mj{6m} leaves us a kanchan and we lose the chance of drawing \mj{16m} to avoid kanchan wait, so there is no reason to keep 6 blocks. Considering that \mj{68p} can be upgraded to ryanmen with \mj{5p} draw, we drop the weakest block \mj{12p}.

\tehai{11 668m 12 68p 35s 333z}[][block-5]

Even if the weakest taatsu has 1 dora, the order will not change \citep{yuusee-5block}. The \mj{79s} kanchan in \ref{mj:block-make5c} has a dora, but we have nice shapes in the rest of our hand. Thus, we drop \mj{7s} first, and drop \mj{33p} pair if later we draw another copy of \mj{9s}.

\tehai{234788m 3367p 3479s}[Dora \mj{9s}][block-make5c]

The reasons why we keep toitsus over kanchans \multicitep{uzaku-haikouritsu, pp.~30--33; yuusee-kanchan-toitsu} are:

\begin{itemize}
\item By keeping two pairs, we keep the possibility of their ryanmen improvements and therefore the possibility of perfect 1-away
\item If our pairs are outside, they come out more likely, especially for honor tile pairs
\item Two pairs means that our waits becomes \textit{atosuji}\footnote{When our discards after riichi makes our wait suji, it is called atosuji. The part \textit{ato} means ``behind, after'' in Japanese.} twice likely
\item Pairs do not use neighboring tiles, thus making them no-chance\footnote{If there are 4 copies of \mjfn{3p}, it makes \mjfn{1p} wait or \mjfn{2p} wait \textit{no-chance} since the ryanmens waiting on \mjfn{1p} or \mjfn{2p} requires 1 copy of \mjfn{3p}. It makes \mjfn{13p} kanchan also impossible. It, however, does not negate the possibility of \mjfn{11p} or \mjfn{22p} shanpon wait.} more likely
\end{itemize}

In terms of ryanmen improvement and complex taatsu improvement, an outside kanchan and one pair are weaker than two pairs of 2 (or 8). When we have two pairs, we keep ryanmen improvement and complex taatsu improvement of each pair. On the contrary, when we only have a pair and a kanchan, we must keep the pair as our head, so we lose its improvement; our only improvement becomes the improvement of kanchan. Table~\ref{tab:24p2288s} lists acceptance and all improvement of \mj{24p22s} (a kanchan and a pair) and \mj{2288s} (pairs of 2 and 8.)

\begin{table}[ht]
  \centering
  \begin{tabular}{lll}
    Shape & \mj{24p22s} & \mj{2288s} \\
    \hline
    Acceptance & \mj{3p} 4 tiles & \mj{28s} 4 tiles \\
    Ryanmen improvement & \mj{5p} 4 tiles & \mj{37s} 8 tiles \\
    Ryankanchan improvement & \mj{6p} 4 tiles & None \\
    Kanchan-toitsu improvement & None & \mj{46s} 8 tiles
  \end{tabular}
  \caption{Comparing \mj{24p22s} and \mj{2288s}.}
  \label{tab:24p2288s}
\end{table}

Even for pairs like \mj{1199s} which lack any ryanmen improvement, due to them being the outmost number tiles and twice likely becoming atosuji than a kanchan, they are decent as final waits. In this regard, the only cases where we would keep kanchan is when we have yakus such as tanyao or sanshoku involved, or when table situation makes pairs less favorable than kanchan.

\textbf{When we decide to reduce our hand to 5 blocks, keep safe tiles over the block to drop} \citep[][pp.~152--153]{suphx}. At the moment of drawing \mj{6z} in \ref{mj:block-6}, we have 6 blocks. We will not touch yakuhai pairs that give us value and speed; in that regard, we have no incentive to keep \mj{68s}. Instead of keeping \mj{68s} and maintaining 6 blocks until we complete another block, we must drop it immediately to make room for \mj{1z} safe tile. Drop \mj{6s}.

\tehai{456m 2245p 68s 1226 -6z}[South-2, East seat, turn 4, dora \mj{6p}][block-6]

When we are comparing ryanmens and kanchan, when they have the same number of acceptance left, \textbf{keep kanchan} \citep[][pp.148--149, p.~218]{fkt}. In Situation \ref{rittai:block-16}, a lot of manzu tiles has been cut, and both of \mj{34m} ryanmen and \mj{79m} kanchan in our hand accepts 4 tiles. Here, kanchan is better than ryanmen because,

\begin{itemize}
\item We can draw \mj{6m} and upgrade it into a ryanmen.
\item Even if all copies of \mj{6m} are cut, kan \mj{8m} becomes no-chance and more likely to come out.
\item Occasional increase in value due to the fu given by kanchan.
\end{itemize}

Even though 234 sanshoku seems plausible, with 3 copies of \mj{2m} already cut, and \mj{34p} ryanmen still uncertain which side of it is getting filled up, we opt to drop \mj{34m} here.

\begin{rittai}[h]
  \centering
  
  \drawrittai{
    jicha/seat = N,
    jicha/hand = 3479m 34p 234s 77z -55'5z,
    jicha/river = X 9m XXXX,
    jicha/riverb = XX,
    %
    toimen/river = XXXXXX,
    toimen/riverb = 2m X,
    %
    kami/river = 9m X 2m X 36m,
    kami/riverb = 66m,
    %
    shimo/river = X 2m XXXX,
    shimo/riverb = XX 5m,
    %
    kyoku = East round,
    dora = 3p
  }
  
  \caption{}
  \label{rittai:block-16}
\end{rittai}


\subsubsection*{Comparing Similar Shapes}

Even for shapes in the same class, since their acceptances are same (without table condition,) their differences lie in their improvement and shapes around them. In particular, we must also take into consideration how bad it backfires when we drop these taatsus. For each class of shapes, we have following merits \multicitep{suphx, pp.~152--153; yuusee-follow}:

\begin{enumerate}
\item (Kanchan, penchan) remaining number of acceptance
\item (Ryanmen, penchan) follow-up of acceptance
\item (Kanchan) follow-up of ryanmen-improvement, ryankanchan-improvement draws
\item (Toitsu) whether or not we can slide towards dora
\item (Toitsu) whether or not, when reduced to a floater, it has improvement
\end{enumerate}

There are two penchans \mj{89m} and \mj{12s} in \ref{mj:block-follow1}. To prepare for the moment we pon \mj{5z}, we must keep another pair \mj{33m}, so we cannot drop \mj{3m}. In this hand, after dropping \mj{12s}, its acceptance \mj{3s} can still be incorporated by its neighboring kanchan, which forms \mj{357s} ryankanchan. On the other hand, drawing \mj{7m} after dropping \mj{89m} is useless; it is just a floater with furiten risk. In other words, dropping \mj{12s} loses the least. Drop \mj{12s} \citep{yuusee-follow}.

\tehai{33489m 057p 1257s 55z}[East-1, North seat, turn 3, \mj{7z}][block-follow1]

\ref{mj:block-follow2} is also choosing between \mj{89m} and \mj{12s}. Since we have \mj{55m} pair, the acceptance \mj{7m} of penchan \mj{89m} gives us \mj{557m} kanchan-toitsu, and its kanchan-improvement draw \mj{6m} gives \mj{556m} ryanmen-toitsu. When there is no other information available to us, judging on how we can still incorporate accepted tiles and improvement draws, \mj{89m} becomes the weakest penchan. However, when \mj{3s} is once cut, or when we immediately draw \mj{7m} after dropping \mj{9m}, we still keep manzu and drop \mj{12s} \citep{yuusee-follow}.

\tehai{5589m 123678p 1288s}[East-1, East seat, turn 2, dora \mj{8s}][block-follow2]

\ref{mj:ist-1-1} is choosing between \mj{56m} and \mj{67p}. With \mj{123m} shape, when we drop \mj{56m} and draw \mj{4m}, we can still slide towards \mj{4m} and drop \mj{1m}. Drop \mj{56m} \citep{gendai-iishanten-yojouhai}.

\tehai{12356m2267p23456s}[][ist-1-1]

\ref{mj:block-9} is choosing between \mj{99p} and \mj{11s}. With \mj{11234s} shape, when we draw another dora \mj{4s}, we can drop \mj{1s} and maintaining a head in our hand. Therefore, drop \mj{9p} \citep[][p.~20]{uzaku-nankiru}.

\tehai{45m 67899p 1123467s}[East-1, West seat, turn 7, dora \mj{4s}][block-9]


%---------------------------------------------
\subsubsection{Dealing with Multiple Pairs}

We have discussed situations where we would destroy pairs or keep them in Section \ref{subsubsec:toitsu}. However, the arguments in Section \ref{subsubsec:toitsu} are made in the context of 5 blocks. In the following section, we will look at pairs in the context of 6 block theory.

\textbf{When there are three pairs, it is generally a good idea to destroy one} \citep{gendai-toitsu}, \textbf{since keeping pair has less acceptance than keeping ryanmen or kanchan} \citep[p.~34][]{uzaku-haikouritsu}. We valuate the strength of a pair according to the following merits \multicitep{gendai-toitsu; uzaku-haikouritsu, p.~169; uzaku-nankiru, p.~8},

\begin{itemize}
\item The number of improvement draws
\item Involvement in yakus
\item Strength as final waits
\end{itemize}

There are three pairs in \ref{mj:toitsu-1} and we must drop one. Considering that we can get \mj{11355m} shape with \mj{3m} draw, \mj{11p} becomes the most useless pair. Drop \mj{1p}.

\tehai{1155m11p678p34789s}[][toitsu-1]

In the case of \ref{mj:toitsu-2}, we keep \mj{22p} for ittsu chance with \mj{13p} draw, and drop \mj{11s} for slight chance of tanyao.

\tehai{77m 22456789p 1134s}[][toitsu-2]

Dropping either \mj{1s} or \mj{2s} from \ref{mj:toitsu-3} gives us the same number of 18 tiles of acceptance. In the case when \mj{78p} is completed first, \mj{1s7z} wait is more likely to win than \mj{2s7z}. Therefore, drop \mj{2s}.

\tehai{06778p 1122345s 77z}[][toitsu-3]

\textbf{Keep pairs that connect to other shapes} \citep[][pp.~78--79]{fkt}. We are choosing between \mj{7m8p5s} from \ref{mj:toitsu-5}. It seems that, since we already have 3 copies of \mj{7m}, \mj{77m} pair has the least acceptance. However, when we draw \mj{4m}, we can turn \mj{7m} into anko and incorporate \mj{4m} into our hand as \mj{46m} kanchan, which is impossible if we drop \mj{7m} first. \mj{55s} pair does not only has \mj{0s} acceptance but also \mj{1s} improvement. In this regard, drop \mj{8p} pair that has the least improvement.

\tehai{67778m 2388p 23455s}[Turn 7, dora \mj{2z}][toitsu-5]

Likewise, when choosing one from \mj{9m3p6p} to drop from \ref{mj:toitsu-6}, we drop \mj{9m} pair \citep[][pp.~80--81]{fkt}.

\tehai{0699m 12333668p 67s}[East round, non-dealer, dora \mj{1p}][toitsu-6]


%===============================================================
\section{Building Yaku} \label{sec:yaku}

Most of Section \ref{sec:tile-eff} is about choosing the choice that gives us the most tiles of acceptance with good quality. In cases where choices lead to the same number of acceptance, we then start discussing the expected value of our choices. In this section, we are going to focus more on the value of our hand, and, to some extent, to defy tile efficiency and choose actions that will not give us the most acceptance.

In a ruleset with yakus such as dora, akadora, ippatsu, and uradora, players can easily reach mangan without putting too much thought in building their hand; all they have to do is to be a little bit luckier with how the wall is built. In a ruleset where these opportunistic yakus do not exist, or when we are in dire need of points, making rii-nomi hands seems equal to digging our own grave, and we certainly do not want to sacrifice value potential just to reach rii-nomi faster.

Mahjong is a game of building mangan. For 30 fu hands, we start at 1000; for each additional 1 han we have, our hand value doubles, until it reaches 7700 for 4 han (or 8000 if it rounds up in some rulesets.) Likewise for other hands with different fu numbers, our hand value doubles as han number goes up until it reaches mangan. Therefore, when we have 1 or 2 han confirmed in our hand, spending effort on building hand value becomes worthwhile.

When we already have 3 han or more, yakus are still important -- it gives us the opportunity to dama or go open. Since the value of our hand is close to mangan cap, the benefit of riichi (potential ippatsu and uradora) becomes marginal, which then does not justify the decrease in win rate. Only riichi as last resort -- no yaku, and must push.

For a concrete example of how much we can sacrifice our efficiency for value, see \ref{mj:3900}. If we drop \mj{3p}, as is prescribed by tile efficiency, we end up 3900 points. Unless we are in huge lead in south game, we can opt to drop \mj{45p}, and keep ryanmen ryanmen for 7700 \citep[][pp.~98--99]{fkt}.

\tehai{40m 3345p 22334s -66'6z}[Turn 8, dora \mj{3p}][3900]


%---------------------------------------------
\subsection{Yakuhai} \label{sub:yakuhai}

Yakuhais refer to a family of honor tiles that gives us 1 han when we have 3 copies of them. Yakuhais include,

\begin{itemize}
  \item Sangenpai: any of haku \mj{5z}, hatsu \mj{6z}, or chun \mj{7z}.
  \item Bakaze: the same wind as the round. For east games, it would be \mj{1z}.
  \item Jikaze: the same wind as your own seat. For a player in South seat, it would be \mj{2z}.
\end{itemize}

Honor tiles that are not included in the previous categories are called \textit{otakaze}.

When we have a pair of yakuhais in our hand, we can call pon when one copy of it is dropped and we have a yaku confirmed instantly. When to call pon or not will be left for the section of call judgement. In this section, we are going to deal with lone yakuhai tiles and compare it with other floaters.

To put it simply, \textbf{it depends.} While we always teach beginners to drop honor tiles first, there are certainly situations where we would keep honor tiles and drop 19 or even 28 floaters first. Efficiency wise, honor tiles can only accept same kinds of tiles, so their acceptance always is less than a floating 1 or 9. However, the value and potential increase in speed it brings us when it pairs up makes it more versatile than a floating 1 or 9, whose job is to either provide a head or to create a penchan or kanchan.

The assessment of pure tile efficiency and neef of speed depends on table situation, hand value, hand shape, and of course the person making the assessment. Different people and AIs have their own rules and preferences. The list below shows different kinds of preferences from 6 sources.

\begin{itemize}
  \item \citet{senba-yakuhai}
  \begin{itemize}
    \item Hand is in good shape (one pair and two good shape blocks), 19 floater $>$ otakaze $>$ yakuhai
    \item Otherwise, yakuhai $>$ otakaze $>$ 19 floaters
  \end{itemize}
  \item \citet{tsuchida-yakuhai}
  \begin{itemize}
    \item Hand close to winning, or hand we want to win: 19 floaters $>$ otakaze $>$ yakuhai
    \item Hand that does not seem close to winning: otakaze $>$ yakuhai $>$ 28 floaters $>$ 19 floaters
  \end{itemize}
  \item NAGA's preference according to \citet{naga-yakuhai}
  \begin{itemize}
    \item Pinfu hand: otakaze $>$ yakuhai
    \item Otherwise: yakuhai $>$ 19 floaters $>$ otakaze
  \end{itemize}
  \item \citet{shibukawa-yakuhai}
  \begin{itemize}
    \item The 1 in 1345 $>$ yakuhai $>$ otakaze
    \item Otakaze will be your first drop most of the time
  \end{itemize}
  \item \citet{yoteru-yakuhai}
  \begin{itemize}
    \item One dora: 19 floaters $>$ yakuhai $>$ otakaze
    \item No dora: yakuhai $>$ 19 floaters $>$ otakaze
  \end{itemize}
  \item \citet{chigo-yakuhai}
  \begin{itemize}
    \item For headless hands, drop yakuhais first.
  \end{itemize}
\end{itemize}

To sum these up, it would be three kinds of orders, depending on our hand:

\begin{enumerate}
\item It is slow to remain closed, or it would very likely end up valueless rii-nomi: yakuhai $>$ otakaze $>$ 19 floaters
\item Pinfu or tanyao hands, or headless hands: 19 floaters $>$ otakaze $>$ yakuhai
\item Otherwise, hand is not too good nor too bad: 19 floaters $>$ yakuhai $>$ otakaze
\end{enumerate}


\subsubsection*{Comparing Yakuhais and Floaters}

A more detailed comparison between yakuhais and 19 and 28 floaters is as follows \multicitep{yuusee-yakuand1; yuusee-yakuand2; zero-teyaku, pp.~14--15},

\begin{itemize}
  \item Keep yakuhais over 19 floaters when
  \begin{itemize}
    \item Hand is too far away from tenpai when it stays closed (the number of ryanmens and mentsus are no more than 2.)
    \item The floater is suji to another floater, such as \mj{14p}.
    \item We don't have enough blocks and we have a pair.
    \item We have 5 blocks but one of those is bad shape.
    \item Pinfu does not seem doable.
    \item Floaters are not related to dora.
  \end{itemize}
  \item Keep yakuhais over 28 floaters when
  \begin{itemize}
    \item There are more than 2 bad shapes and no completed mentsu.
    \item Terminal tiles are dora.
  \end{itemize}
\end{itemize}

For hands such as \ref{mj:yakuhai-1} which only has two ryanmens, it would take long to complete the kanchans or to improve them, not to mention \mj{9m1s} still could give us penchans; by the time they are filled up, our opponents already tsumo'd their riichi. Keeping yakuhais here not only potentially speeds up when they pair up, but they also are generally safer than number tiles. Therefore, we opt to drop either of \mj{9m1s} here \citep{yuusee-yakuand1}.

\tehai{1249m 078p 156s 1557z}[East-2, South seat, turn 1, dora \mj{8s}][yakuhai-1]

Hand \ref{mj:yakuhai-2} has only 4 blocks. Efficiency wise, we could use \mj{789p} acceptance of \mj{9p} floater. However, drawing a taatsu around \mj{9p} means that \mj{0s} becomes surplus, and we end up dora 1 kanchan/penchan riichi of 2600 if we cannot improve the taatsu in time. Value wise, \mj{9p} is the most useless tile; we would like to see draws near \mj{0s}, or drawing another copy of \mj{56z} to give us good wait 2000 \citep{yuusee-yakuand1}.

\tehai{45677m 239p 0777s 56z}[East-1, East seat, turn 3, dora \mj{5m}][yakuhai-2]

Since there is only 1 ryanmen and 1 mentsu in \ref{mj:yakuhai-3}, it meets the criteria of keeping yakuhai over terminal floaters and the choices are \mj{1p} and \mj{1s}. When we drop \mj{1s}, its acceptance \mj{2s} cannot be picked up by its neighboring taatsu \mj{45s}. If we drop \mj{1p}, \mj{2p} forms a ryankanchan with \mj{46p} and \mj{3p} is the ryanmen improvement of that kanchan. Therefore, we drop \mj{1p} which backfires the least \citep{yuusee-yakuand1}.

\tehai{124789m 146p 145s 1z}[East-1, South seat, turn 4, dora \mj{7m}][yakuhai-3]

For once cut yakuhais, we still keep them over terminal floaters when our hand meets the criteria. Drop \mj{1s} from~\ref{mj:yakuhai-4}.

\tehai{34067m 1168p 16s 167z}[South-1, East seat, dora \mj{6s}, \mj{1z} is once cut][yakuhai-4]

For the same \ref{mj:yakuhai-5}, if dora is \mj{7s}, NAGA chooses to drop \mj{6z}. Otherwise, it chooses to drop \mj{9s}.

\tehai{447m 45p 1109s 44467z}[East-4, North seat][yakuhai-5]

For hands such as \ref{mj:yakuhai-6} which has no mentsu or ryanmen or whatsoever, if we just goes by pure efficiency, we end up rii-nomi, possibly bad wait also. For hopeless hands like this, we can even drop the penchan taatsu and go back in shanten. Either we draw another copy of \mj{167z}, or we get rid of terminal tiles and honor tiles altogether and go for tanyao \citep[][p.~15]{zero-teyaku}.

\tehai{4688m 22p 35589s 167z}[South seat, turn 1, dora \mj{8p}][yakuhai-6]

For hands like \ref{mj:yakuhai-7} where there is no completed mentsu or even ryanmens, we can even drop \mj{2m2s} \citep{yuusee-yakuand2}.

\tehai{268m 455899p 279s 56z}[East-4, West seat, turn 4, dora \mj{1p}][yakuhai-7]

When we have terminal tile doras in our hand, confirming yaku becomes first priority. Therefore, we drop \mj{8s} and rely on \mj{56z} in \ref{mj:yakuhai-8} \citep{yuusee-yakuand2}.

\tehai{2367m 11699p 348s 56z}[East-1, West seat, turn 3, dora \mj{1p}][yakuhai-8]

However, if our hand looks like \ref{mj:yakuhai-9} which has more ryanmens and completed blocks, going full efficiency is better than rely on a single yakuhai. Drop \mj{5z} \citep[][p.~15]{zero-teyaku}.

\tehai{18m 23667p 557999s 5z}[South seat, turn 1, dora \mj{9s}][yakuhai-9]


\subsubsection*{Yakuhai Attack}

Yakuhai attack refers to the strategy where we drop yakuhais first before they get to pair up in others hands, or when others hands are still in a state unready to call \citep[][pp.~16--17]{zero-teyaku}. Dropping yakuhai first means that our hand is blessed with good shapes, or yakuhais contribute little to our hand. In other words,

\begin{itemize}
  \item Our hand is full of shuntsus or taatsus and there are no pairs; calling pon on yakuhais does not speed up our hand.
  \item Our hand consists of only tanyao tiles if we drop yakuhai tiles.
  \item Our hand is ``in good shape,'' and to be in good shape, it has to
  \begin{itemize}
    \item Have more than 3 ryanmens or mentsus combined \citep[][p.~17]{zero-teyaku}
    \item Have 1 pair and two good shape blocks (including kanchan-toitsus) \citep{senba-yakuhai}
  \end{itemize}
\end{itemize}

If our hand meets any of the aforementioned conditions, we can drop yakuhais before otakaze. Note that otakaze can serve as our head for pinfu. For example, for hands like \ref{mj:yakuhai-attack-1} where tanyao is very likely, or for hands like \ref{mj:yakuhai-attack-2} which is in good shape according to \citet{senba-yakuhai}\footnote{\mjfn{1s} would be the first to go according to \citet{yuusee-yakuand1}: not enough blocks, having one pair. Some see hands like this doable, some see this far away from tenpai. This is why different folks have different preferences when it comes to yakuhai and terminal tiles.}, we can drop yakuhais first.

\tehai{33468m 22268p 08s 25z}[East round, West seat, turn 1][yakuhai-attack-1]

\tehai{1156m 23p 1779s 1347z}[East-1, North seat, turn 1, dora \mj{5m}][yakuhai-attack-2]

We already have two yakuhai pairs in \ref{mj:yakuhai-attack-3}, more than enough to end the game and take first place. It means that \mj{7z} does not contribute to our victory and, if we hold it for too long and drop it too late, we may end up letting others call it. Therefore we drop \mj{7z} here first despite an otakaze \mj{3z} in our hand \citep[][pp.~181--182]{kawamura-naki}.

\tehai{136m 78p 389s 223557z}[All-last, non-dealer, haipai, dora \mj{4z}, 500 points to overtake first place][yakuhai-attack-3]

Hand \ref{mj:yakuhai-attack-4} is the moment we pon dora \mj{5z}. The lone yakuhai \mj{7z} again contributes very little to our hand, since its value is already mangan. At this rate, the demerit of letting others call and possibly allowing them to win before we could is greater than the merit of value or defense. Drop \mj{7z}, even though \mj{9m} is close to useless also \citep[][pp.~183--184]{kawamura-naki}.

\tehai{1369m 247p 389s 7z -55'5z}[East-1, non-dealer, turn 1, dora \mj{5z}][yakuhai-attack-4]

% omitting shonpai jihai and 1ce cut jihai


\subsubsection*{Lone Dora Yakuhai}

\textbf{If there are other yakus for our hand, drop lone dora yakuhai before anything else. Otherwise, do not drop it easily} \citep{yoteru-lonedorayakuhai}. As mentioned earlier, if our hand already has strength points such as tanyao or pinfu, we do not particularly need dora yakuhai to pair up, and dropping them early on reduces the probability of them getting called pon. Therefore, for hands like \ref{mj:lonedorayaku-2} with tanyao and dora and pinfu potential, we drop \mj{5z} first.

\tehai{34m 468p 334056s 125z}[Dora \mj{5z}][lonedorayaku-2]

For \ref{mj:lonedorayaku-1}, however, it has no pinfu, no tanyao, and no dora whatsoever, so we keep \mj{5z} and other honor tiles. Drop any of \mj{8p125s}.

\tehai{24668m 8p 12578s 125z}[Dora \mj{5z}][lonedorayaku-1]

For \ref{mj:lonedorayaku-3} which is mentanpin\footnote{Short for menzen (riichi), tanyao, and pinfu.} 1-away, we would still drop dora \mj{5z} first, even if we draw this scary thing this late.

\tehai{23467m 34588p 56s 3z -5z}[Turn 7, dora \mj{5z}][lonedorayaku-3]

If our hand already has mangan value, it is better to drop lone dora yakuhai. \ref{mj:lonedorayaku-4} already has honitsu aka and chun, which is 8000, and getting \mj{5z} paired up only gives us 4000 extra points. If we cannot get it paired up in time and someone calls riichi, it would be very hard to drop it against riichi. Keep it if you are falling behind and are in dire need of points.

\tehai{11340899s 125z -7'77z}[East-1, South seat, turn 7, dora \mj{5z}][lonedorayaku-4]


%---------------------------------------------
\subsection{Tanyao}

\subsubsection*{Dropping Terminal Block from 6 Blocks}

\textbf{In most cases, keeping 5 tanyao blocks and drop the non-tanyao one is correct}\multicitep{gendai-tanyao; zero-teyaku, p.~22}, even if it is not the choice that gives us the most acceptance. If we drop either \mj{4s} or \mj{6s} from \ref{mj:tanyao-1} we will have maximal acceptance of 16 tiles. However, we will only get pinfu with \mj{5s} draw; any other tiles will only give us rii-nomi. If we drop \mj{9s}, we have tanyao confirmed, and we can open our hand. Drop \mj{9s}.

\tehai{23467m 345p 446699s}[][tanyao-1]

Likewise, drop \mj{9m} from \ref{mj:tanyao-2}. Efficiency wise, we would drop \mj{2p}. However, with only one ryanmen in our hand, we could end up waiting with kanchan wait. Therefore, \mj{9m} is the best course of action here, giving us tanyao and the opportunity to get rid of kanchans by calling or having any of those tiles paired up.

\tehai{35799m 2467p 24888s}[East-1, South seat, turn 3, dora \mj{8s}][tanyao-2]

Drop \mj{9m} from \ref{mj:tanyao-3}. Efficiency wise, \mj{3m} is correct, which loses acceptance of \mj{2m} dora and gives us only riichi pinfu. If we keep it and drop \mj{9m} instead, we have tanyao and possibly dora or sanshoku.

\tehai{356789m 346788p 45s}[East-1, South seat, turn 3, dora \mj{2m}][tanyao-3]

Efficiency wise, dropping \mj{7s} from \ref{mj:tanyao-4} gives us the maximal acceptance of 28 tiles and also the risk of ending up shanpon wait rii-nomi if we complete ryanmen first. We drop \mj{2p} here first, and when we draw any useful souzu tiles, we can drop \mj{1p} pair entirely \citep[][pp.~68--69]{fkt}.

\tehai{56m 11256p 2246778s}[Turn 4, dora \mj{6z}][tanyao-4]

We already have enough blocks in \ref{mj:tanyao-4a}. Going full efficiency means that we end up riichi dora 1. Note that, in this hand, we have strong yonrenkei and a strong floater \mj{7p}, and its still early, so we can afford going back in shanten and go for tanyao. Drop \mj{3z} \citep[pp.~70--71]{fkt}.

\tehai{3406m 7p 2346688s 33z}[East round, North seat, turn 4, dora \mj{1p}][tanyao-4a]

When any of the following conditions is met, do not force tanyao \citep[][p.~23]{zero-teyaku},

\begin{itemize}
\item We have enough good shapes.
\item Forcing tanyao increase the probability of ending up bad wait
\end{itemize}

If we drop \mj{11m} from \ref{mj:tanyao-5}, we most likely end up waiting on kan \mj{3m}. Ryanmen + ryanmen is sufficient for us to stay closed, so we drop \mj{2s} here and go full efficiency.

\tehai{1156m 34678p 24666s}[East-1, South seat, turn 3, dora \mj{8p}][tanyao-5]

Likewise, dropping kanchan from \ref{mj:tanyao-6} leaves us 1-away with good shapes, so we do not need to drop the only \mj{99m} pair and force tanyao and slow down our hand. We can either drop \mj{4p} or \mj{8p}; drop \mj{4p} if we are in the middle of a round for safety reasons, and drop \mj{8p} for possible iipeiko improvement with \mj{3p} draw.

\tehai{4599m 234468p 3478s}[East-1, South seat, turn 3, dora \mj{8s}][tanyao-6]

Likewise, for \ref{mj:tanyao-7}, we go full efficiency and drop the entire ryankanchan. Although this hand looks like \ref{mj:tanyao-2}, since \ref{mj:tanyao-2} is too far away from menzen tenpai, we had to force tanyao. For cases like \ref{mj:tanyao-7}, we are closer to menzen tenpai with 3 ryanmen blocks and not in need to force tanyao.

\tehai{35799m 2367p 23888s}[East-1, South seat, turn 3, dora \mj{8s}][tanyao-7]

Likewise, forcing tanyao in \ref{mj:tanyao-8} by dropping \mj{13s} leaves us pairs and kanchan and makes us lose the acceptance of dora \mj{5s}. Efficiency wise, keeping two pairs is better than a kanchan and a pair in terms of possible ryanmen improvement draws, so we drop \mj{68m} here \citep[][pp.~32--33]{fkt}.

\tehai{23468m 0577p 12334s}[Turn 5, dora \mj{5s}][tanyao-8]


\subsubsection*{Forcing Tanyao when Agaritoppu}

If we are agaritoppu, we must win our hand as soon as possible, and tanyao becomes our most reliable tool. For abysmal hands like \ref{mj:tanyao-force-1}, we are forced to call, and \mj{89m} penchan becomes completely useless that we can drop \mj{9m} first. After that, if we draw any of \mj{678m}, drop \mj{9p}. We can also call any of \mj{368p257s} and drop \mj{9p} (in the case of \mj{6p} chii, we cannot drop \mj{9p}; drop \mj{8p} instead.) \citep[][pp.~11--12]{kawamura-naki}.

\tehai{489m 247889p 22468s}[All-last, West seat, turn 1, dora \mj{3z}, 500 points to overtake first place][tanyao-force-1]

Hand \ref{mj:tanyao-force-2} is hopeless, and we still want to get that first place. As always, tanyao is our most reliable tool; yakuhai is possible too. In this regard, since we already have 4 tanyao blocks, the \mj{4z} pair is the most useless. Drop \mj{4z}. Same for \ref{mj:tanyao-force-3} which has two ryanmens \citep[][pp.~27--30]{kawamura-naki}.

\tehai{2488m 358p 46s 44567z}[All-last, South seat, turn 1, dora \mj{3z}, 500 points to overtake first place][tanyao-force-2]

\tehai{3488m 568p 46s 44567z}[All-last, South seat, turn 1, dora \mj{3z}, 500 points to overtake first place][tanyao-force-3]

For \ref{mj:tanyao-force-4} where 234 sanshoku seems viable, \mj{44z} becomes useful, \mj{8p567z} less so. Still, we keep \mj{8p} and drop \mj{567z} first since \mj{8p} still accepts \mj{678p} 11 tiles for tanyao. If we can't draw around \mj{8p}, pon \mj{4z} and force sanshoku instead \citep[][pp.~29--30]{kawamura-naki}.

\tehai{2488m 348p 24s 44567z}[All-last, South seat, turn 1, dora \mj{3z}, 500 points to overtake first place][tanyao-force-4]

Hand \ref{mj:tanyao-force-5} has two penchans and no completed blocks, making the 500 point difference far far away from it seems. Still, our dim hope lies in tanyao. In that regard, both \mj{1m9p} are useless. Drop \mj{1m} first, since, even if we draw \mj{3m} and \mj{23m} ryanmen is a furiten taatsu, we can still chi \mj{4m} to get rid of this taatsu. Dropping \mj{9p} makes \mj{78p} ryanmen furiten and it would be very hard to get rid of the furiten taatsu while also incorporating \mj{7p}. Do not immediately drop \mj{2m} after \mj{1m}; \mj{2m} draw which gives us pairs backfires a lot \citep[][pp.~53--54]{kawamura-naki}.

\tehai{12578m 24688s 2489p}[All-last, turn 1, non-dealer, dora \mj{3z}, 500 points to overtake first place][tanyao-force-5]


\subsubsection*{Forcing Tanyao when Menzen is Also Close}

Hand \ref{mj:tanyao-menzen-close-1} looks far better than previous hands and menzen tenpai is within reach. Even so, opening our hand with tanyao is still faster. In particular, in \ref{mj:tanyao-menzen-close-1}, there are \mj{2s8p} pairs and \textbf{2(8) comes out fairly easily.} Thus, drop \mj{1m} here and force tanyao \citep[][pp.~41--42]{kawamura-naki}.

\tehai{123457m 56688p 224s}[All-last, turn 1, non-dealer, dora \mj{3z}, 500 points to overtake first place][tanyao-menzen-close-1]

If our hand is better such as \ref{mj:tanyao-menzen-close-2} which is perfect 1-away if we drop \mj{7m}, dropping \mj{7m} is fine for very early turns.

\tehai{123457m 22678p 667s}[All-last, turn 1, non-dealer, dora \mj{3z}, 500 points to overtake first place][tanyao-menzen-close-2]

Hand \ref{mj:tanyao-menzen-close-3} looks like we can force tanyao by dropping \mj{1p} pair, but doing so will lose \mj{178p} 8 tiles of acceptance. Notice that, \mj{5s} drop not only give us the widest 1-away which accepts \mj{36m1578p} 20 tiles, 456 sanshoku is also possible with \mj{6m5p} chi and then \mj{1p} drop. Therefore, drop \mj{5s} here \citep[][pp.~43--44]{kawamura-naki}.

\tehai{45m 11234688p 4556s}[All-last, turn 1, non-dealer, dora \mj{3z}, 500 points to overtake first place][tanyao-menzen-close-3]

Likewise, for \ref{mj:tanyao-menzen-close-4}, \mj{1p} drop renders \mj{6m5s} useless, and it actually increases our shanten, unlike \mj{1p} drop from \ref{mj:tanyao-menzen-close-3}. Widest 1-away we can get from \ref{mj:tanyao-menzen-close-4} is 16 tiles. In that case, it is better to drop \mj{6m} and maintain 1-away with 12 tiles of acceptance, and after we draw or chi any of \mj{3467s}, we can opt to drop \mj{1p} entirely. This is similar to \ref{mj:tanyao-4} \citep[][pp.~43--44]{kawamura-naki}.

\tehai{34688m 11234p 4556s}[All-last, turn 1, non-dealer, dora \mj{3z}, 500 points to overtake first place][tanyao-menzen-close-4]

Hand \ref{mj:tanyao-menzen-close-5} has 4 pairs, and chiitoi seems closer than normal hand, but in agaritoppu, tanyao should still be the first priority. Still, to make chiitoi our backup plan, \mj{9m} drop then becomes the best course of action. If we want to keep the possibilities of a normal hand and a chiitoi hand and drop \mj{2m2s}, our hand becomes way slower \citep[][pp.~57--58]{kawamura-naki}.

\tehai{247889m 2277p 2488s}[All-last, turn 1, non-dealer, dora \mj{3z}, 500 points to overtake first place][tanyao-menzen-close-5]

Hand \ref{mj:tanyao-menzen-close-6} is mostly tanyao tiles and forcing tanyao by dropping \mj{9p} seems fine. However, the pairs of this hand are all middle tiles, so they would come out harder than usual. Thus, for \ref{mj:tanyao-menzen-close-6}, drop \mj{5s} and keep menzen possibility while also not overlooking 345 sanshoku chance \citep[][pp.~59--60]{kawamura-naki}.

\tehai{355m 357789p 22455s}[All-last, turn 1, non-dealer, dora \mj{3z}, 500 points to overtake first place][tanyao-menzen-close-6]

Hand \ref{mj:tanyao-menzen-close-7} too has middle tile pairs, and even if we drop \mj{9p}, the acceptance of \mj{778p} overlaps with that of \mj{45p}. It is better to simple take the 1-away with \mj{7p} drop \citep[][pp.~59--60]{kawamura-naki}.

\tehai{34577m 457789p 468s}[All-last, turn 1, non-dealer, dora \mj{3z}, 500 points to overtake first place][tanyao-menzen-close-7]

If we change the hand a little bit and make it \ref{mj:tanyao-menzen-close-8} which now has pairs of 2(8) tiles, we can force tanyao with it. Drop \mj{9p} \citep[][pp.~59--60]{kawamura-naki}.

\tehai{34588m 347889p 468s}[All-last, turn 1, non-dealer, dora \mj{3z}, 500 points to overtake first place][tanyao-menzen-close-8]


\subsubsection*{Keeping Tanyao as Plan B}

Before turn 7 or until someone opens their hand, we can keep tanyao possibility in our hand instead of safe tile. This particularly works well for cases where we only have one tanyao bottleneck such as \mj{23m} ryanmen. Scoring tanyao not only increases our hand value, in cases where we are agaritoppu or we have dora 3 it would also increase our win rate by allowing us to open our hand \citep[][pp.~155--158]{kawamura-naki}.

When we look at hands like \ref{mj:tanyao-9}, we should not only think that \mj{2367m} draws give us perfect 1-away, but also keep in mind that any of \mj{345678p345678s} gives us tanyao chance. This tanyao chance is particularly helpful when we are agaritoppu or we have dora 3. Suppose that we draw \mj{3p} and our hand becomes \ref{mj:tanyao-9a}, we can chi \mj{4p} and drop \mj{3m}. Even with \mj{3s} draw as in \ref{mj:tanyao-9b}, we can chi \mj{4s} and drop \mj{3m}.

\tehai{2367m 23477p 067s 3z}[][tanyao-9]

\tehai{2367m 233477p 067s}[][tanyao-9a]

\tehai{2367m 23477p 3067s}[][tanyao-9b]

If our hand has no dora, the increase in value (1000 or 1300) brought by tanyao does not justify the loss in defense. Thus, for \ref{mj:tanyao-9c}, we keep \mj{3z} and tsumogiri\footnote{\textit{Tsumogiri} refers to the action of drawing a tile and immediately discarding it. The word is a compound of \textit{tsumo} ``draw'' and \textit{giri}, derived from \textit{kiru} ``to cut.''} anything other than \mj{2367m7p}.

\tehai{2367m 23477p 567s 3z}[Dora \mj{4z}][tanyao-9c]

Likewise, hand \ref{mj:tanyao-10a} has a bottleneck \mj{78s}, so we keep \mj{4s} here in hope for tanyao and sanshoku. For hand \ref{mj:tanyao-10b}, however, since its tanyao is already confirmed, and there is already dora 2, there is no need for sanshoku, and \mj{4s} becomes useless. Drop \mj{3z} in \ref{mj:tanyao-10a} and \mj{4s} in \ref{mj:tanyao-10b}.

\tehai{23467m 23455p 478s 3z}[Dora \mj{5p}][tanyao-10a]

\tehai{23467m 23455p 467s 3z}[Dora \mj{5p}][tanyao-10b]

Previously we learned that, in the context of 6 blocks, drop \mj{9m} from the skipping pairs \mj{7799m} is usually better especially when there is tanyao chance. Even when we compare skipping pairs to floaters, we still must always keep tanyao in mind. For hands from \ref{mj:tanyao-11a} to \ref{mj:tanyao-11c}, we are comparing \mj{7799m} to a floater \mj{7s}. Even though our hand already has 6 blocks and \mj{7s} does not contribute to acceptance, with this particular shape \mj{123457s}, we can opt to chi \mj{6s} and drop \mj{1s}, so it contributes to hand value and speed.

\textbf{If the hand has 3900 or 3900 potential after we open it, drop} \mj{9m} \citep[][pp.~165--170]{kawamura-naki}. Hand \ref{mj:tanyao-11a} has dora 2 and is already 3900 if we open our hand, and drawing any of the aka doras gives us 7700. Hand \ref{mj:tanyao-11b} only has dora 1, but its \mj{45m} ryanmen accepts dora \mj{6m} and \mj{34p} accepts aka dora \mj{0p}, so it still has 3900 potential. Hand \ref{mj:tanyao-11c}, however, only has one dora and dora \mj{3z} is useless, so this hand has no 3900 potential and is thus not fit for calls, and we drop \mj{7s} instead.

\tehai{457799m 34p 123457s}[Dora \mj{7m}][tanyao-11a]

\tehai{457799m 34p 123407s}[Dora \mj{6m}][tanyao-11b]

\tehai{457799m 34p 123407s}[Dora \mj{3z}][tanyao-11c]

Hand \ref{mj:tanyao-chance-1} looks like pinfu hand. Except for \mj{34568m4567p} which give us tenpai or perfect 1-away or headless 1-away with anko, keep in mind that \mj{23p} draws allow us to pon. In flat situation, \mj{23p} is useless, but it is crucial in all-last agaritoppu situations (or dora 3) \citep[][pp.~17--18]{kawamura-naki}.

\tehai{4588m 12356p 234s 4z}[All-last, non-dealer, haipai, dora \mj{3z}, 500 points to overtake first place][tanyao-chance-1]


\subsubsection*{Deciding between Tanyao and Dabuton}

\textbf{If menzen tenpai is close, drop dabuton }\citep[][pp.~138--139]{fkt}. If we drop \mj{11z} from \ref{mj:tanyao-12}, we still have enough good shapes for tenpai. Dabuton and dora 1 is 5800, but if we stay closed, with riichi tanyao dora 1 we score 7700, and with ippatsu, tsumo, or uradora we can reach 11600. Therefore, value wise, dropping \mj{11z} and go for tanyao is better.

\tehai{22456m 22240p 67s 11z}[East round, East seat, turn 6, dora \mj{3z}][tanyao-12]

%%%% dabuton and ryanmen, not dabuton and tanyao
%
%\textbf{仕掛けが入ったらダブ東で12000を目指す}\citep[][pp.~140--141]{fkt}。手牌\ref{mj:tanyao-10}は立直や裏ドラなどがないため、両面両面で行くと2900や都合よくて赤ドラ引いて5800。一方、\mj{1z} を鳴いたら打点が2900から11600になるので段違いです。
%
%\tehai{22456m 40p 67s 11z -6'66z}[東場　東家　6巡目　ドラ \mj{3z}][tanyao-10]

If we already opened our hand like \ref{mj:tanyao-13}, without opportunistic yakus such as ippatsu, the value of dabuton dora 1 becomes greater than tanyao dora 1. Even if we cannot reach tenpai in time and someone calls riichi, \mj{1z} is still generally safer than other tanyao tiles. We drop \mj{3p} first. After that, compromise on 2900 if the ryanmens get filled first, or drop \mj{67s} when we draw \mj{6p} \citep[][pp.~143--144, p.~218]{fkt}.

\tehai{22456m 30p 67s 11z -2'22p}[East round, East seat, turn 6, dora \mj{3z}][tanyao-13]


\subsubsection*{Deciding between Tanyao and Yakuhai}

\textbf{Anything other than ryanmen ryanmen or dora 3 means we keep yakuhai pair} \citep[][pp.~177--180]{kawamura-naki}. For hands from \ref{mj:tanyao-14a} to \ref{mj:tanyao-14c}, we just called \mj{7s} chi on turn 7 to get rid of a kanchan, and we can either drop \mj{7z} for tanyao or drop any of the other taatsus for chun. Some may be thinking that if \mj{7z} does not come out this late, \textit{mochimochi}\footnote{Derived from Japanese \textit{motsu} ``to possess,'' it refers to a state where two opponents have two copies of one kind of yakuhai and thus they can never call their yakuhais.} is very likely the case, and we cannot win our hand without dropping \mj{7z}. Paradoxically, dropping \mj{7z} means that it will get picked by our opponent and make it harder to win our hands.

Therefore, there is little merit of dropping \mj{7z}. If it really is mochimochi, we are giving others advantage; if it is not, we lose defense when facing riichi. Only drop \mj{7z} pair if our taatsus are all ryanmens or we have dora 3 and we must get tanyao confirmed. Drop \mj{6m} from \ref{mj:tanyao-14a}, \mj{7z} from \ref{mj:tanyao-14b}, and \mj{7z} from \ref{mj:tanyao-14c}.

\tehai{23468m 56p 22s 77z -7'68s}[Turn 7, dora \mj{3m}][tanyao-14a]

\tehai{23467m 56p 22s 77z -7'68s}[Turn 7, dora \mj{3m}][tanyao-14b]

\tehai{23468m 06p 22s 77z -7'68s}[Turn 7, dora \mj{2s}][tanyao-14c]


%---------------------------------------------
\subsection{Pinfu}

\subsubsection*{Forcing Pinfu when Our Hand Has No Dora}

\textbf{Force pinfu when our hand has no dora or other yakus} \citep[][pp.~200--201]{doi-prob}. Hands like \ref{mj:pinfu-1} very likely end up rii-nomi with little value, the cheapest 1300. Efficiency wise, the \mj{3z} that we just drew is useless. If we drop \mj{3z} here, considering that unseen doras are more likely end up in others hands as the game progresses, most likely we will face riichis from others that have values higher than average\footnote{If all \mjfn{8p} are seen, it is another story.}. From the perspective of defense and hand value, we drop \mj{3m} here.

\tehai{334789m 56p 11788s -3z}[East-1, West seat, turn 5, dora \mj{8p}][pinfu-1]

\ref{mj:pinfu-2} is a similar hand, but we have 1 han from tanyao. With 1 han in our hand, we can become more aggressive and keep maximal acceptance. Drop \mj{3z}.

\tehai{334678m 56p 22788s -3z}[East-1, West seat, turn 5, dora \mj{8p}][pinfu-2]

We must drop any of the ryanmens of \ref{mj:pinfu-3}. In a ruleset with akadoras, 2--5 or 5--8 waits are stronger than 3--6 or 4--7 waits, so we drop \mj{45s} here. If our only pair \mj{9m} is not dora, drop \mj{4s} first, so that when we luckily draw \mj{0s}, we can drop \mj{9m} instead. If \mj{9m} is dora, drop \mj{5s} first, so that when we luckily draw \mj{0s}, we can choose other ryanmens to drop instead \citep[][pp.~13--14]{kawamura-naki}.

\tehai{3499m 12367p 45789s}[][pinfu-3]


\subsubsection*{Deciding between Pinfu and Yakuhai}

If we have \ref{mj:pinfu-4} in a flat situation, since it already has good shapes, we can drop \mj{7z} and go for menzen riichi with pinfu. If we are agaritoppu, pinfu itself becomes less important and being able to use \mj{7z} becomes more important. Our instinct is to drop any of the ryanmens and keep two pairs, but we should drop \mj{2s} pair instead -- so as to not lose the chance to directly call riichi on any draws of \mj{69m47p}. After dropping \mj{2s} pair, any of \mj{78m56p} draws gives us perfect 1-away \citep[][pp.~179--180]{kawamura-naki}.

\tehai{12378m 56p 22678s 77z}[Dora \mj{3m}][pinfu-4]


%---------------------------------------------
\subsection{Iipeiko} \label{subsec:iipeiko}

Since iipeiko worth only one han and is limited to menzen hands only, usually we do not go against efficiency and choose bad shapes to force iipeiko. Shapes such as \mj{33445p} only have a fifty-fifty chance of giving us iipeiko. Shapes such as \mj{33455p} guarantee us iipeiko when our hand is completed, but it is very hard to get this \mj{35p} kanchan filled since we already have one copy of its acceptance. Therefore, we keep ryanmens over \mj{33455p}, and we keep other kanchans over \mj{33455p} when one more \mj{4p} is dropped \citep{gendai-ipk}. In this regard, we drop \mj{79m} from~\ref{mj:ipk-0} \citep[][pp.~121--122]{ooi-gundou}.

\tehai{77899m 11166p 2378s}[North seat, turn 6, dora \mj{8p}][ipk-0]

When our hand has no dora or has dora 3, we can lean a little bit towards iipeiko chance when choosing which floater to keep. Hand \ref{mj:ipk-1} has three floaters \mj{26m8p}. In terms of acceptance, \mj{2m} is weaker since we already have \mj{13m} in our hand and it is less likely to draw them than to draw \mj{79p}. However, drawing \mj{13m} gives us iipeiko chance and makes our hand worth of fighting when it has no dora and allows us to dama when we have dora 3. Therefore, keep \mj{2m} when we have no dora or dora 3. When we have dora 1 or 2, we should just go straight for riichi and drop \mj{2m} for maximal acceptance \citep[][p.~32]{zero-teyaku}.

\tehai{12236m 2348p 35799s}[\mj{1m} is once cut][ipk-1]

However, for hands like \ref{mj:ipk-2} which does not have its head confirmed, aiming iipeiko with \mj{1223m} becomes less feasible because, if we arrive at 1-away with ryanmens, we keep \mj{22m} and drop \mj{13m} to confirm our head. Drop \mj{8p}.

\tehai{12237m 2348p 23079s}[\mj{1m} is once cut, dora \mj{9s}][ipk-2]


\subsubsection*{Three Skipping Pairs}

For shapes such as \mj{113355p} where we have three pairs, efficiency wise, dropping the middle tile gives us the most acceptance as is shown in Table~\ref{tab:santobitt-1}. However, if we keep this 6-tile shape, we have iipeiko upon \mj{24p} draw. Therefore, the discussion of this shape will involve how much we want iipeiko even if that means sacrificing efficiency \multicitep{zero-teyaku, p.~35; yoteru-tobitoitsu; uzaku-haikouritsu, pp.~114--117}.

\begin{table}
  \centering
  \begin{tabular}{lll}
    Shape & Acceptance of \mj{3m} drop & Acceptance of \mj{8s} drop \\
    \hline
    \mj{113355m78s} & \mj{1245m69s} 20 tile & \mj{12345m} 14 tile \\
    \mj{113355m68s} & \mj{1245m7s} 16 tile & \mj{12345m} 14 tile \\
    \mj{113355m89s} & \mj{1245m7s} 16 tile & \mj{12345m} 14 tile \\
    \mj{113355m88s} & \mj{1245m8s} 14 tile & \mj{12345m} 14 tile \\
  \end{tabular}
  \caption{Acceptance of shapes resulting from one tile drop from various combinatinos of three skipping pairs and taatsus.}
  \label{tab:santobitt-1}
\end{table}

For example, the drop candidates of \ref{mj:santobitt-1} are \mj{3m9p}. Their merits are,

\begin{itemize}
  \item Merits of \mj{3m} drop
  \begin{itemize}
    \item \mj{45m} draws give us \mj{1m9p} shanpon wait which has high win rate
  \end{itemize}
  \item Merits of \mj{9p} drop
  \begin{itemize}
    \item Iipeiko chance on \mj{24m} draw
    \item After done dropping this pair, we have space for safe tiles
  \end{itemize}
\end{itemize}

If our only goal is to riichi, drop \mj{3m} and hope that our wait is \mj{1m} and/or \mj{9p}. If there are one or more \mj{1m9p} visible on the table which makes our wait near dead, or we are in a point situation which favors dama, then drop \mj{9p} and rely on iipeiko chance \citep[][p.~35]{zero-teyaku}.

\tehai{113355m 123 99p 567s}[][santobitt-1]

Hand \ref{mj:santobitt-2} has its three skipping pairs closer to the middle. Unlike the possible \mj{1m9p} shanpon wait in \ref{mj:santobitt-1}, the best shanpon wait we can get from \ref{mj:santobitt-2} is \mj{2m1p} and it already has lower win rate than \mj{1m9p}. In such cases, we lean towards to dropping \mj{1p} \citep{yoteru-tobitoitsu}.

\tehai{224466m 11p 789s 123s}[][santobitt-2]

For hands like \ref{mj:santobitt-3} which has tanyao confirmed, we can still either aim for iipeiko, which has theoretically the highest expected value, or drop \mj{4p} and hope for ryanmen improvement of \mj{7s}. Since we have tanyao confirmed, we can call and wait for improvements even after we reach tenpai. The only ryanmen improvement of \mj{224466p} is \mj{7p}, while \mj{22466p77s} improves on \mj{7p6s} (\mj{78s} cannot accept \mj{9s} if we force tanyao.) Therefore, in situations such as agaritoppu where hand value is not to be maximized, drop \mj{4p}; otherwise drop \mj{7s} \citep[][p.~117]{uzaku-haikouritsu}.

\tehai{678m 224466p 23477s}[][santobitt-3]


%---------------------------------------------
\subsection{Honitsu and Chinitsu}

In the following cases, we ignore general tile efficiency and force honitsu,

\begin{itemize}
\item Acceptance is narrow as there are many bad shapes and thus far from menzen tenpai \citep{shibukawa-honitsu}
\item Hand value is mangan (or near mangan) even after call \citep{hirasawa-honitsu}
\item Hand is still cheap (below 2 han) if we go straight for riichi \citep{hirasawa-honitsu}
\item We have two pairs of yakuhai, which we call \textit{yakuyaku} \citep[][pp.~96--97]{zero-teyaku}
\item We have one pair of yakuhai and the rest of our hand is abysmal \citep[][pp.~97--98]{zero-teyaku}
\end{itemize}

Honitsu can only use tiles from one suit and honor tiles, so, in terms of pure efficiency, forcing honitsu sacrifices efficiency. Generally, this loss in efficiency is justified by its value (2 han minimum, easily becomes mangan with dora or yakuhai) and its defense (honor tiles generally have lower deal-in rate.) When we have abysmal haipais, forcing honitsu is usually a good option, value wise and defense wise.


\subsubsection*{Forcing Honitsu in Case of Yakuyaku}

Mahjong is a game of building mangan. If we have two yakuhai pairs, we can maximize our hand value and easily reach mangan by forcing honitsu \citep[][pp.~96--97]{zero-teyaku}. The value of hand \ref{mj:honitsu-1} is 1000 or 2000 if we just follow pure efficiency. If we force honitsu, it can be 3900 or 7700, way bigger than before. Drop \mj{8s}.

\tehai{7m 12456p 228s 44557z}[East-1, North seat, turn 1, dora \mj{3z}][honitsu-1]

Hand \ref{mj:honitsu-2} has a completed souzu block. If we are super early in a round like before turn 4 and yakuhai comes out, we can call pon and drop the entire souzu block to force honitsu \citep[][p.~97]{zero-teyaku}. This is the extent we can deny pure efficiency and maximize hand value.

\tehai{12456p 123s 44557z}[East-1, North seat, haipai, dora \mj{3z}][honitsu-2]

Hand \ref{mj:honitsu-yakuyaku-1} has 6 blocks. With 4 pairs in our hand, we can reach mangan with either honitsu or toitoi. In this regard, the weakest tiles are \mj{56m35z}. From there, we have two options \citep[][pp.~98--103]{tetsusenryaku},

\begin{itemize}
\item Drop \mj{6m} for \mj{0m} acceptance and maximal toitoi and honitsu chance.
\item Drop \mj{5z} if we are okay with a 2000 hand. Drop \mj{5z} before \mj{3z} because \mj{3z} is jikaze.
\end{itemize}

\tehai{56m 22p 2499s 113566z}[East-1, West seat, turn 1, dora \mj{8p}][honitsu-yakuyaku-1]

Hand \ref{mj:honitsu-yakuyaku-2} is similar to \ref{mj:honitsu-yakuyaku-1}, but its manzu part becomes a kanchan. In this case, \mj{46m} kanchan becomes even weaker and we can drop it without much doubt \citep[][pp.~98--103]{tetsusenryaku}.

\tehai{46m 22p 24699s 11566z}[East-1, West seat, turn 1, dora \mj{8p}][honitsu-yakuyaku-2]


\subsubsection*{Think Twice When There is No Yakuhai Pair}

When there are no yakuhai pairs in our hand, forcing honitsu will mostly give us 2000 or 3900 with one dora. This holds as well for cases where we have otakaze pairs. Therefore, we are less incentivized to force honitsu for offense. For defense purposes, though, honitsu still remains an option.

We already have a head and iipeiko in \ref{mj:honitsu-3}. Having iipeiko already confirmed means that if we reach riichi this hand we already have 2 han. Taking the possibility of tanyao and pinfu into consideration, staying menzen has more value than forcing honitsu and opening our hand, especially when there are no yakuhai pairs and the only thing that could increase our value is aka dora. Still, there is a slim chance of chinitsu. Therefore, we drop \mj{37z} first and keep \mj{1m} \citep{shibukawa-honitsu}.

\tehai{144556688m 8p 26s 37z}[South-1, South seat, turn 1, dora \mj{5p}, last place][honitsu-3]

Hand \ref{mj:honitsu-4} has 5 lone honor tiles. Drawing tiles near floating number tiles is generally faster than drawing duplicate copies of honor tiles. In addition, there are no head confirmed in our hand; even if we draw duplicates of honor tiles, calling pon on them does not speed our hand up. Therefore, drop honor tiles; drop yakuhais first, in particular, since otakaze pairs can still contribute to pinfu \citep{yuusee-honitsu}.

\tehai{48m 23457p 25s 12456z}[East-1, West seat, turn 1, dora \mj{1m}][honitsu-4]

Hand \ref{mj:honitsu-5} and \ref{mj:honitsu-6} are both hands with narrow acceptance and mostly consisting of bad shapes of non-dominant suit. While it is true that we do not have yakuhai pairs and increasing our hand value is difficult, defense wise the otakaze pairs provide us a lot of safety. Therefore, we force honitsu for defense purpose here \citep{shibukawa-honitsu}.

\tehai{114478m 246s 44567z}[East-3, East seat, turn 2, dora \mj{2m}, \mj{4m} is once cut][honitsu-5]

\tehai{35m 57p 123449s 2244z}[East-3, East seat, turn 2, dora \mj{8s}, \mj{4z} is once cut][honitsu-6]

Hand \ref{mj:honitsu-alllast} has no yakuhai pair at all. Still, this is agaritoppu, and we must force some yaku to take first place. For this kind of messed up hand, honitsu becomes our hope, with chanta also slightly probably. Therefore, drop \mj{67m}, and prepare to call everything possible \citep[][pp.~37--38]{kawamura-naki}.

\tehai{67m 13p 112279s 1445z}[All-last, West seat, turn 1, dora \mj{3z}, 500 points to overtake first place][honitsu-alllast]


\subsubsection*{When There is Only One Yakuhai Pair}

Having only one yakuhai pair means that we are in the middle ground between the sweetspot of honitsu value increase and menzen riichi. Deciding whether to force honitsu or not will largely depend on how good the rest of our hand is.

With one yakuhai pair and two lone yakuhai tiles, we are in a state referred to as ``ready for yakuyaku.'' In this state, we can tolerate dropping floating 1289 tiles of lesser suits \citep[][p.~97]{zero-teyaku}. In this line of thought, we drop \mj{8p} from \ref{mj:honitsu-7} and \mj{8s} from \ref{mj:honitsu-8}.

\tehai{1126m 8p 56889s 1566z}[East-1, North seat, turn 1, dora \mj{2s}][honitsu-7]

\tehai{34m 1455689p 8s 4556z}[East-1, North seat, turn 1, dora \mj{2m}][honitsu-8]

Likewise, we drop \mj{28s} from \ref{mj:honitsu-9} and get prepared for pinzu honitsu \citep{hirasawa-honitsu}.

\tehai{347m 34479p 28s 1157z}[Dora \mj{1p}][honitsu-9]

On the contrary, we do not force honitsu for \ref{mj:honitsu-10}. We only have one lone yakuhai tile and it is difficult to make more blocks in manzu suit. Drop \mj{6z} \citep[][p.~98]{zero-teyaku}.

\tehai{35999m 379p 48s 5556z}[East-1, North seat, turn 1][honitsu-10]

Hand \ref{mj:honitsu-11} is also not fit for honitsu. It is true that we have one yakuhai pair and two lone yakuhai tiles, but the shapes in lesser suits are all ryanmen. Therefore, drop \mj{67z} \citep[][p.~98]{zero-teyaku}.

\tehai{67m 23p 223056s 5567z}[East-1, North seat, turn 1, dora \mj{1s}][honitsu-11]

Hand \ref{mj:honitsu-12} has no shapes in suits other than pinzu. Even though we have iipeiko, with \mj{7z} pair in our hand, its value is mostly 2600. If we force honitsu, we will have at least 3900, and drawing another copy of jikaze \mj{4z} increases our mangan chance. Duplicate copies of \mj{3z} can also allow us pivot to toitoi. Considering also the slim chance of \mj{0m}, we drop \mj{7s} first \citep[][pp.~104--109]{tetsusenryaku}。

\tehai{5m 22333447p 7s 3477z}[East-1, North seat, turn 1, dora \mj{6z}][honitsu-12]


\subsubsection*{Drop The Suit with Second Most Tiles First}

Hand \ref{mj:honitsu-13} is the moment when we called pon on the first drop of \mj{4z} from oya. All of \mj{6m46p} are useless since we are aiming for honitsu here. If we drop \mj{6m} here, soon we will start dropping pinzu tiles, making others realize that we are waiting on souzu tiles. Therefore, the trick here is to drop pinzu tiles first, and delay the timing when opponents figure out what the suit we are waiting on is \citep[][pp.~118--119]{suzuki}. If unfortunately we draw manzu tiles, just tsumogiri.

\tehai{6m 46p 14467s 336z -444'z}[East-1, North seat, turn 1, dora \mj{2m}][honitsu-13]

Hand \ref{mj:honitsu-14} is the moment right after we called pon on dabuton. In cases where there is no difference in the number of pinzu and souzu tiles. drop the outermost tiles to make our river look like we are just following pure efficiency and not forcing honitsu. Drop \mj{8p} \citep[][pp.~120--121]{suzuki}.

\tehai{1179m 68p 36s 447z -111'z}[East round, East seat, turn 2][honitsu-14]


%---------------------------------------------
\subsection{Yakuyaku}

\textbf{Actively seek ways to make yakuyaku mangan }\citep{yoteru-yakuyaku}. As mentioned previously, yakuyaku gives us 2 hans that is the sweetspot where, for 2 additional hans, our value increases from 2000 to 7700, a whopping 5700 points of difference. Honitsu is the most common way to achieve mangan, but keep in mind that toitoi is also an option.

Hand \ref{mj:yy-1} has two yakuhai pairs and \mj{34s} is our only ryanmen. If we keep \mj{34s} and follow pure efficiency, we will mostly end up 2000 (1000 if we do not get to pon one of the yakuhai pairs.) If our hand is slow anyway, it is better to force honitsu to maximize our value. We can drop any of \mj{68p34s}.

\tehai{2579m 68p 34s 135577z}[East-1, South seat, turn 1, no dora][yy-1]

Hand \ref{mj:yy-2} has four pairs and toitoi becomes more viable than honitsu. Middle tiles are harder to call than outside tiles, so we drop \mj{3m} here first.

\tehai{223m 119p 28s 135577z}[East-1, South seat, turn 1, no dora][yy-2]


\subsubsection*{Cases Where We Have One Dora}

With one dora in our hand, the increase in hand value now becomes limited. Suppose that we have all yakuhai ponned, our hand value is already 3900 minimum, and forcing honitsu will only make it 8000, since its value is capped at mangan. Therefore, we do not force 2 han yakus, and for hands like \ref{mj:yy-3} where we also have a ryanmen in lesser suits, it is perfectly fine to proceed with 3900. Drop \mj{13z}.

\tehai{67m 1330p 46s 135577z}[East-1, South seat, turn 1, dora \mj{2m}][yy-3]

Likewise, we do not force honitsu in \ref{mj:yy-4}. However, there is also a great possibility of 123 sanshoku and chanta which are 1 han yakus that brings us mangan. Drop \mj{4s}.

\tehai{137m 3799p 124s 5577z}[East-1, South seat, turn 1, dora \mj{1m}][yy-4]

Hand \ref{mj:yy-5} indeed has one dora and two yakuhai pairs, but note that we are close to tenpai yet none of our yakuhai pairs are ponned, which means that, when we reach tenpai, we either have one pair as our head, or wait on yakuhai shanpon, and our hand value is still 2000 minimum. Therefore, we force honitsu here and drop \mj{67s}.

\tehai{13m 67s 5577z -6'0799'9m}[East-1, South seat, turn 7, dora \mj{8p}][yy-5]

Hand \ref{mj:yy-6} likewise is a hand where we did not get to pon our yakuhai pairs in time. Drop \mj{4s} and force toitoi for mangan.

\tehai{405s 25577z -1'11p888'm}[East-1, South seat, turn 7, dora \mj{8p}][yy-6]


\subsubsection*{Cases Where We Advance Yakuyaku Without Forcing Yakus}

When our hand is slow, the increase in value from forcing 2 han yaku outweighs loss in speed, but this is not the case when our hand already has enough good shapes like \ref{mj:yy-7}. Drop \mj{2z} and aim for a fast 2000 to get to next round.

\tehai{2467m 223p 56s 25577z}[East-1, South seat, turn 1, dora \mj{1m}][yy-7]

If we are in huge lead and just want to move to the next round, or we are agaritoppu, forcing honitsu or toitoi is also not a good option.


%---------------------------------------------
\subsection{Sanshoku Doujun}

Asada Tetsuya once said, ``look for sanshoku\footnote{There is another yaku called \textit{sanshoku doukou} which requires three triples with the same number, i.e., \mjfn{333m}, \mjfn{333p}, and \mjfn{333s}. Since sanshoku doukou is far less common than sanshoku doujun, in most cases, \textit{sanshoku} refers to sanshoku doujun.} as soon as you see haipai.'' In the early days of mahjong when riichi (in the sense that we declare our tenpai in the middle of a round), ippatsu, uradora and stuff didn't exist, building your hand with yakus such as sanshoku was important. As riichi and uradora became widespread, the importance of sanshoku dropped, and it is not as popular nowadays. Some would even go so far as to say that sanshoku is a ``luck'' yaku.

Still, being a 2 han yaku (1 han if our hand is open) and a yaku that only uses 3 blocks, sanshoku is still a useful yaku. The following criteria will help us decide when to force sanshoku \citep{gendai-ssk}:

\begin{itemize}
\item How easy it is to get sanshoku confirmed.
\item Loss of acceptance in the case when it is hard to confirm and we have to force it.
\item The merit of increase in value with 2 hans of sanshoku or just the yaku itself.
\end{itemize}

When our hand has no dora, building value to make it worth fighting becomes crucial. When our hand has a lot of dora but also far away from tenpai, having a yaku to back up our calling is also important. In these two cases, we tend to force sanshoku.

If our hand is close to tenpai, or our head is not confirmed, sanshoku becomes harder to score. Hand \ref{mj:ssk-1} already has 7 tiles necessary for 123 sanshoku, but there is no confirmed head. The most likely outcome is drawing some good shape around \mj{6778s} yonrenkei and drop \mj{1m} to confirm \mj{22m} as our head. Going for pinfu is more efficient \citep[][p.~67]{zero-teyaku}.

\tehai{122m 12357p 236778s}[South seat, turn 5, dora \mj{2p}][ssk-1]

Hand \ref{mj:ssk-2} too has two blocks completed. With our hand being headless, we cannot have all 123 sanshoku parts without destroying other blocks \citep[][p.~67]{zero-teyaku}.

\tehai{235678m 13567p 123s}[South seat, turn 5, dora \mj{8m}][ssk-2]


\subsubsection*{Force Yaku to Avoid Rii-Nomi}

Hand \ref{mj:ssk-3} and \ref{mj:ssk-4} would be perfect 1-away if we drop \mj{7p}. For \ref{mj:ssk-3}, since our value is 2600 minimum (riichi tanyao,) we are at that sweetspot of increase in value brought by riichi. Therefore, we prioritize speed over value and take perfect 1-away. For \ref{mj:ssk-4}, however, our hand value is 1300 minimum and increase in value brought by riichi is smaller. The benefit of 2 han from 567 sanshoku therefore justifies the loss of \mj{16m} acceptance after we drop \mj{6m} \citep{clearrain-ssk}.

\tehai{22566m 457p 222567s}[Turn 6, no dora][ssk-3]

\tehai{11566m 457p 222567s}[Turn 6, no dora][ssk-4]

Dropping either \mj{4m} or \mj{6s} from \ref{mj:ssk-5} gives us the same number of acceptance. According to Table~\ref{tab:fkgtt-2}, dropping \mj{6s} and confirming ryanmen gives us the most good-wait acceptance. However, dropping \mj{4m} confirms 234 sanshoku. In terms of expected value, confirmed sanshoku is greater, so we drop \mj{4m} here \citep[][p.~65]{zero-teyaku}.

\tehai{244m 234p 234667s 44z}[South seat, turn 3, dora \mj{3z}][ssk-5]

For the same tiles, if we have two doras like \ref{mj:ssk-6}, we should just go for good waits and drop \mj{6s} \citep[][p.~65]{zero-teyaku}.

\tehai{244m 234p 234667s 44z}[South seat, turn 3, \mj{4z}][ssk-6]

If our hand has only one dora, there is no big difference in terms of expected return. Adjust your judgement according to actual table situation.


\subsubsection*{Go Big or Go Home with Bad Haipai}

Hand \ref{mj:ssk-7} is 4-away and has only one good shape. If our hand is messed up anyway, we can try to force yakus to increase value or move on with quick hands. From this hand, we can see yakuhai and tanyao. Note that with 5 tiles of 123 sanshoku already in our hand, sanshoku is also viable, especially when we do not have enough blocks. Therefore, drop the most useless tile \mj{9s}. Even though it is in a \mj{899s} block, notice that it is not part of any of the yakus mentioned previously \citep[][p.~64]{zero-teyaku}.

\tehai{1356m 124p 15899s 56z}[South seat, turn 1][ssk-7]


\subsubsection*{Keep Seeds of Sanshoku Over Bad Blocks}

Hand \ref{mj:ssk-8} has 6 blocks. According to 6 block theory, we drop floater \mj{8s} and take maximal acceptance of 22 tiles. However, dropping \mj{2p} and making lower\footnote{The ``upper'' part of a suit refers to tiles with greater number and the ``lower'' part refers to tiles with smaller number.} pinzu part one block only lose 2 tiles of acceptance, and we get to keep the ``seed'' of 789 sanshoku. Therefore dropping \mj{2p} is good \citep[][p.~64]{zero-teyaku}. Another way to look at this is to treat \mj{8s} as a block, so we are still maintaining 6 blocks in our hand, and we are ready to drop \mj{8s} when we are closer to tenpai.

\tehai{79m 2244789p 458s 33z}[South seat, turn 3, dora \mj{9p}][ssk-8]

Likewise, in the case of \ref{mj:ssk-9}, it would be very difficult to make two mentsus from \mj{2446m}, so we drop \mj{2m} and keep \mj{7p} with allows us to aim for 789 sanshoku \citep{gendai-ssk}.

\tehai{2446789m 7p 113578s}[][ssk-9]


\subsubsection*{6 Block Theory Sanshoku Edition}

If you have read \ref{subsec:6-tile}, you will notice that dropping \mj{8m} from \ref{mj:ssk-10} will give us maximal acceptance as well as guaranteed pinfu. Notice that we also have 7 tiles required for 789 sanshoku as well, but since our shapes in upper parts of suits are ryanmens, we could easily end up not getting sanshoku, and keeping \mj{78m78s} makes us los 3 less tiles of acceptance than \mj{8m} drop. Therefore, we will keep \mj{78m} only when we have no dora and have to chase for value; if we have at least one dora, we should drop \mj{8m} \citep{gendai-ssk}.

\tehai{2334578m 789p 5578s}[][ssk-10]

Hand \ref{mj:ssk-11} has 6 blocks \mj{334m} + \mj{234p} + \mj{46p} + \mj{88p} + \mj{34s} + \mj{88s}. According to 6 block theory, we drop \mj{4p} from the weakest block and proceed with 5 blocks which gives us the maximal chance of perfect 1-away. Notice that we have 7 tiles required for 234 sanshoku. By dropping \mj{3m} early and confirming \mj{34m}, the chance of 234 sanshoku increases\footnote{Because with \mjfn{334m} shape, drawing \mjfn{3m} withh ruin this ryanmen. In case where there are no head in the rest of our hand, we must drop \mjfn{4m} and sacrifice this ryanmen.}, and it becomes easier for others to drop sotogawa \mj{2m} \citep[][p.~65]{zero-teyaku}.

\tehai{334m 2344688p 3488s}[South seat, turn 4, dora \mj{4z}][ssk-11]

Hand \ref{mj:ssk-12} has 6 blocks and 6 tiles required for 567 sanshoku. Unlike previous cases where we can reduce taatsus involved in sanshoku to simpler taatsus, we must either drop a sanshoku taatsu entirely or drop \mj{4p} which reduced our perfect 1-away chance. On top of that, we already have \mj{0m} in our hand, so there is no need to force sanshoku here. Drop nidouke \mj{67p} \citep{gendai-ssk}.

\tehai{06m 34467p 1112367s}[][ssk-12]

Hand \ref{mj:ssk-13} looks like it has 567 sanshoku confirmed, but to score 567 sanshoku we have to keep three all kanchans, which requires incredible luck to get them filled. Therefore, unless we absolutely have to force some value out of this hand, follow pure efficiency and drop \mj{5m} \citep{gendai-ssk}.

\tehai{5788m 223457p 3357s}[][ssk-13]


\subsubsection*{Choosing between Floaters}

Hand \ref{mj:ssk-14} is at 1-away. Since we already have \mj{789m789p}, keeping \mj{9s} allows us to draw \mj{78s} and score sanshoku. However, compared to other floaters \mj{3m3p}, \mj{9s} has 8 less tiles of acceptance, and it cannot form a ryanmen. Therefore, we will only keep \mj{9s} when (1) we have no dora and must force value or (2) it is still early in a round. When we have at least 2 doras, or it is already in the middle of a round, simply go for maximal efficiency and drop \mj{9s}. If we only have one dora, either is fine \citep{yoteru-ssk}.

\tehai{3789m 3789p 229s 333z}[South seat][ssk-14]

Hand \ref{mj:ssk-15} is further away from tenpai than \ref{mj:ssk-14}. Likewise, keep \mj{9s} when we have no dora. When we have dora 3, we can also rely on sanshoku to give us yaku to open our hand and speed it up. In other cases, drop \mj{9s} \citep{yoteru-ssk}.

\tehai{3789m 3379p 3559s 33z}[South seat][ssk-15]

Hand \ref{mj:ssk-16} has only 4 blocks and currently does not have any yaku confirmed. Dropping \mj{1s} means that, if we dont get \mj{12p} penchan filled up first, we will end up rii-nomi. For hands without any dora, drawing around \mj{7s} would not be desirable, so we drop it here and keep \mj{1s} and aim for 123 sanshoku \citep{gendai-ssk}.

\tehai{123567m 1255p 17s56z}[][ssk-16]


\subsubsection*{Choosing between Seeds of Sanshoku}

\textbf{Sanshoku with a floater as its middle tile is weak} \citep[][pp.~65--66]{zero-teyaku}. Take \ref{mj:ssk-17} for example. \mj{3p} can be the seed of lower part sanshoku and \mj{7p} is dora, so we cannot drop either. Thus, it becomes a decision between \mj{1m} and \mj{4m}, which is equivalent to choosing between 123 and 234 sanshoku. As shown in Table~\ref{tab:ssk-1}, the greatest difference is \mj{1p} draw which, if we dropped \mj{1m} previously, gives us no yaku at all. Dropping \mj{4m} always gives us a yaku, so we drop \mj{4m} and aim for 123 sanshoku. We thus conclude that, with \mj{3p} floater, 234 sanshoku is weaker than 123 sanshoku.

\tehai{123405m 37p 122334s}[South seat, turn 5, dora \mj{7p}][ssk-17]

\begin{table}[h]
  \centering
  \begin{tabular}{lll}
     & Drop \mj{1m} & Drop \mj{4m} \\
    \hline
    \mj{1p} draw & No yaku & Sanshoku \\
    \mj{2p} draw & Pinfu, (sanshoku) & Pinfu, (sanshoku) \\
    \mj{4p} draw & Pinfu, (sanshoku) & Pinfu
  \end{tabular}
  \caption{Comparing \mj{1m} drop and \mj{4m} drop from \ref{mj:ssk-17}. Sanshoku in parentheses means that it is not confirmed due to ryanmen waits.}
  \label{tab:ssk-1}
\end{table}

Hand \ref{mj:ssk-18} and \ref{mj:ssk-19} are both tenpai if we drop the floater, which gives us dora 1 bad wait. We have learned previously that, ryanmens formed around \mj{7m} only gives 678 sanshoku chance, and any kanchan formed from it is undesirable. Just drop \mj{7m} from \ref{mj:ssk-18} and accept tenpai. On the other hand, \mj{6m} accepts \mj{578m} all of which give us yaku. Considering also possible \mj{4p} draw, we actually have 16 tiles of upgrade if we drop \mj{1p}, a reasonable trade-off at turn 3. Therefore, we deny this tenpai in \ref{mj:ssk-19}. 678 sanshoku is weaker and does not justify for going back in shanten if the floater is the middle one \mj{7m} \citep[][p.~66]{zero-teyaku}.

\tehai{1237m 13678p 67778s}[South seat, turn 3, dora \mj{1m}][ssk-18]

\tehai{1236m 13678p 67778s}[South seat, turn 3, dora \mj{1m}][ssk-19]


\subsubsection*{Ryoutenbin}

\textit{Ryoutenbin} is a Japanese term which generally translates to ``sitting on the fence,'' and it refers, in the context of mahjong, to a situation where, with 13 tiles, there are possibilities of two yakus that cannot coexist. Which yaku we eventually score depends on the subsequent draws. There are ryoutenbins of sanshoku and ittsu and also ryoutenbins of different sanshokus.

If we drop \mj{4m} from \ref{mj:ssk-20}, with the shape \mj{5678m67p67s}, drawing \mj{5p5s} gives us 567 sanshoku chance, and drawing \mj{8p8s} gives us 678 sanshoku chance. This shape is therefore the ryoutenbin of 567 and 678 sanshoku. While it is true that dropping \mj{67p} gives us sanmenchan and maximal acceptance, in a flat situation, we can sacrifice 3 tiles of acceptance and take sanshoku ryoutenbin \citep{rb1}.

\tehai{45678m 3367p 34567s}[][ssk-20]

Hand \ref{mj:ssk-21} becomes ryoutenbin of 123 and 234 sanshoku if we drop \mj{3p}. While it indeed is 4 tiles less of acceptance than \mj{1s} drop which gives us perfect 1-away, our sanshoku chance also increases, which gives us way higher expected value than \mj{14s} drop\footnote{The difference in expected value justifies loss of acceptance up to dora 2. When we have dora 3, the difference in expected value is very small and we should go for maximal acceptance.}. Drop \mj{3p} \citep{clearrain-ssk}.

\tehai{23067m 23399p 1234s}[Turn 6][ssk-21]

Hand \ref{mj:ssk-22} also becomes ryoutenbin of 345 and 456 sanshoku if we drop \mj{4p}. In this case, we do not have pinfu, thus there is a higher chance of undesirable rii-nomi, which further justifies loss in 4 tiles of acceptance in exchange of increased sanshoku chance \citep{gendai-ssk}.

\tehai{3456m 445 999p 45 88s}[][ssk-22]

Hand \ref{mj:ssk-23} also becomes ryoutenbin of 567 and 678 sanshoku if we drop any of \mj{2m3p}. However, since this is a headless hand with anko, dropping \mj{2m3p} actually loses 11 tiles of acceptance, a wide gap in acceptance that do not justify increased sanshoku chance. Considering that it is more likely for 8 of 678 sanshoku to come out and complete our hand, we drop \mj{5s} here \citep{clearrain-ssk}.

\tehai{22267m 33367p 5678s}[Turn 6][ssk-23]

If we drop \mj{3p} from \ref{mj:ssk-24} for 567 and 678 sanshoku ryoutenbin, we lose \mj{4p} which gives us mentanpin dora 1 mangan tenpai. In previous cases where we would trade 4 tiles of acceptance for sanshoku chance, those 4 tiles do not contribute to value, but this \mj{4p} puts us from 3900 to 7700, an increase too great to lose. Therefore, drop \mj{8m}, since 567 sanshoku is more likelier than 678 sanshoku with \mj{067p} block \citep{gendai-ssk}.

\tehai{566678m 306778p 67s}[][ssk-24]


%---------------------------------------------
\subsection{Ikkitsuukan}

\textit{Ikkitsuukann} in Japanese means ``to strike through in one breath,'' and refers to a yaku we score when we have all nine tiles of a suit. It is often shortened to \textit{ittsu}, or translated into English as \textit{pure straight.} Ittsu has the same value as sanshoku, which is 2 han for menzen and 1 han for open hands. Since ittsu requires all nine tiles from a suit, there are only 3 possible patterns of ittsu (1--9 tiles in manzu, souzu, and pinzu), less than sanshoku (123, 234, 345, 456, 567, 678, 789 tiles), it is a less common yaku.


\subsubsection*{When Ittsu is Viable}

We can go for ittsu if our hand has no dora, it is still far away from tenpai, and we have at least 5 tiles in a suit. In such case, we keep floaters in the dominant suit and drop other floaters, even if they have more acceptance. Take \ref{mj:ittsu-1} for example where we are choosing one from \mj{7m9p8s}. Acceptance wise, \mj{9p} is the weakest, but this hand has no tanyao no pinfu no dora whatsoever. Thus, keep \mj{7m} for acceptance of \mj{6m} dora, and keep \mj{9p} for a slim chance of pinzu ittsu. Drop \mj{8s} which is the one that contributes the least value \citep[][p.~68]{zero-teyaku}.

\tehai{3457m 2345679p 118s}[Dora \mj{6m}][ittsu-1]

We must drop either of \mj{1m1s} from \ref{mj:ittsu-2}. Manzu ittsu is further than \ref{mj:ittsu-2} with only 5 manzu tiles, but it is still slightly better than \mj{1s} which gives us nothing but the chance of a penchan or a kanchan. Likewise, drop \mj{2p} from \ref{mj:ittsu-3} \citep[][p.~69]{zero-teyaku}.

\tehai{14578m 1136p 1899s 5z}[Dora \mj{5m}][ittsu-2]

\tehai{267899m 2678p 1399s}[][ittsu-3]

When comparing \mj{12m} penchan and \mj{13s} kanchan in \ref{mj:ittsu-4}, efficiency wise, it is better to drop penchan and wait for ryanmen improvement of kanchan. However, dropping \mj{12m} means that we lose ittsu chance on \mj{45m} draws. Therefore, we drop \mj{1s} \citep[][pp.~68--69]{zero-teyaku}.

\tehai{1267899m 678p 1399s}[No dora][ittsu-4]

Dropping ryanmen in favor of penchan may sound counterintuitive, but for cases such as \ref{mj:ittsu-5} where we have no dora and must force ittsu, two penchans that confirm ittsu is better than a penchan and a ryanmen that give us only rii-nomi. Therefore, drop \mj{45s} \citep{gendai-ittsu}. If our hand has 1 dora, follow pure efficiency and drop any of the penchans.

\tehai{1245689m 99p 45567s}[No dora][ittsu-5]


\subsubsection*{Cases Where Ittsu is Not Viable}

We rarely go for ittsu when we have 1 dora, or we can have tanyao confirm if we remove non-tanyao blocks, even if we have 7 tiles in a suit. While it is true that ittsu has more value than tanyao when stay closed, the shapes in the ittsu suit are fixed and cannot upgrade to ryanmens. If we have at least 1 han, this inflexibility becomes undesirable.

Hand \ref{mj:ittsu-6} has 7 distinct manzu tiles, but dropping \mj{9m} will not make us lose the acceptance of  \mj{27m}, and \mj{7m} allows us to drop \mj{1m} and confirm tanyao. It is thus the right choice to drop \mj{9m}. \citep[][p.~69]{zero-teyaku}。

\tehai{1345689m 33468p 06s}[][ittsu-6]

Likewise, if there is 1 dora in \ref{mj:ittsu-7}, drop \mj{9m} and go straight forward. Only drop \mj{246p} ryankanchan when there is no dora \citep[][p.~70]{zero-teyaku}.

\tehai{1345689m 11246p 78s}[][ittsu-7]

Dropping \mj{19m} from \ref{mj:ittsu-8} gives us tanyao, so we do not need ittsu here. Drop \mj{19m} \citep[][p.~70]{zero-teyaku}.

\tehai{1345679m 3557p 455s}[][ittsu-8]

Hand \ref{mj:ittsu-9} too does not need ittsu if dropping \mj{12m} already gives us pinfu \citep{gendai-ittsu}.

\tehai{1256789m 99p 45567s}[][ittsu-9]


\subsubsection*{Golden 1-Away}

Golden 1-away is a subset of kuttsuki iishanten that is also a ryoutenbin of sanshoku and ittsu. Take \ref{mj:golden-1} for example. If we drop \mj{6s}, \mj{16p} draws give us pinzu ittsu chance while \mj{79s} draws give us 789 sanshoku chance. Taking the iishanten with ittsu and sanshoku chances has higher expected value than \mj{8m} drop which gives us only pinfu, so we drop \mj{9s} here \citep{rb1}.

\tehai{8m 2345789p 556789s}[][golden-1]

Hand \ref{mj:golden-2} too is golden 1-away. Golden 1-away does not always come in shapes like \ref{mj:golden-1} where there is clearly a floating \mj{8m}. For the combination of yonrenkei and ryanmen-toitsu, \ref{mj:golden-2} is also a kuttsuki iishanten.

\tehai{234m 2346789p 344s}[][golden-2]

Hand \ref{mj:golden-3} becomes golden 1-away if we drop \mj{6s}. Dropping \mj{3m} gives us kuttsuki iishanten with two yonrenkei shapes, whose 53 tiles of acceptance is undeniably the most, but its bad wait rii-nomi draws is also the most \citep{shibukawa-ssk}. Since \mj{344m} ryanmen implies 345 sanshoku chance, and we need pinzu yonrenkei for ittsu, \mj{6s} becomes the most valueless tile.

\tehai{344m 2345789p 3456s}[East-1, East seat, turn 5, dora \mj{2z}][golden-3]

Hand \ref{mj:golden-4} becomes golden 1-away if we drop \mj{6m}. However, dropping \mj{3p} is not too bad either, since it still gives us a wide acceptance of mentanpin (unless with unfortunate \mj{1m19s} draws.) Therefore, drop \mj{3p} for the best chance of good wait tenpai \citep{gendai-ittsu}.

\tehai{345688m 3p 2345678s}[][golden-4]

However, if the pair becomes \mj{9m} and tanyao is gone, drop \mj{6m} for value\footnote{The result I got from simulator.}.

\tehai{345699m 3p 2345678s}[][golden-5]


\subsubsection*{Choosing between Sanshoku and Ittsu}

Since both sanshoku and ittsu uses 3 blocks, we cannot have both when we are dropping one block from a hand with 6 blocks. In most cases, \textbf{keeping sanshoku is better} \citep{gendai-ittsu} since sanshoku can coexist with tanyao.

Hand \ref{mj:ittsu-10} looks like 789 sanshoku or pinzu ittsu. However, it has 6 blocks, and we must drop one block from any of \mj{78m}, \mj{23p}, and \mj{56p}. To complete ittsu, we need exactly \mj{14p}, which is one more kind of tiles than 789 sanshoku and more difficult to score. Therefore, we drop ittsu here and drop \mj{5p}. If we draw \mj{6s} and drop \mj{9s9p}, we can even have 678 sanshoku and tanyao at best.

\tehai{3378m 2356789p 789s}[East round, North seat, turn 7, dora \mj{8s}][ittsu-10]

Hand \ref{mj:ittsu-11} already has pinzu ittsu confirmed. If we drop any of the ryanmens and have the other one filled before \mj{79p} does, we end up 5200. On the other hand, if we drop \mj{9p}, we have 123 and 234 sanshoku ryoutenbin, and \mj{4m4s} draws give us tanyao chance. Therefore, we do not go for ittsu here. Drop \mj{9p}.

\tehai{23m 12345679p 2377s}[East round, West seat, turn 8, dora \mj{3p}][ittsu-11]


%---------------------------------------------
\subsection{Toitoiho}

\textit{Toitoiho}, or \textit{toitoi} for short, has 2 hans regardless of hand being open or not, making it highly compatible with other yakus such as yakuhai and easily giving us high value. However, since toitoi hands mostly consist of pairs and triplets, when our opponents call riichi, we have less kinds of tiles and thus less possible safe tiles \citep{yoteru-toitoi}. In this section, we will look at cases and see if it is worth sacrificing tile efficiency for value; for cases where others drop tiles and we are presented the option to open hand and go for toitoi, see Section 3.


\subsubsection*{From 1300 to 5200}

If we drop \mj{6m} from \ref{mj:toitoi-1}, its value is just 1300. If we force toitoi, we have half the tenpai tiles but 4 times the value, a good trade-off. Drop \mj{5m} \citep[][p.~57]{zero-teyaku}。

\tehai{566m 22288p 222s -666'z}[Dora \mj{1z}][toitoi-1]

Even if we have two open sets and we are more exposed to the danger of riichi, we can still force toitoi. Drop \mj{5m} from \ref{mj:toitoi-2} \citep[][p.~57]{zero-teyaku}.

\tehai{566m 88p 222s -22'2p666'z}[Dora \mj{1z}][toitoi-2]

If we already have mangan with ryanmen tenpai, there is no need to foice toitoi for haneman. If we halve our tenpai tiles, our value only becomes 1.5 times higher. Drop \mj{6m} \citep[][p.~57]{zero-teyaku}.

\tehai{566m 88p 222s -22'2p666'z}[Dora \mj{2p}][toitoi-3]

Hand \ref{mj:toitoi-4} too has 4 pairs and we can go for toitoi. Drop \mj{3s7p7m} \citep[][p.~57]{zero-teyaku}.

\tehai{788m 799p 113s 33z -666'z}[Dora \mj{8s}][toitoi-4]

Hand \ref{mj:toitoi-5} is the moment when we draw \mj{2p}. Seeing \mj{123p} already completed, people would drop \mj{2p} without thinking too much. Do keep in mind that, if we keep \mj{2p}, there is a chance for toitoi and 5200 points. If it is before turn 4, drop \mj{3p} and go get that 5200. If it is later than that, drop \mj{2z}, and keep toitoi route as our backup plan \citep[][p.~58]{zero-teyaku}.

\tehai{8899m 123p 233z -2p -666'z}[East round, South seat, turn 3, dora \mj{8s}][toitoi-5]

Hand \ref{mj:toitoi-6} has still needs one more pair for toitoi. Efficiency wise, \mj{3z} would be the right choice, but we most likely end up 1300. The point here is to keep toitoi possible without sacrificing efficiency too much. In this regard, we cannot drop any of the pairs. \mj{3z} would be a good pon candidate when it pairs up or a good tile for defense. \mj{7p} then becomes the least useful tile; indeed, it accepts \mj{8p}, but dropping it also makes \mj{9p} to come out more easily. Therefore, we drop \mj{7p} first; if we draw another pair, we can go drop \mj{45s} \citep[][p.~58]{zero-teyaku}.

\tehai{223m 11799p 45s 3z -77'7z}[Dora \mj{8s}][toitoi-6]


\subsubsection*{Ryanmen 5200 is Sufficient}

\textbf{Ryanmen 5200 is better than bad wait 8000}, which is a theory proposed by Fukuchi Makoto \citep{yoteru-toitoi}. For hands like \ref{mj:toitoi-7} which already has dora 2 and chun, and all the fu from terminal tile triplets making it 40 fu, we can just go take ryanmen wait 5200. Halving our wait tiles just for 2800 points increase is not worth. Drop \mj{5s}.

\tehai{111m 05p 056s -999'm 77'7z}[][toitoi-7]

On the other hand, \ref{mj:toitoi-8} has one less dora and only has 2600. It is perfectly fine to take toitoi for mangan. Drop \mj{6s}.

\tehai{111m 05p 556s -88'8s 77'7z}[][toitoi-8]


\subsubsection*{Forcing Toitoi}

When we do not have all 5 pairs necessary for toitoi, force it anyway if our hand meet any of the following two conditions \citep[][p.~59]{zero-teyaku}:

\begin{itemize}
\item Forcing toitoi gives us 5200 or higher value
\item There are 4 pairs and there is at most 1 middle tile pair among them
\end{itemize}

Hand \ref{mj:toitoi-9} has 4 pairs and no ryanmen, so it would be very slow even if we follow pure efficiency, and we would end up 1000 from \mj{7z} pon. Therefore, we force toitoi with this hand, and drop tiles from the middle. We call pon on any tiles possible.

\tehai{2288m 1146p 18s 5677z}[Dora \mj{6s}][toitoi-9]

Likewise, force toitoi with \ref{mj:toitoi-10}. To deliberately set up sotogawa\footnote{\textit{Sotogawa} in Japanese means ``outside,'' and it refers to the outside of early discards in mahjong context. Sotogawa is generally considered safe because early discards are often not connected or related to other shapes in our hand, and cases where we drop \mjfn{3m} from \mjfn{334m} early are rare (but not impossible.) The act of deliberately dropping \mjfn{3m} from \mjfn{334m} early is called \textit{sakigiri}.}, drop \mj{7m} first.

\tehai{22799m 1179p 1s 5677z}[Dora \mj{6s}][toitoi-10]

We just called \mj{9m} with \ref{mj:toitoi-11} at the premise of toitoi. From opponents' perspective, if someone calls \mj{9m} pon and drops \mj{3p}, our visible parts would look like manzu honitsu, which makes \mj{1p} to come out more likely. Therefore, after we call pon of a suit, we drop a tile from a different suit. If we call \mj{1p} pon first, we drop \mj{7m} instead.

\tehai{1227m 1139p 677z -99'9m}[Dora \mj{6s}][toitoi-11]


\subsubsection*{Cases Where We Do Not Go for Toitoi}

Even though \ref{mj:toitoi-12} has 4 pairs, we do not force toitoi with this hand. The greatest hurdle of toitoi, in this case, is that we cannot use double dora \mj{0m}. Therefore, toitoi will only give us 2600. Forcing tanyao would be 3900 at the very least and would be a far better option. Since honor tiles are generally safer, we drop \mj{1s} pair first \citep[][pp.~128--129]{fkt}.

\tehai{0m 45666p 112246s 44z}[East round, South seat, turn 3, dora \mj{5m}][toitoi-12]


%---------------------------------------------
\subsection{Chiitoitsu}

\textit{Chiitoitsu}, or \textit{chiitoi} for short, is a yaku that was invented in the United States of America and later became a part of Riichi Mahjong ruleset. To get ready for chiitoitsu, our hand must have 6 pairs and 1 lone tile, which also means that chiitoitsu is not compatible with other yakus that consist of sequences and triplets. On top of that, chiitoi 1-away is usually just us waiting on 3 lone tiles which give us 9 tiles of acceptance at most. Except for rare cases where we can still keep both normal hand 1-away and chiitoi 1-away, we usually do not force chiitoi. Even when our hand has 3 or 4 pairs, we still try to grasp the slim chance of a normal hand.


\subsubsection*{From Three Pairs to Chiitoi}

For hands like \ref{mj:chit-1} which has 3 pairs and is pretty much hopeless, chiitoi can be our plan B. Chiitoi 3-away accepts 21 tiles, while dropping any of \mj{3p5z} and going for normal hands also accepts \mj{15m8p456s3z} 22 tiles. On top of that, going by full efficiency will most likely only leave us bad wait dora 1. Therefore, we will not touch pairs in this hand, and start dropping from the middle tile \mj{3p} \citep[][pp.~82--83]{zero-teyaku}。

\tehai{1146m 379p 3557s 335z}[East round, North seat, turn 3, dora \mj{9p}][chit-1]

We draw \mj{9p} in the next turn and \ref{mj:chit-1} becomes \ref{mj:chit-2}. Again, we do not force chiitoi here; drop the weakest kanchan \mj{46m} \citep[][p.~83]{zero-teyaku}.

\tehai{1146m 799p 3557s 335z}[East round, North seat, turn 3, dora \mj{9p}][chit-2]


\subsubsection*{Forget about Chiitoitsu When Our Hand is Full of Good Shapes}

For both \ref{mj:chit-3} and \ref{mj:chit-4}, we \textit{could} drop any of the tiles with a single copy and make our hand ryoutenbin of chiitoitsu 1-away and normal hand 2-away, but we \textit{should not} do that with \ref{mj:chit-3} which has 3 ryanmen-toitsus. With \mj{7z} anko, we can open our hand and proceed just fine; we are 2-away form normal hand indeed, but it is faster than chiitoi. On the other hand, for \ref{mj:chit-4} which has one kanchan-toitsu, dropping \mj{1p} loses \mj{2p} acceptance but gains \mj{7m6p} acceptance, so it is actually 2 more tiles of acceptance. Therefore, we confirm a ryanmen in \ref{mj:chit-3} and drop \mj{1p} from \ref{mj:chit-4} \citep[][p.~83]{zero-teyaku}.

\tehai{788m 334677p 05s 777z}[East seat, turn 5, dora \mj{2s}][chit-3]

\tehai{788m 133677p 05s 777z}[East seat, turn 5, dora \mj{2s}][chit-4]

Hand \ref{mj:chit-5} has neighboring pairs and our acceptance is less ideal, but we still exclude chiitoi from the possibilities early in a round and drop \mj{4m}. When we start discarding second row discards, drop \mj{78m} and advance normal hand and chiitoi in parallel \citep[][p.~84]{zero-teyaku}.

\tehai{334478m 22388s 777z}[Dora \mj{1p}][chit-5]

Hand \ref{mj:chit-6} has a kanchan and is much less ideal than previous hands. No matter how early it is, drop \mj{7m} and advance normal hand and chiitoi in parallel \citep[][p.~84]{zero-teyaku}.

\tehai{334479m 22388s 777z}[Dora \mj{1p}][chit-6]

As shown in Table \ref{tab:chit-7}, dropping either of \mj{8p6s} from \ref{mj:chit-7} will give us the same number of 17 tiles of acceptance. If we drop \mj{8p}, the neighboring pairs \mj{2233p} is forced to serve as head candidates and \mj{14p} draws cannot make it mentsu + ryanmen. Still, \mj{1p} draw cancels tanyao out, so it doesn't hurt much when it backfires; the draws around \mj{5667s} would hurt more than that. Therefore, we drop \mj{8p} here; if we happen to reach chiitoi tenpai, take it \citep[][p.~18]{uzaku-nankiru}.

\tehai{2233556678p 5667s}[East-1, East seat, turn 5, dora \mj{4m}][chit-7]

\begin{table}[ht]
  \centering
  \begin{tabular}{ll}
    Drop & Acceptance \\
    \hline
    \mj{8p} & \mj{2347p57s} 17 tiles \\
    \mj{6s} & \mj{123457p} 17 tiles
  \end{tabular}
  \caption{Acceptance of \ref{mj:chit-7}.}
  \label{tab:chit-7}
\end{table}

Hand \ref{mj:chit-8} has enticing dora dabuton. However, we are already at turn 9, and the later it is the harder it is for dora to come out. Thus, we would take chiitoi as our plan B. Dropping either of \mj{4m7p} gives us ryoutenbin of normal hand of chiitoitsu, and the difference lies in the shapes after we call \mj{1z} pon, as shown in Table~\ref{tab:chit-8}. If we keep \mj{567p} mentsu, it would be difficult for \mj{3344p} to function as ryanmens; conversely, and somewhat unexpectedly, dropping \mj{67p} in tandem leaves us \mj{34p} ryanmen + \mj{55s} pair + \mj{344m} ryanmen-toitsu. Therefore, drop \mj{7p}, which gives us chiitoi possibility and the best acceptance after \mj{1z} pon \citep[][p.~20]{uzaku-nankiru}.

\tehai{344m 3344567p 55s 11z}[East-1, East seat, turn 9, dora \mj{1z}][chit-8]

\begin{table}[ht]
  \centering
  \def\arraystretch{1.2}
  \begin{subtable}[t]{0.48\linewidth}
    \centering
    \begin{tabular}{ll}
      \multicolumn{2}{l}{\mj{44m 3344567p 55s -1'11z}} \\[3pt]
      Drop & Acceptance \\
      \hline
      \mj{4m} & \mj{23458p5s} 17 tiles \\
      \mj{3p} & \mj{4m2458p5s} 17 tiles \\
      \mj{4p} & \mj{4m2358p5s} 17 tiles \\
      \mj{5s} & \mj{4m23458p} 17 tiles
    \end{tabular}
    \caption{With \mj{3m} already dropped.}
  \end{subtable}
  %
  \begin{subtable}[t]{0.48\linewidth}
    \centering
    \begin{tabular}{ll}
      \multicolumn{2}{l}{\mj{344m 334456p 55s -1'11z}} \\[3pt]
      Drop & Acceptance \\
      \hline
      \mj{4m} & \mj{25m235p5s} 19 tiles \\
      \mj{3p} & \mj{245m25p5s} 19 tiles \\
      \mj{6p} & \mj{245m25p5s} 19 tiles \\
      &
    \end{tabular}
    \caption{With \mj{7p} already dropped.}
  \end{subtable}
  \caption{Shapes of \ref{mj:chit-8} after \mj{1z} pon.}
  \label{tab:chit-8}
\end{table}


\subsubsection*{Setting Up First Row Discards}

Efficiency wise, it should be \mj{8m} drop from \ref{mj:chit-9}. However, this hand is very likely rii-nomi, so we would like to have chiitoi as our backup plan. We can drop either of \mj{59s} to keep acceptance in souzu suit, or we can drop \mj{7s}, and make \mj{9s} look like sotogawa, which increases our win rate \citep[][pp.~84--85]{zero-teyaku}.

\tehai{8899m 33445p 579s 33z}[South seat, turn 4, dora \mj{1p}][chit-9]

Likewise, drop \mj{3m} first from \ref{mj:chit-10}, even though both \mj{36m} are equally bad for a chiitoi tanki wait.

\tehai{136m 88p 6699s 22334z}[East seat, turn 4, dora \mj{1s}][chit-10]

Note that in both cases we are still at turn 4 where our discards are very likely considered sotogawa. If it is already second row discards, sotogawa trap becomes ineffective, and we can stop thinking about how our river looks.


\subsubsection*{Passive Chiitoi}

Chiitoi hands, much like toitoi hands, lack defense. When we are in the middle of a round where it is likely that our opponents would call riichi next turn, the floater we keep in our hand would be more defense oriented. Therefore, while \mj{2z} is twice cut and gives us furiten if we end up waiting on it, we should still drop \mj{6p} so that, even if some one calls riichi, we can still drop safe tiles while maintaining chiitoi 1-away \citep[][p.~85]{zero-teyaku}.

\tehai{499m 67799p 11779s 2z}[East seat, turn 10, dora \mj{4m}, \mj{2z} is twice cut and furiten][chit-11]


\subsubsection*{Choosing between Chiitoi and Toitoi}

According to \citet{gendai-toitoi},

\begin{itemize}
  \item The following conditions favor toitoi,
  \begin{itemize}
    \item We have an anko.
    \item Yakuhai or honitsu is feasible.
    \item We have terminal pairs, or pairs close to terminals.
    \item Agaritoppu in all last
  \end{itemize}
  \item The following conditions favor chitoi,
  \begin{itemize}
    \item We have a pair which is already twice cut
    \item We have no anko
    \item We have no other yaku
    \item We have a lone dora
    \item Our pairs are middle tiles
    \item We are in the lead and would like to avoid unnecessary risk
  \end{itemize}
\end{itemize}

Hand \ref{mj:chit-toi-1} looks like \mj{8s} drop for \mj{7z} pon or \mj{7z} drop for pinfu. \mj{8s} drop would only leave us a 1000 hand, so unless we are agaritoppu, it would not be ideal. \mj{7z} drop is the menzen way, but it could also be slow while none of the menzen yakus such as pinfu and iipeiko are guaranteed. In this regard, the possibilities of chiitoi and toitoi would give us maximal expected value.

For chitoi and toitoi, \mj{4m34p} would be equally useless, but since we still have \mj{7z} pair, to not completely confine ourselves to chitoi, consider the shapes after we actually call \mj{7z} pon, \mj{34p} would counterintuitively be the best drop here \citep{gendai-toitoi}.

\tehai{33477889m 34p 88s 77z}[Dora \mj{4z}][chit-toi-1]

Hand \ref{mj:chit-toi-2} has both conditions that favor toitoi and conditions that favor chitoi. Considering that, in the case we end up chitoi tanki, our win rate plummets, so we drop \mj{8s} and go for toitoi unless we are in the lead \citep{gendai-toitoi}.

\tehai{33388m 1155p 8s 4566z}[Dora \mj{4z}][chit-toi-2]

\begin{itemize}
  \item Conditions that favor toitoi:
  \begin{itemize}
    \item We have an anko
    \item We have yakuhai pair \mj{6z}
    \item We only have one middle tile pair \mj{55p}
  \end{itemize}
  \item Conditions that favor chitoi:
  \begin{itemize}
    \item We have a lone dora
  \end{itemize}
\end{itemize}

For hands with two anko and four pairs, our choice will be based on (1) if there is yakuhai pair and (2) how many ideal chitoi tanki wait tiles are left. If we drop \mj{3m8s} from \ref{mj:chit-toi-3}, any draws other than \mj{8m15p6z} would give us instant chiitoi tenpai; on the other hand, if we drop \mj{5p}, we can either ponten\footnote{A portmanteau of pon and tenpai, meaning that we reach tenpai after we call pon and drop a tile.} or even go for suuanko. Since we already have \mj{66z} pair, toitoi would be the best course.

\tehai{33388m 1155p 888s 66z}[Dora \mj{4z}][chit-toi-3]

%=============================================================
\section{Call Judgement}

Calls, melds, or \textit{fuuro}, refer to \textit{chi}, \textit{pon}, and \textit{kan}, the methods you claim a tile from other players discards. When we call and open our hands, the state of \textit{menzen} expires, and menzen yakus such as riichi iipeiko and pinfu also become impossible, thus decreasing the expected value of our hand. We cannot touch the open part of our hand, so we will defend against riichi with fewer tiles as well. The merits of calling are thus speed and confirmation of yaku, as summarized in Table~\ref{tab:naki-compare}.

\begin{table}[ht]
  \centering
  \begin{tabular}{lll}
    & Open & Menzen \\
    \hline
    Speed & $\bigcirc$ & $\times$ \\
    Value & $\times$ & $\bigcirc$ \\
    Defense & $\times$ & $\bigcirc$
  \end{tabular}
  \caption{Comparison between open hands and menzen hands \citep[][p.~87]{hirasawa-naki}.}
  \label{tab:naki-compare}
\end{table}

Call judgement is based on the following factors:

\begin{itemize}
  \item How fast we can reach tenpai by opening our hand
  \begin{itemize}
    \item Does the hand ``look like menzen''?
    \item Can we get rid of bad shapes with calls?
    \item How good and how far from tenpai is the result shape?
    \item Are there any other fast opponents?
  \end{itemize}
  \item How much defense there is after shortening our hand
  \item How much the value of our hand decreases (except for cases such as agaritoppu where value is relevant)
  \item Whether the open part would alarm other players (and whether they care)
  \item Turn; the later it is, the better it would be for us to open our hand for tenpai
\end{itemize}

It is very difficult for humans to perform real-time calculations and weigh these factors, and the decision process can be laden with preferences and biases and easily affected by many other things even outside of mahjong table, so often different people have different call judgement. This section is compiled with material from various sources and put in a way as coherent as possible, but there still could be some disagreements between these sources. We could have some leeway in call judgement in very specific conditions, but the general guidelines should stay the same.

As a rule of thumb, if any of the following conditions are met, go ahead and make the call \citep[][pp.~36--41, p.~45, pp.~58--61]{hirasawa-naki}:

\begin{itemize}
  \item The increase in value brought by riichi is limited. That is, there are no menzen yakus such as iipeiko and pinfu.
  \item The value of our hand is already mangan.
  \item We can get rid of bad shape or drained up ryanmens.
\end{itemize}

If all of the following conditions are met, do not call \multicitep{hirasawa-naki, p.~48; rb1, p.~216}:

\begin{itemize}
\item Cheap; if we open our hand, its value is no more than 2000.
\item Far away; there is no head, or it has so many bad shapes that it is very likely we have to make 3 calls.
\item There are no safe tiles.
\end{itemize}


%--------------------------------------------
\subsection{Tanyao}

When we open our hand with tanyao as its main yaku, it is called \textit{kuitan}. Since kuitan is completely made of middle tiles, while it is very easy to score, it also comes with a risk of low defense. For players with menzen preference, the criteria for kuitan would be higher, and cheap kuitan are rarely forced except for situations such as agaritoppu. As a rule of thumb, the criteria for kuitan are \multicitep{yoteru-badnaki; tetsusenryaku, pp.~170--175},

\begin{itemize}
\item Ponten or chiten for hands worth 1000 points
\item Call all bad waits for 3900 hands; leave ryanmen calls for mid round.
\end{itemize}


\subsubsection*{Only Call for Tenpai for Hands Worth 1000 Points}

Except for agaritoppu, regardless of the shapes of the rest of our hand, regardless of the callable part is bad shape or not, refrain from early cheap calls \multicitep{yoteru-badnaki; tetsusenryaku, pp.~170--175; hirasawa-naki, pp.~127--128}. Therefore, do not open \ref{mj:naki-tanyao-1000-1} or \ref{mj:naki-tanyao-1000-2} or \ref{mj:naki-tanyao-1000-3}.

\tehai{24m 3468p 3358899s}[East-1, non-dealer, turn 4, dora \mj{7z}][naki-tanyao-1000-1]

\tehai{23467m 246p 35889s}[West seat, turn 4, dora \mj{1s}][naki-tanyao-1000-2]

\tehai{367m 234499p 3377s}[East-1, West seat, turn 5, dora \mj{3z}][naki-tanyao-1000-3]

If we can get rid of bad shape and also get tenpai, it is fine to call. We can call \mj{3m7p} from \ref{mj:naki-tanyao-1000-4} \citep{yoteru-badnaki}.

\tehai{24m 23468p 34588s 4z}[East-1, non-dealer, turn 4, dora \mj{7z}][naki-tanyao-1000-4]


\subsubsection*{Hands Worth 2000 are Tricky}

Generally we can treat hands worth 2000 as if they are worth 1000 only \citep[][p.~179]{tetsusenryaku}. That is, do not open unless opening our hand gives us tenpai. This 1000 points of gap is not too much of a difference.

For hands with one dora, we are at the sweetspot of reaching mangan with riichi, so we would still be inclined to stay menzen even at 1-away. According to \citet[][p.~33, p.~38]{hirasawa-naki},

\begin{itemize}
\item If we have other menzen yakus, stay menzen until the end of second row discards.
\item If there are no other menzen yakus and there is no much difference between menzen and open in terms of value, open our hand as soon as we enter second row discards.
\end{itemize}

Hand \ref{mj:naki-1} is perfect 1-away and has pinfu chance which could make it 7700, or 5200 if there is no pinfu anyway. If we open our hand here, it will only be 2000 or 3900 if we draw any of the remaining aka doras. Thus we should not open this hand until about turn 12, unless we are agaritoppu, there are opponents who already openend their hands, or the ryanmens are drained up.

\tehai{405667m 22667p 34s}[East-1, West seat, dora \mj{3z}][naki-1]

If our hand is ryanmen + kanchan 1-away, we can call kanchan part regardless of turn, and leave ryanmen part after about turn 9 \citep[][p.~40]{hirasawa-naki}. For \ref{mj:naki-4}, we can chi \mj{4s} at any turn, or chi \mj{58p} when its turn 9 or later.

\tehai{405667m 2267p 35s 2z}[East-1, North seat, dora \mj{3z}][naki-4]

Hand \ref{mj:naki-8} has two kanchans, and one may be tempted to call \mj{7m7p} chii. However, note that this hand has no head. If we make any call, we have 2 less kinds of head candidate, and it becomes more likely that we have to sacrifice ryanmens for heads. Do not open headless hands easily \citep[][p.~55]{hirasawa-naki}.

\tehai{3468m 3468p 067s 34z}[East-1, West seat, dora \mj{2z}][naki-8]


\subsubsection*{Call All Bad Shapes for Hands Worth 3900}

\textbf{Call all bad shapes for hands worth 3900} \citep[][pp.~176--181]{tetsusenryaku}. For hands like \ref{mj:naki-tanyao-3900-1} which is laden with bad shapes or hands like \ref{mj:naki-tanyao-3900-2} which seems close to menzen riichi, \mj{4s} is a must call. For \ref{mj:naki-tanyao-3900-1}, we can also call \mj{35p} chi. We can also pon \mj{8s} to make our hand mangan.

\tehai{23468m 246p 35889s}[West seat, turn 4, dora \mj{8s}][naki-tanyao-3900-1]

\tehai{23467m 334p 35889s}[West seat, turn 4, dora \mj{8s}][naki-tanyao-3900-2]

We can make calls to confirm tanyao even if it does not advance our hand in terms of shanten number \citep[][pp.~127--128]{hirasawa-naki}. If we pon \mj{7s} with \ref{mj:naki-tanyao-3900-3} and drop \mj{9p}, we are still 2-away, but our acceptance is wider with \mj{3m4p} floaters instead of \mj{58m9p37s} when we stay menzen.

\tehai{367m 234499p 3377s}[East 1, West seat, turn 5, dora \mj{3s}][naki-tanyao-3900-3]


\subsubsection*{Pon Dora with Hands Already Worth 7700 if Your Opponents Would Not Mind}

The merits of dora pon include increase in value as well as wider acceptance when our hand is laden with bad shapes. We have already seen cases such as \ref{mj:naki-tanyao-3900-1} and \ref{mj:naki-tanyao-3900-2} where dora pon increase our value from 3900 to 7700. For \ref{mj:naki-tanyao-3900-1} specifically, if we deliberately get rid of the head, our acceptance becomes \mj{68m246p35s}, which is wider than \mj{7m35p4s} if we stay closed. For \ref{mj:naki-tanyao-3900-2}, \mj{33p} can still serve as our head, so we can pon dora \mj{8s} without slowing down our hand.

If our hand is already mangan (or 7700), this merit diminishes, and in a table where your opponents would try to sabotage your visibly expensive hands, the increase in value no longer justifies the decrease in win rate. For hands such as \ref{mj:naki-tanyao-7700-1} which is already full of bad shapes, a wider acceptance is still favorable. For \ref{mj:naki-tanyao-7700-2} though, we should only call \mj{4s} chi and save us unnecessary attention \citep[][pp.~188--195]{tetsusenryaku}.

\tehai{23468m 246p 3088s 4z}[West seat, turn 4, dora \mj{8s}][naki-tanyao-7700-1]

\tehai{23467m 334p 3088s 4z}[West seat, turn 4, dora \mj{8s}][naki-tanyao-7700-2]


\subsubsection*{Calls that Save Us from Bad Wait and Dropping Dora}

If we draw \mj{25s} into \ref{mj:naki-tanyao-3} first, we end up awkward tanki. If kamicha drops \mj{4m}, we can chi, drop \mj{9p}, and make our hand a flexible kuitan 1-away. Any of \mj{4578p25s} will give us good wait tenpai \citep{nemata-133}.

\tehai{35m 067789p 34678s}[Dora \mj{4z}][naki-tanyao-3]

If we stay closed with \ref{mj:naki-tanyao-4}, \mj{0s} is very likely useless and we end up dora 1 bad wait riichi. Therefore, we should opt to chi \mj{7p6s}, drop \mj{89s}, and go for kuitan dora 2 \citep{nemata-133}.

\tehai{2344588m 68p 0789s}[Dora \mj{6p}][naki-tanyao-4]

Likewise, double dora \mj{0p} is very likely useless for \ref{mj:naki-tanyao-5} and we could end up dora 1 bad wait riichi. If any of \mj{27m6s} comes out, claim it, and drop the entire \mj{1s} anko. It may sound bizarre, but considering it could be tanyao dora 3 if we get to keep \mj{0p}, it should be worth going for \citep{nemata-133}.

\tehai{2234077m 0p 11157s}[Dora \mj{5p}][naki-tanyao-5]

Hand \ref{mj:naki-3} is already tanyao dora 3, which we want to win no matter what, but is hindered by \mj{12p} penchan. Calls that confirm tanyao would help us tremendously even if they don't advance our shanten number. Thus, call any \mj{2457m}, and call \mj{3p} with \mj{24p} kanchan. \mj{58s} chi is also fine \citep[][p.~61]{hirasawa-naki}.

\tehai{3406m 12456p 3367s}[East-1, West seat, dora \mj{3s}][naki-3]


\subsubsection*{Tanyao Calls in All-Last}

When we are in all-last and there is only 500 points left for us to overtake first place or get out of last place, tanyao calls become immensely important. While it is true that each call sacrifices one turn of tsumo, our hand is no longer confined to menzen and is thus faster than menzen. There is no need to consider defense as well since it is better to fight for a slim chance for a better placement than to fold and completely give up the chance of winning. The following hands assume that we absolutely to win 500 points in all-last for a better placement unless otherwise stated.

Chi \mj{6p} with \mj{57p} and drop \mj{9p} from \ref{mj:naki-tanyao-oorasu-1}. Likewise, chi \mj{36p} with \mj{45p} and drop \mj{9p} from~\ref{mj:naki-tanyao-oorasu-2}. Both hands are perfect 1-away, and we stay menzen in flat situations. In all-last and in dire need to win any hand, value becomes irrelevant and we'd rather get to fastest tenpai with any yaku than wait for menzen riichi which is slow and has lower win-rate. Thus, the point here is to find ways to get rid of non-tanyao tiles, or to slide completed blocks such as \mj{789p} for tanyao \citep{mark2-tanyao-naki}.

\tehai{34m 567789p 45588s}[All-last, dora \mj{1s}][naki-tanyao-oorasu-1]

\tehai{34m 456789p 45588s}[All-last, dora \mj{1s}][naki-tanyao-oorasu-2]

Keep in mind that dropping \mj{9p} after \mj{6'78p} chi is \textit{kuikae} which is forbidden in most rulesets. The reason why \ref{mj:naki-tanyao-oorasu-1} and \ref{mj:naki-tanyao-oorasu-2} are special is because \mj{789p} is connected to other completed blocks. Do better see this, we consider shapes minus the terminal tile \mj{9p}. The shape \mj{56778p} can be either \mj{567p} + \mj{78p} or \mj{57p} + \mj{678p}. Likewise, \mj{45678p} can be \mj{45p} + \mj{678p}. However, if we chi \mj{7p} with \mj{56p} from \mj{567789p} and drop \mj{9p}, we are left a furiten ryanmen-toitsu instead of a completed block, so this is not a call that really advances our hand. Table~\ref{tab:kuikae} lists all possible shapes and calls to allow us pivot to tanyao.

\begin{table}[ht]
  \centering
  \begin{tabular}{lll}
    Shape & Combination & Call \\
    \hline
    \mj{123345p}    & \mj{234p} + \mj{35p} & \mj{4p} chi \\
    \mj{123456p}    & \mj{234p} + \mj{56p} & \mj{47p} chi \\
    \mj{123444p22s} & \mj{234p} + \mj{44p} + \mj{22s} & \mj{4p2s} pon \\
    \mj{12456p}     & \mj{24p} + \mj{56p} &  \mj{347p} chi
  \end{tabular}
  \caption{Calls that allow us to swap blocks without committing \textit{kuikae} \citep[][p.~20]{kawamura-naki}. In all cases, we can drop \mj{1p} after.}
  \label{tab:kuikae}
\end{table}

There are indeed cases where we only have an isolated \mj{789p} block and must sacrifice one safe tile for tanyao like \ref{mj:naki-tanyao-oorasu-3}. Chi \mj{6p} with \mj{78p} and \mj{3z} \citep{mark2-tanyao-naki}.

\tehai{34678m 789p 4588s 3z}[All-last, dora \mj{1s}][naki-tanyao-oorasu-3]

\textbf{Always call shimocha's discards and be cautious about calling kamicha's discards} \citep[][p.~19--20]{kawamura-naki}. If we call shimocha's discards, we will only let shimocha draw one more turn. If we call kamicha's discards instead, only our turn of draw is skipped. In other words, calling shimocha's discards benefits our opponents the least. For hands like \ref{mj:naki-tanyao-oorasu-4}, we can absolutely call \mj{8m2p} from shimocha, as well as any other stuff that we'd usually hesitate to call for any other hands. This, of course, is not to say that we cannot call from kamicha; it is a thing to be kept in mind and taken into consideration all the time.

\tehai{4588m 1223p 45678s}[All-last, dora \mj{3z}, \mj{9s} is twice cut][naki-tanyao-oorasu-4]

Chi \mj{47m} and pon \mj{2m8p} with \ref{mj:naki-tanyao-oorasu-5} is a must, but also keep in mind that we can also chi \mj{36p58s} ryanmens. Some people may be wary of the possibility of ending up bad waits if we get rid of good shapes, but even then, the bad waits are 2(8) tiles, not too bad as final waits. On the contrary, the bad waits in \ref{mj:naki-tanyao-oorasu-6} are middle tile pairs, so do not chi \mj{25p58s} ryanmens with \ref{mj:naki-tanyao-oorasu-6} unless they are drained up \citep[][pp.~21--22]{kawamura-naki}.

\tehai{1223456m 4588p 67s}[All-last, dora \mj{3z}][naki-tanyao-oorasu-5]

\tehai{1234566m 3466p 67s}[All-last, dora \mj{3z}][naki-tanyao-oorasu-6]

\mj{37m7p} chi is a must for \ref{mj:naki-tanyao-oorasu-7}. However, besides those, \mj{245s} is also callable choices. To get two tanyao blocks out of \mj{2234s}, we can see it as \mj{22s} + \mj{34s} or \mj{23s} + \mj{24s}. In the former case, \mj{25s} chi becomes obvious, but what is better than \mj{2s} chi is \mj{2s} \textit{pon} which leaves us a ryanmen that accepts 5 copies rather than a pair that only accepts 1 copy. If kamicha drops \mj{5s} instead, chi leaves us a \mj{2s} pair that still accepts 2 copies of \mj{2s}. In the case of \mj{23s} + \mj{24s}, the true bottleneck is \mj{23s} ryanmen which gives us furiten risk when we draw \mj{1s}, so \mj{4s} is a must and \mj{3s} is a no-go. Notice that manzu part is a sankanchan, but calling \mj{5m} chi destroys manzu part and leaves us two floating tiles \citep[][pp.~55--56]{kawamura-naki}.

\tehai{2468m 68p 2234s 234z}[All-last, dora \mj{1p}][naki-tanyao-oorasu-7]


% Perhaps this section is too advanced?
%
%\subsubsection*{上家の変則な捨て牌を利用する}
%
%\textbf{上家が変則な捨て牌をするなら通常より積極的に鳴く}\citep{hououlab-1}。立体~\ref{rittai:early-mid} のようなたった2巡目で上家から \mj{6p} が来たら、手牌の中に辺張の \mj{12p} があるにもかかわらず、\mj{0677p}（＝\mj{07p}＋\mj{67p}）から鳴く。
%
%\begin{rittai}[h]
%  \centering
%  
%  \drawrittai{
%    jicha/seat = 東,
%    jicha/hand = 688m 120677p 456s 1z,
%    jicha/river = 9s 9m,
%    jicha/point = 24200,
%    %
%    toimen/river = 2z 5z,
%    toimen/point = 3400,
%    %
%    kami/river = 6z 6p,
%    kami/point = 21300,
%    %
%    shimo/river = 1z 9^p,
%    shimo/point = 51100,
%    %
%    dora = 4z
%  }
%  
%  \caption{上家がたった2巡目で中張牌の \mj{6p} を切る場合}
%  \label{rittai:early-mid}
%\end{rittai}
%
%立体~\ref{rittai:early-mid-2} は立体~\ref{rittai:early-mid} の次巡。ASAPINは \mj{456s} の面子をわざと壊して \mj{3s} をチーした。この鳴きは向聴数を戻すに見えるけど、次の巡目搭子を作る浮き牌は \mj{6m6s} になる。
%
%\begin{rittai}[h]
%  \centering
%  
%  \drawrittai{
%    jicha/seat = 東,
%    jicha/hand = 688m 267p 456s 1z -6'07p,
%    jicha/river = 9s 9m 1p,
%    jicha/point = 24200,
%    %
%    toimen/river = 2z 5z 9s,
%    toimen/point = 3400,
%    %
%    kami/river = 6z 3^s,
%    kami/point = 21300,
%    %
%    shimo/river = 1z 9^p 7z,
%    shimo/point = 51100,
%    %
%    dora = 4z
%  }
%  
%  \caption{次巡で上家が \mj{3s} を切る場合}
%  \label{rittai:early-mid-2}
%\end{rittai}
%
%立体~\ref{rittai:early-mid-3} も同じ上家が早めの巡目で中張牌を切り始める場面。そしてASAPINも亜両面から鳴く。しかもこの手牌はすでにドラドラなので、スルーするわけもない\citep{hououlab-1}。
%
%\begin{rittai}[h]
%  \centering
%  
%  \drawrittai{
%    jicha/seat = 北,
%    jicha/hand = 267m 3058p 4406 88s,
%    jicha/river = 9^m 1p,
%    jicha/point = 15900,
%    %
%    toimen/river = 1p 1^z 1^m,
%    toimen/point = 60700,
%    %
%    kami/river = 9p 85s,
%    kami/point = 25600,
%    %
%    shimo/river = 271^s,
%    shimo/point = 17800,
%    %
%    dora = 7z
%  }
%  
%  \caption{上家が3巡目で中張牌の \mj{5s} を切る場合}
%  \label{rittai:early-mid-3}
%\end{rittai}

%--------------------------------------------
\subsection{Yakuhai}

\textbf{In most cases, call yakuhai pon.} In any of the following cases, refrain from calling \citep[][p.~103]{hirasawa-naki}:

\begin{itemize}
\item Our yakuhai pair is the only head and calling it would slow us down (unless our hand is full of bad shapes.)
\item Our hand is still cheap and far away after calling
\end{itemize}

Unlike kuitan, if we skip one yakuhai call, it would be very difficult for the last tile to come out again. To have our yaku confirmed, the criteria for yakuhai pon is generally lower than kuitan.


\subsubsection*{Yakuhai Pair is The Only Head}

For hands like \ref{mj:naki-7a} which has only one pair \mj{66z}, it does not make our hand faster if we pon it. No matter how many times it is cut, do not pon \mj{6z} \citep[][p.~104]{hirasawa-naki}.

\tehai{3457m 67p 2346s 266z}[East 1, West seat, dora \mj{3m}][naki-7a]

For hands such as \ref{mj:naki-7b}, even though it is close to riichi, its value is 2600 minimum, which is just 600 points higher than open hand. If it is mid round, we must pon \mj{6z}; calling pon in early turns would not hurt too much \citep[][pp.~105--106]{hirasawa-naki}.

\tehai{345m 667p 23467s 66z}[East-1, West seat, dora \mj{3m}][naki-7b]

Hand \ref{mj:naki-7c} has a bad shape \mj{89s}, so getting rid of the only head will also make \mj{89s} part of our acceptance. In cases like this where we have bad shapes in our hand, deliberately getting rid of the head actually speeds our hand up.

\tehai{123m 23467p 89s 277z}[East-1, West seat, turn 5, dora \mj{4p}][naki-7c]

The same principle applies to shapes after yakuhai pon. If we pon \mj{2m} with \ref{mj:naki-20}, our acceptance changes from \mj{6m3p} to \mj{567m234p}. This is thus a must call. If our hand has one ryanmen like \ref{mj:naki-21}, or if our hand has an anko like \ref{mj:naki-22}, calling pon is still better \citep[][p.~199]{hirasawa-naki}.

\tehai{12257m 24p 789s -5'55z}[East-1, West seat, dora \mj{5z}][naki-20]

\tehai{12257m 45p 789s -5'55z}[East-1, West seat, dora \mj{5z}][naki-21]

\tehai{12257m 45p 888s -5'55z}[East-1, West seat, dora \mj{5z}][naki-22]

Hand \ref{mj:2-6-7} is already worth 7700 points. We should pon \mj{7z} first to have our yaku confirmed \citep[][pp.~164--167]{tetsusenryaku}.

\tehai{307m 4789p 207s 277z}[East-1, West seat, turn 4, dora \mj{8p}][2-6-7]


\subsubsection*{Cheap and Far Away}

If our hand is 3-away, is cheap, has no good shape, and has no pair other than yakuhai, refrain from calling \citep[][p.~110]{hirasawa-naki}. For hands like \ref{mj:naki-7d}, \mj{7z} pon is a big no-go.

\tehai{135m 4789p 268s 277z}[East-1, West seat, turn 3, dora \mj{7m}][naki-7d]


\subsubsection*{Atozuke}

\textit{Atozuke} refers to the state where only a part of our waits give us a yaku. For example, \ref{mj:atozuke-0a} waits on \mj{7s7z}, but only \mj{7z} tsumo or ron allows us to win this hand. In a ruleset called \textit{ariari}\footnote{The first \textit{ari} refers to kuitan, and the second \textit{ari} refers to atozuke. Some rulesets ban both, and those rulesets are called \textit{nashinashi}. By the way, there is no such thing as \textit{arinashi} or \textit{nashiari} which only bans one of kuitan and atozuke.}, atozuke is allowed and all fine and dandy.

The issue would be how and why we approach atozuke. For example, is there anything we can call other than \mj{7z} in \ref{mj:atozuke-0b}? Do we call any of \mj{35m2p} from kamicha? Would the number of dora in our hand change how much we want to go for atozuke?

\tehai{234m 340p 77s 77z -2'13p}[][atozuke-0a]

\tehai{246m 13340p 77s 577z}[][atozuke-0b]

Some people may shun atozuke, afraid of the unstable state and attention from other opponents. However, know that even if your hand looks like anything but kuitan, your opponents do not always try to strangle you for that. The criteria for atozuke is as follows \citep[][pp.~146--151]{tetsusenryaku},

\begin{itemize}
  \item Always call bad shapes (unless someone calls riichi first) in any of the following situations,
  \begin{itemize}
    \item Agaritoppu
    \item First place in south round
    \item Hands with mangan value
  \end{itemize}
  \item If we are not in the lead and our hand is not mangan either,
  \begin{itemize}
    \item Call only when it gives us 2000 1-away
    \item Call only when it gives us 3900 2-away
  \end{itemize}
\end{itemize}

For hands like \ref{mj:atozuke-1}, its value does not justify loss in defense after \mj{3m3p} call except for agaritoppu.

\tehai{24779m 1267p 5s 477z}[Turn 4, dora \mj{7s}][atozuke-1]

For hands like \ref{mj:atozuke-2} which has 1 dora, opinions may vary from person to person. Since calling \mj{37m2p} only gives us 2000 2-away and thus does not meet the criteria shown previously, skipping those should be fine.

\tehai{24779m 2267p 5s 477z}[Turn 4, dora \mj{2m}][atozuke-2]

If we have 2 doras like \ref{mj:atozuke-3}, calling any of \mj{37m2p} is fine.

\tehai{24779m 2206p 5s 477z}[Turn 4, dora \mj{2m}][atozuke-3]

If opened parts such as \mj{2'13s} from \ref{mj:atozuke-4} would alarm your opponents and make them hesitate to drop live dora \mj{7z}, \mj{2s} chi would not be wise. If your opponents do not care even a little bit, \mj{2s} chi is fine.

\tehai{23899m 45p 136s 477z}[East-1, West seat, turn 4, dora \mj{7z}][atozuke-4]

Do not call \mj{9p} chi from \ref{mj:atozuke-5}. Even though we have dora 3 and we would like to move forward as fast as possible, \mj{9p} chi does not get rid of bad shapes, and it reduces the number of tiles that could give us another pair. Calling any of the kanchan is acceptable \citep[][pp.~116--117]{fkt}.

\tehai{268m 3078p 1357s 66z}[Non-dealer, turn 2, dora \mj{6z}][atozuke-5]

If we have 3 pairs other than yakuhai pair, we can call any of them. With 4 pairs in our hand, toitoit is also within reach; atozuke with yakuhai even if we do not get to score toitoi is also fine. Call \mj{2p} pon with \ref{mj:atozuke-6} \citep{yuusee-atozuke}.

\tehai{37m 22378p 1122s 55z}[East-2, turn 1, dora \mj{4z}][atozuke-6]

Situation \ref{rittai:atozuke-7} shows two copies of \mj{1p} already seen. Kamicha drops a tile; it could be \mj{1p} or \mj{4p} (or any other things.) If it is \mj{4p}, since \mj{23p} is not critical, we can skip it just fine, and \mj{14p} are likely useless for kamichi who will drop those later. However, if kamicha drops \mj{1p} instead, it is another story. If we choose to skip it, there is only one copy of \mj{1p} left somewhere in the wall, and our \mj{23p} ryanmen becomes almost as weak as a kan \mj{4p}. While it is true that a visible \mj{1'23p} may alarm others and make \mj{7z} less likely to come out, skipping \mj{1p} also weakens our waits. Likewise, we should call outside tiles such as 2 and 3 more compared to their inside suji counterparts (5 and 6) \citep[][pp.~131--132]{kawamura-naki}.

\begin{rittai}[h]
  \centering
  \drawrittai{
    jicha/seat = S,
    jicha/hand = 2468m 236p 3688s 77z,
    jicha/river = 32z,
    %
    toimen/river = 2z 1p,
    %
    kami/river = 4z 9s ?,
    %
    shimo/river = 1p 4z,
    %
    dora = 8s,
  }
  \caption{Call kamicha's \mj{1p}, but do not call \mj{4p}}
  \label{rittai:atozuke-7}
\end{rittai}


%--------------------------------------------
\subsection{Honitsu}

Calling \mj{1p} pon with \ref{mj:naki-honitsu-1} is fine, while some people may choose not to. We already have 4 blocks necessary for honitsu. Another draw of \mj{7z} gives us huge increase in value. If someone calls riichi first, we can drop \mj{5z} and buy us two turns \citep[][pp.~112--123, p~217]{fkt}.

\tehai{116m 1123468p 557z}[East round, East seat, turn 6, dora \mj{3z}][naki-honitsu-1]

Hand \ref{mj:naki-honitsu-2} is worse than \ref{mj:naki-honitsu-1}. The dealer's first drop is \mj{3z}. We can either pon \mj{3z} as a bluff and force honitsu, or skip it and hope for \textit{kyuushukyuuhai}\footnote{If we have 9 or more kinds of terminal and honor tiles in our hand with our first draw, without any other calls interrupting us, we can declare kyuushukyuuhai and force abortive draw.} or at least closed honitsu chiitoi. For rulesets which require more defense, skipping \mj{3z} pon is better \citep[][pp.~124--125, p.~217]{fkt}.

\tehai{7m 14p 114s 1233467z}[Non-dealer, haipai, dora \mj{2z}][naki-honitsu-2]

We can pretty much call anything with \ref{mj:naki-honitsu-3}. The visible pinzu taatsus of this hand is \mj{34p} and \mj{79p}, so it is clear that we can call \mj{258p}. We can also call dora \mj{3z}. However, \mj{3p} chi is also a viable option in early turns. Even though \mj{1334p} is a little bit awkward, if kamicha drops \mj{3p} this early, it is very likely that \mj{2p} is useless to him and will drop it later. With 3 copies of \mj{3p} visible, it would also be difficult for other opponents to make use of \mj{2p}, making it a good wait \citep[][pp.~121--123]{kawamura-naki}.

\tehai{24m 13344579p 336z}[East-1, south seat, early turns, dora \mj{3z}][naki-honitsu-3]


%--------------------------------------------
\subsection{Toitoiho}

If we are chiitoi 1-away and also have a yakuhai pair, generally calling is better. The more 1289 pairs we have in our hand, the better is to call. Therefore, with \ref{mj:2-8-1} if someone drops \mj{4p}, we must call it and call any other pairs \citep[][pp.~204--211]{tetsusenryaku}.

\tehai{1167m 2244p 788s 77z}[East-1, West seat, turn 4, dora \mj{8p}][2-8-1]

If we force toitoi with \ref{mj:2-8-2}, we will not be able to make us of dora \mj{3m} and our hand will only be 2600. Unless we want to confirm a yaku as soon as possible, do not open this hand.

\tehai{3m 3399p 11889s 224z}[East-1, West seat, turn 5, dora \mj{3m}][2-8-2]

If we are given 4 pairs in the haipai \ref{mj:2-8-3} and oya drops \mj{1m} first, calling \mj{1m} pon is fine. If your opponents get alarmed easily, skipping pon is also fine. For some people, 1 yakuhai pair does not qualify as a toitoi hand \citep[][pp.~132--133, p~217]{fkt}.

\tehai{117m 1266p 224s 557z}[Non-dealer, dora \mj{6s}][2-8-3]

Hand \ref{mj:2-8-4} has 4 pairs and one dora anko, but 3 of its pairs are middle tiles and hard to call. There are people who wants to call \mj{9m} pon on turn 1, and there are also people who wants chitoi even at turn 6 \citep[][pp.~134--135, p.~217]{fkt}.

\tehai{24499m 222577p 55s}[East round, non-dealer, dora \mj{2p}][2-8-4]

%--------------------------------------------
% Skipping the following subsections - either too short or not very well elaborated
%\subsection{片アガリ}
%\subsection{食い伸ばし}
%\subsection{食い替え}

%--------------------------------------------
\subsection{Kan}

Any of the following calls are collectively referred to as kan:

\begin{itemize}
\item Daiminkan: claiming others discards by opening an anko in our hand.
\item Kakan: when we draw the 4th copy of a tile which we already pon'd. This is sometimes referred to as \textit{shouminkan}.
\item Ankan: when we draw the 4th copy of a tile which we already have an anko of.
\end{itemize}

An opened quadruplet is called \textit{minkantsu} or \textit{minkan} for short, and a closed quadruplet is called \textit{ankantsu} or \textit{ankan} for short. 

The main benefits of kan is to increase hand value, either through the kan dora that is revealed after the kan call is completed, or through the additional fu from the quadruplet. For example, a terminal anko gives 8 fu, and a terminal ankan gives 32 tu. Another benefit of kan is an additional draw, which is why when we are tenpai with ryanmen kan is usually better. Kan dora also benefits our opponents, and it is three times likely that someone will get kandora than we getting it. Therefore, kan far away from tenpai is not wise and often advised against.


\subsubsection*{Daiminkan}

When 3 out of the following 4 criteria is met, we can commit daiminkan \citep[][p.~140]{kawamura-naki}:

\begin{enumerate}
\item We are in last place and in dire need of points
\item It is unlikely for us to stay menzen (for example, only one completed block and we'd call everything else)
\item By calling, we can get faster than others
\item At least 2 hans
\end{enumerate}

The reason why 2 hans would be the sweetspot of kan is due to the value brought by the kan dora. Lets say we only have a chun ankou and call daiminkan with it, with 1 kan dora our hand value jumps from 1300 to 2600 (or 2000 to 3900 for dealer.) Considering the value increase of others who likely have doras, this 1300 points of different is easily outnumbered.

For hands like \ref{mj:kan-1}, since we have 1 dora, in last place, and all other pairs are perfect for calls, we can commit \mj{7z} daiminkan \citep[][pp.~139--140]{kawamura-naki}.

\tehai{11688m 99p 06s 6777z}[South-3, dealer, early to mid turn, last place, dora \mj{3z}][kan-1]

Hand \ref{mj:kan-2} has no dora, but we can force honitsu with the hand to make it at least 3 han, and any kan dora increase our hand value by a lot. Therefore, we can kan \mj{7z}. However, if any of \mj{19p7z} becomes new dora, we don't need to force honitsu anymore and we still have mangan, in which case just follow pure efficiency \citep[][pp.~141--142]{kawamura-naki}.

\tehai{34m 11246899p 777z}[Last place, dora \mj{3z}][kan-2]

Hand \ref{mj:kan-3} has 2 dora, but our hand is also close to menzen tenpai, so we should not kan \mj{7z} unless someone calls or it is mid turn. Even so, only kan shimocha's \mj{7z} and skip toimen and kamicha's draws. If we only have one aka dora, the value brought by riichi is bigger than kan, so do not kan at all \citep[][pp.~141--142]{kawamura-naki}.

\tehai{11m 34078p 40s 4777z}[Last place, dora \mj{3z}][kan-3]


\subsubsection*{Kakan}

Calling kakan when we are in ryanmen tenpai is perfectly fine, but it would be awkward for tanki wait in Situation \ref{rittai:kakan}. Opinions vary \citep[][pp.~198--199, p.~222]{fkt}.

\begin{rittai}[h]
  \centering
  \drawrittai{
    jicha/seat = N,
    jicha/hand = 456m 340p 2334s -7z -7'77z,
    jicha/river = 1m 9s 1s 8s 3z 6s,
    %
    toimen/river = 3z 4z 1z 9p 8m 4s,
    toimen/riverb = 5s,
    %
    kami/river = 1p 8p 1z 1m 7m 9m,
    %
    shimo/river = 2z 1s 2s 8p 8m 6z,
    %
    dora = 2m
  }
  \caption{}
  \label{rittai:kakan}
\end{rittai}


\subsubsection*{Ankan}

\textbf{Call ankan when our hand is 1-away} \citep[][pp.153--157]{ooi-gundou}. Hand \ref{mj:kan-4} is still 2-away, so we don't call kan yet here; drop \mj{5z} and see how things goes. If it becomes \ref{mj:kan-4a} which is 1-away, we can kan \mj{4p}.

\tehai{66778m 444478p 56s 5z}[South seat, turn 7, dora \mj{7m}][kan-4]

\tehai{667789m 444478p 56s}[South seat, turn 9, dora \mj{7m}][kan-4a]

However, for hands like \ref{mj:kan-5} which is also chitoi 1-away, kan makes chitoi impossible. Therefore, unless chitoi is completely out of the question, refrain from kan and maintain normal hand 1-away and chitoi 1-away \citep[][pp.~163--164]{ooi-gundou}.

\tehai{11688m 2222p 44566s}[West seat, turn 9, dora \mj{2s}][kan-5]


%--------------------------------------------
\subsection{Canceling Ippatsu and Shifting Haitei}

\textbf{The only effective ippatsu canceling is when someone can overtake us with haneman tsumo} \citep[][p.~188]{hirasawa-naki}. In other words, when we are in first place, have a lot of safe tiles, and second place who is 14000 below us calls riichi, we can go make any calls to cancel their ippatsu. In other cases, deliberately shortening our hand and giving up our tenpai chance would not worth it.

\textbf{Always shift haitei} \citep[][p.~190]{hirasawa-naki}, especially when our hand is hopeless. However, shifting haitei only makes a difference of 100 points in terms of expected value \citep[][p. 188]{miinin-choice}. In Situation~\ref{rittai:haitei} where we have enough genbutsu, if we drop \mj{1m} we won't be able to chi \mj{3m}, and dropping any of \mj{27s} makes us unable to call genbutsu from kamicha. Therefore, drop \mj{2m} here. Since 2 copies of \mj{2m} are already on the table, there is no way we can pon it.

\begin{rittai}[h]
  \centering
  
  \drawrittai{
    jicha/seat = N,
    jicha/hand = 12206m 23556p 247s -8s,
    jicha/river = XXXXXX,
    jicha/riverb = XXXXXX,
    jicha/riverc = XXX,
    %
    toimen/river = 3z 9s 2m 2z 6z 4'm,
    toimen/riverb = 4^p 5^z 7^3^s 3^m 2^z,
    toimen/riverc = 2^s 7^z 1^2^m,
    %
    kami/river = XXXXXX,
    kami/riverb = XXXXXX,
    kami/riverc = XXXX,
    %
    shimo/river = XXXXXX,
    shimo/riverb = XXXXXX,
    shimo/riverc = XXXX,
    %
    dora = 2p
  }
  
  \caption{When choosing a genbutsu to drop, choose the one that gives us the most chances to call and shift haitei. Drop \mj{2m}.}
  \label{rittai:haitei}
\end{rittai}














\nocite{*}
\bibliographystyle{jecon}
\bibliography{refmj}


\end{document}